% 本文件由 LaTeX 排版 Agent 根据有效 translation.json 自动生成。
% author=Athanasius; work_id=2808; title=驳亚略派

\chapter{\WorkTitle{驳亚略派}}

% structure: p0001 source=h1 target=omitted reason=page-title-already-rendered-by-parent

\Emph{第一部分} \newline{} \Emph{第二部分}\par

\SourceRecord{Translated by M. Atkinson and Archibald Robertson.；Nicene and Post-Nicene Fathers, Second Series；Vol. 4.；Edited by Philip Schaff and Henry Wace.；Buffalo, NY: Christian Literature Publishing Co.,；1892.；<http://www.newadvent.org/fathers/2808.htm>.}

% structure: p0001 source=h1 target=section reason=page-title

\section{\WorkTitle{反驳亚略派}{（第一部分）}}

% structure: p0002 source=h2 target=subsection reason=relative-heading-rank; base=h2

\subsection{第一章 导言}

1. 我曾以为，在提供了那么多证明我清白的证据之后，我的敌人们会不再追究，并且会因他们对别人所进行的诬告而自责。然而，他们并未因此羞愧，尽管他们的诬告已被清楚地揭穿，却若无其事地坚持他们对我的诽谤，声称整个事件应当重新审理（并非为了让自己受到审判——他们对此避之不及——而是为了骚扰我，扰乱单纯之人的心神）。所以，我认为有必要在你们面前为自己辩护，好使你们不再倾听他们的怨言，而是谴责他们的邪恶和无耻的诽谤。我仅向你们这些心地真诚的人呈上辩护；至于那些好争辩的人，我满怀信心地诉诸于我所拥有的确凿证据。因为我的案件无需再次审判；因为已经审判过了，而且不止一两次，而是许多次。首先，它在我自己的国家，由近一百名当地的主教组成的会议中审理过；第二次在\Place{罗马}，当时由于\Person{优西比乌}的来信，他们和我们都被召来，共有五十多位主教聚会；第三次在\Place{撒底卡}召开的大会上，这是由最虔诚的皇帝\Person{君士坦提乌斯}和\Person{君士坦斯}下令召开的，那次会议上我的敌人们被斥责为诬告者而受到贬黜，而有利于我的判决得到了超过三百名主教的赞成，这些主教来自\Place{埃及}、\Place{利比亚}、\Place{彭塔坡里}、\Place{巴勒斯坦}、\Place{阿拉伯}、\Place{伊苏里亚}、\Place{塞浦路斯}、\Place{潘菲利亚}、\Place{吕基亚}、\Place{加拉太}、\Place{达契亚}、\Place{默西亚}、\Place{色雷斯}、\Place{达达尼亚}、\Place{马其顿}、\Place{伊庇鲁斯}、\Place{塞萨利}、\Place{亚该亚}、\Place{克里特}、\Place{达尔马提亚}、\Place{希西亚}、\Place{潘诺尼亚}、\Place{诺里克}、\Place{意大利}、\Place{皮塞嫩}、\Place{托斯卡纳}、\Place{坎帕尼亚}、\Place{卡拉布里亚}、\Place{阿普利亚}、\Place{布鲁提亚}、\Place{西西里}、整个\Place{阿非利加}、\Place{撒丁尼亚}、\Place{西班牙}、\Place{高卢}及\Place{不列颠}等省份。\par

此外，还有\Person{乌尔撒奇乌斯}和\Person{瓦伦斯}的证言，他们先前曾诬告我，后来却改变了主意，不仅赞成了有利于我的判决，还承认他们自己以及我的其余敌人都是诬告者；一个人作出这样的改变和这样的翻供，自然使\Person{优西比乌}及其同伙受到谴责，因为他们正是与他们一起密谋了反对我的诡计。如今，既然事情已经由这么多卓越的主教根据如此清楚的证据审理并决定了，每个人都会承认再讨论下去是多余的；否则，如果此时此刻再启动调查，可能又得再次讨论再次调查，这样的琐屑之事将永无止境。\par

2. 这么多主教的裁决足以使那些仍然想对我提出某种指控的人哑口无言。然而，当我的敌人也作出对我有利、对他们自己不利的证言，声明针对我的诉讼是一场阴谋时，谁还会不引以为耻而继续怀疑呢？法律要求凭两三个证人的口供（\BibleRef{申 17:6}），审判才能定案；我们在这里既有这么多证人对我有利，更有来自敌人的证据补充；因此那些仍然反对我的人不再重视他们自己的任意判断，而是诉诸暴力，不以公正的推理，反想加害于那些揭露他们的人。因为这正是使他们恼怒的主要原因：他们暗中策划、在角落里自行设计的伎俩，被\Person{瓦伦斯}和\Person{乌尔撒奇乌斯}揭露并公之于众；因为他们心里明白，他们的翻供，一方面证明了他们所加害之人的清白，另一方面定了他们自己的罪。\par

事实上，这导致了他们在前述的\Place{撒底卡}会议上受到贬黜，这是理所当然的；因为正如古时的\Person{法利塞人}，当他们为\Person{保禄}辩护时（\BibleRef{宗 23:9}），便完全揭露了他们和\Person{犹太人}一同构陷他的阴谋；又如蒙福的\Person{达味}，当迫害者承认说：\BibleQuote{我犯了罪，我儿达味！}（\BibleRef{撒上 26:21}）时，便证明他受了不义地迫害；这些人也是如此；他们被真理挫败，便提出请求，将一份文书呈交给\Place{罗马}主教\Person{犹利乌}。他们也写信给我，要求与我修好，尽管他们曾散布关于我的种种流言；或许现在他们羞愧难当，因为他们看到那些他们企图靠主的恩宠加以毁灭的人，仍然活着。与此举动一致的是，他们也将\Person{亚略}和他的异端绝罚；因为他们知道，\Person{优西比乌}及其同伙是由于自己的错误信仰，而非别的什么，才结党反对我；所以当他们决定承认对我的诬告时，立刻也放弃了那个敌基督的异端，他们正是为此才编造了那些不实之词。\par

以下是各位主教在历次大会上所写的支持我的信函，首先是\Place{埃及}主教们的信函。\par

\Emph{埃及大会通谕}\par

在\Place{亚历山大}聚集的圣会议，由\Place{埃及}、\Place{底巴依德}、\Place{利比亚}和\Place{彭塔坡里}各地的主教组成，致书于散居各地的天主教会的主教们，就是在主内所亲爱的、极其渴慕的弟兄们，问候你们。\par

最亲爱的弟兄们，我们原本可以就\Person{尤西比乌斯}及其同党对\Person{亚大纳削}的阴谋提出辩护，申诉他所遭受的苦难，并揭露他们的一切诬告，无论是在他们阴谋之初，还是在他抵达\Place{亚历山太}之时。但当时情况不许可，你们也知道这一点；后来，在\Person{亚大纳削}主教归来以后，我们以为他们会因明显的不义而自觉羞愧、蒙受耻辱：因此我们决意保持沉默。然而，在他遭受了所有这些严酷苦难之后，在他退居\Place{高卢}之后，在他流寓异乡、远离故土之后，在他因他们的诽谤而九死一生——幸亏皇帝仁慈，他才幸免——这样的痛苦，即使最残忍的敌人也当满足，他们却仍然不知羞耻，再次对教会和\Person{亚大纳削}肆行侮慢；因他得救而心生恼怒，竟敢策划更凶恶的阴谋陷害他，并且准备控告，毫不惧怕圣经上记载的话：\BibleQuote{作假见证的，必不免受罚；说谎的嘴，杀死灵魂。}\BibleRef{箴 19:5}因此我们再也不能缄默，惊异于他们的邪恶，以及他们在阴谋中所表现的那永无止息的争竞之心。\par

你们看，他们不住地用新的报告扰乱帝王的听闻，反对我们；他们不住地写致命的信，为要毁灭那位与他们不虔不义为敌的主教。他们再次写信给皇帝们控告他；再次阴谋陷害他，以从未发生的屠杀指控他；再次企图流他的血，以从未犯过的谋杀罪控告他（前次若非我们有仁慈的皇帝，他们早已用诽谤置他于死地）；再次急切地，至少是可以这样说，要求将他放逐，而他们却装作哀怜那些据称被他放逐者的不幸。他们在我们面前悲叹从未发生之事，而且不满足于已加诸他身上的，还想给他加上其他更残酷的待遇。他们就是如此温和、如此慈悲、如此公正；更确切地说（因为要讲实话），他们就是如此邪恶、如此恶毒；他们不是凭自己的虔敬和正义，像主教所当行的，而是靠着恐吓和威胁来获取尊敬。他们竟敢在致皇帝的信中倾泻出连教外好争竞的人也不会使用的言语；他们以多起谋杀和屠杀指控他，并且不是向总督或其他上级官员，而是直接向三位\OriginalTerm{奥古斯都}{Augusti}；也不怕路途遥远，只要能够把一切较大的法庭都塞满他们的控告就行。确实，最亲爱的弟兄们，他们的事业就在于控告，且是最具严重性质的控告，因为他们所上诉的法庭是地上最庄严的。他们通过这些调查，除了促使皇帝判处死刑之外，还能有什么别的目的呢？\par

因此，该受悲叹哀悼的，不是\Person{亚大纳削}的遭遇，而是他们自己的行为；人们更当为他们哀哭，因为此类行径正应受人悲悼，正如经上记着说：\BibleQuote{不要为死人哭号，不要为他悲伤；却要为离家远行的人痛哭，因为他不得再见故土。}\BibleRef{耶 22:10}因为他们的整封信无非图谋死亡；他们的目的就是杀人，只要获准便行杀戮，若不获准，便放逐。而皇帝中最虔敬的父亲曾允许他们这样做，他以放逐\Person{亚大纳削}代替死刑，从而满足了他们的狂怒。如今，这种行为连普通基督徒都不会做，几乎连异教徒都少有，更何况是主教，他们自称教导人正义；我们认为你们基督徒的良心必定立刻明白这一点。他们怎能禁止别人控告弟兄，自己却成为控告者，且向皇帝控告？他们怎能教导人怜悯他人的不幸，自己却连我们的放逐都不能满足？因为众所周知，我们这些主教都曾被普遍判以放逐，我们都视自己为被放逐之人；而今，我们再次认为我们是与\Person{亚大纳削}一同重返故乡，不再像从前那样为他哀哭悲伤，反而得了极大的鼓励和恩宠——愿主将这恩宠继续赐予我们，不容\Person{尤西比乌斯}及其同党摧毁。\par

即使他们对\Person{亚大纳削}的指控属实，他们自身也有一项具体的罪状：违背了基督信仰的教训，在他已被放逐、遭受磨炼之后，竟再次攻击他，以谋杀、屠杀和其他罪名控告他，并在皇帝耳旁大造声威，攻击主教们。然而，他们的邪恶何等深重，你们以为他们是何等样人？他们口中所出的每一句话都是虚谎，所提出的每一项指控都是诽谤，无论口中还是笔下，全无一点真实！那么，我们终于要进入这些事端，应对他们最新的指控。这将证明，他们先前在公会议和审判中的表现是无耻的，或者说，他们的话是不实的，此外，也将暴露他们现今所提出的种种。\par

5. 我们确实羞于对这类指控作任何辩护。但由于我们那些鲁莽的指控者抓住任何罪名，声称在\Person{亚大纳削}归来后发生了凶杀和屠戮，我们恳求你们容忍我们的回答，尽管它可能稍长；因为情势所迫。无论是\Person{亚大纳削}本人，还是因他之故，均未发生过任何凶杀；因我们先前说过，控告者迫使我们进行这令人屈辱的辩护。屠杀与囚禁与我们的教会无涉。\Person{亚大纳削}从未将任何人交付刽子手；就他而言，监狱也从未被惊扰。我们的圣所如今，一如往昔，是纯洁的，只以基督的圣血和对祂的虔敬崇拜为荣。\OriginalTerm{司铎}{Presbyter}或\OriginalTerm{执事}{Deacon}均未遭\Person{亚大纳削}杀害；他未曾犯下凶杀，未曾使任何人遭放逐。但愿他们从未对他行如此之事，也未曾让他亲历其苦！此地无人因他而被放逐；除了\Person{亚大纳削}本人，即\Place{亚历山大里亚}的主教，是他被他们放逐，而如今他既已复位，他们又再次图谋将他卷入与先前相同、甚至更为残酷的阴谋之中，鼓动唇舌，编造各种虚假且致命的言辞来攻击他。\par

看，他们如今竟将行政长官的行为归咎于他；尽管他们在信中明确承认，是\OriginalTerm{埃及总督}{Prefect of Egypt}对某些人宣判，如今却毫不羞耻地将此判决归咎于\Person{亚大纳削}；而事实上，那时他尚未进入\Place{亚历山大里亚}，仍在流放地归来的途中。确实，他当时身在\Place{叙利亚}；因为我们必须在辩护中指出他离家路途遥远，故不应为总督或\OriginalTerm{埃及总督}{Prefect of Egypt}的行为负责。但即便假设\Person{亚大纳削}当时身在\Place{亚历山大里亚}，总督的作为又与\Person{亚大纳削}何干？然而，他甚至不在该地；且\OriginalTerm{埃及总督}{Prefect of Egypt}所为并非基于教会理由，而是因某些原因，你们将从记录中得知；我们得知他们所写内容后，便勤勉查考这些记录，并已传送给你们。既然他们如今大声指控某些既非他所为、亦非因他而发生的恶事，仿佛它们确曾发生，又对此等恶事作证，犹如对其存在深信不疑；那么，让他们告诉我们，他们从哪次会议、依据何种证据、经由何种司法调查而得悉此等事？但若他们无法提出任何此类证据，仅凭其一面之词，我们就将此事留给诸位判断，对于他们先前的指控，事情究竟如何发生，以及他们为何如此描述。事实上，这不过是诽谤，是我们敌人的阴谋，是一种狂躁失控的性情，且是为\OriginalTerm{亚略}{Arian}疯人张目的不虔敬，他们疯狂反对真正的虔诚，企图根除\OriginalTerm{正统派}{orthodox}，好使不虔敬的鼓吹者从此可以无所畏惧地宣讲任何他们所喜欢的教义。事情的原委如下：——\par

6. 当\Person{亚略}——\OriginalTerm{亚略}{Arian}疯人的异端即由他得名——因其不虔敬而被可敬的\Person{亚历山大}主教逐出教会时，\Person{优西比乌}及其同伙，作为他那不虔敬的门徒与同伴，自认也被驱逐，便屡次写信给\Person{亚历山大}主教，恳求他勿将异端者\Person{亚略}拒于教会之外。然而，\Person{亚历山大}出于对基督的虔诚，拒绝接纳那不虔敬之人，他们便将怨恨转向当时还是\OriginalTerm{执事}{Deacon}的\Person{亚大纳削}，因为他们细加打听，得知他深得\Person{亚历山大}主教的亲近与器重。及至在\Place{尼西亚}举行的大公会议上，他们见识到他对基督的虔诚，他勇敢地发言反对\OriginalTerm{亚略}{Arian}疯人的不虔敬，此后他们对他的仇恨便大大加深。当天主将他提升至\OriginalTerm{主教圣职}{Episcopate}时，他们积蓄已久的恶意便迸发为烈火，因畏惧他的\OriginalTerm{正统信仰}{orthodoxy}和对他们不虔敬的抵抗，他们（尤其是\Person{优西比乌}，他因意识到自己的恶行而内心不安）便对他施展种种奸诈的图谋。他们使\OriginalTerm{皇帝}{Emperor}对他怀有偏见；屡次以大会威胁他；最后在\Place{提洛}召集会议；直至今日，他们仍不断写信攻击他，如此不肯罢休，甚至挑剔他受任\OriginalTerm{主教圣职}{Episcopate}之事，用尽一切手段表露他们的敌意与仇恨，散布虚假传言，唯一的目的是以此贬损他的品格。\par

然而，他们现今所作的歪曲陈述，恰恰证明了他们先前的说法乃是谎言，且是针对他的一场纯粹阴谋。因为他们声称，\Quote{在\Person{亚历山大}主教去世后，某少数人提及\Person{亚大纳削}的名字，便有六、七位主教秘密地在一处隐秘之地选举了他}：这便是他们写给\OriginalTerm{皇帝们}{Emperors}的内容，竟毫无顾忌地捏造如此弥天大谎。然而，事实是，全体群众和\OriginalTerm{天主教会}{Catholic Church}的所有子民如同一心一体地聚集，呼喊、要求\Person{亚大纳削}作他们教会的主教，将此作为他们向基督公开祈祷的主题，并恳求我们允准，达多日多夜之久，他们既不离开教会，也不容我们离开；这一切我们都是见证人，全城和全省亦然。他们从未说一句反对他的话，正如这些人所歪曲的，而是以他们所能想到的最尊贵的称号来称呼他，称他为良善、虔诚、基督徒、\OriginalTerm{苦修者}{ascetic}、真正的主教。且他是经我们团体的大多数，在全体民众的面前与欢呼声中当选的，我们这些选举他的人也为此作证，我们的见证无疑比那些不在场、如今却散布这些虚假说辞的人更为可信。\par

但\Person{尤西比乌斯}居然挑剔\Person{亚大纳削}的任命---他本人或许从未正式领受过任何任命；即便有过，自己也已使之失效。因为他起初拥有\Place{贝里托}的主教座，却离开了那里，来到\Place{尼科美底亚}。他离开第一处是违反法律的，侵入另一处也是违反法律的；他毫不留恋地抛弃了自己的，又毫无理由地占据了他人的；他因贪图另一个而失去了对第一个的爱心，甚至他所贪求得来的也没有守住。因为，看，他离开了第二处，又去占据别人的主教座，向四周别的城市投以贪婪的目光，以为虔敬在于财富和城邑的宏大，而轻视自己被派定去守护的天主的产业；不知道\BibleQuote{\Emph{哪里}甚至有两三个人因主的名聚集，主就\Emph{在他们中间}}\BibleRef{玛 18:20}；也不想想宗徒的话：\BibleQuote{我不以别人的劳苦为夸耀；}\BibleRef{格后 10:15-16} 也没有留意他所给的诫命：\BibleQuote{你与妻子有了约束吗？不要谋求脱离。}\BibleRef{格前 7:27} 因为这话若适用于妻子，那么它岂不更适用于教会，适用于同一主教职？凡是与这职分相结合的人，都不该寻求另一个，免得他按圣经被证明是犯奸淫的。\par

7. 然而，他虽深知自己的这些恶行，却胆敢挑剔\Person{亚大纳削}的任命---那任命曾得全体的尊荣见证---且竟敢以他的罢免来羞辱他，尽管他自己已被罢免，并有他人接任此职的活证。他自己与\Person{德奥格尼}既被罢免，又如何能罢免别人呢？这由他人接替其位的事实足以证明。因为你们很清楚，接替他们的，在\Place{尼科美底亚}是\Person{安菲翁}，在\Place{尼西亚}是\Person{克瑞斯托}，这全因他们自身的恶行以及与\Person{亚略}派的疯狂之徒的勾结，遭大公会议弃绝。然而，他们虽企图推翻那真正的大公会议，却将他们自己的非法集会冠以那名称；他们不愿大公会议的法令得以施行，却要强行推行他们自己的决议；他们使用大公会议之名，却拒绝服从如此伟大的会议。因此，他们并不在乎大公会议，只是假装如此，为的是根除正统信仰者，废除那真正而伟大的大公会议反对\Person{亚略}派的法令；他们一直在支持\Person{亚略}派，无论现在或从前，都胆敢对主教\Person{亚大纳削}散布这些谎话。他们从前所说，正如他们现在所捏造的，即在他进入时会举行混乱的集会，哀哭悲伤，民众愤然拒受。但事实并非如此，恰恰相反，到处充满了喜乐和欢跃，百姓蜂拥而至，渴望见到他。各教会满溢欢腾，处处向主献上感恩；所有圣职人员和神职人员都怀着如此情感注视着他，以至他们的灵魂充满喜悦，看那天为一生中最幸福的日子。我们又何须提及我们主教们中间洋溢的不可言喻的喜乐呢？因为我们先前已说过，我们自认曾分担他的苦难。\par

8. 这既是众所公认的实情，尽管他们极尽歪曲，他们所夸口的那个会议或审判又有何分量呢？他们既擅自干涉一件他们未曾目击、未曾审查、也未为此召开会议的案件，并且书写时仿佛确信其陈述属实，那么对于他们自称为了审议的事，又怎能取信于人呢？难道不该认为，他们在这件事和那件事上，都是出于对我们的敌意吗？当时所举行的是何种主教会议呢？那是一个追求真理的集会吗？他们中岂不几乎人人都是我们的仇敌吗？\Person{尤西比乌斯}及其党羽对我们的攻击，岂非出自他们对\Person{亚略}派的疯狂的热忱吗？他们岂非煽动其同党吗？我们岂非常常着文反对奉行\Person{亚略}教义的他们吗？\Place{巴勒斯坦} \Place{凯撒利亚}的\Person{尤西比乌斯}岂非被我们的宣信者指控向偶像献祭吗？\Person{乔治}岂非被证明为真福\Person{亚历山大}所罢免吗？他们岂非被指控种种罪行，有的犯这个，有的犯那个吗？\par

这样的人，怎能动念举行反对我们的会议呢？他们怎敢称那次会议为大公会议呢？那次会议由一位伯爵主审，有刽子手出席，由一个庭丁---不是教会的执事---把我们引入法庭；会议上只有伯爵一个人说话，其余的人都缄默不语，更确切地说，都服从他的指示。那些看似应当免职的主教，也因他的意旨而未被免职；而当他下令时，我们便被士兵驱赶---更确切地说，是\Person{尤西比乌斯}及其党羽在发号施令，而伯爵只是俯听其意愿。总之，可亲爱的弟兄，那还算什么大公会议呢？其目的在于按皇帝的喜好流放和杀害。他们的指控又是何性质呢？---因为这更是令人大大惊奇的事。有一位名叫\Person{阿尔塞尼奥}的，他们声称已被谋杀；他们又抱怨一个属于神圣奥迹的圣爵被摔碎了。\par

如今，\Person{阿尔塞尼乌斯}活着，并恳求被接纳进入我们的共融。他无需其他证据证明他仍在世，而是亲自承认，以自己的名义写信给我们的弟兄\Person{亚大纳削}，就是他们一口咬定被他杀害的人。那些不虔敬的卑鄙之徒竟不知羞耻地控告他谋杀了一个与他相距甚远，无论海路或陆路都相隔极大距离，且当时无人知其行踪的人。不仅如此，他们甚至胆大妄为地将他隐藏起来，置于隐匿之处，尽管他并未受到任何伤害；而且，若是可能，他们早就把他转移到另一个世界去了，不，甚至当真让他失去生命，以便无论谋杀陈述是真是假，他们都可以借此确定无疑地毁灭\Person{亚大纳削}。但为此也要感谢\OriginalTerm{天主眷顾}{divine Providence}，祂没有让他们的不义得逞，反而将活生生的\Person{阿尔塞尼乌斯}呈现在众人眼前，这清楚地揭穿了他们的阴谋和诽谤。他并不像躲避杀人犯那样避开我们，也不憎恨我们以为他受到了伤害（因为他确实没有遭受任何损害）；相反，他渴望与我们共融；他希望能被计入我们之中，并已为此写信表明心意。\par

9. 然而，他们依旧设下阴谋对付\Person{亚大纳削}，指控他谋杀了一个仍然活着的人；而正是这同一些人造成了对他流放的判决。因为使他遭流放的，并非皇帝们的父亲，而是他们的诬告。请考虑这是否真实。当没有发现任何不利于我们的同工\Person{亚大纳削}的证据，但那位\OriginalTerm{官员}{Count}仍以暴力威胁他，并激烈地反对他时，主教为躲避这暴力，便上到最虔敬的皇帝那里，抗议那位官员及他们的阴谋，并请求要么召集一次合法的主教会议，要么请皇帝亲自听取他对所控罪名的辩护。于是，皇帝愤怒地写信，传唤他们前来，宣布要亲自审理此案，并为此目的下令召开会议。于是，\Person{欧瑟比乌斯}及其同党上来，用虚假的罪名指控\Person{亚大纳削}，不是用他们曾在\Place{提洛}公布的那些罪名，而是指控他意图扣留运粮船只，仿佛\Person{亚大纳削}是个能假装阻止从\Place{亚历山大港}向\Place{君士坦丁堡}出口粮食的人。\par

我们的一些朋友与\Person{亚大纳削}一同在宫廷，听到了皇帝在接到此报告后发出的威胁。当\Person{亚大纳削}大声抗议这诽谤，并明确声明这不是事实时（因为他辩称，他一个穷人，又身居卑微，怎能做这样的事？），\Person{欧瑟比乌斯}毫不犹豫地当众复述这指控，并发誓说\Person{亚大纳削}是个富人，有权势，什么事都能做到；为的是让人由此认为他曾说过那样的话。这就是这些可敬的主教们对他提出的控告。然而，天主的\OriginalTerm{恩宠}{grace}胜过了他们的邪恶，因为它感动了虔敬的皇帝施以怜悯，皇帝没有处死他，而是判以流放。因此，他们的诬告，而非其他，是造成此事的缘由。因为皇帝在先前的信函中抱怨过他们的阴谋，谴责过他们的诡计，并斥责\OriginalTerm{梅肋提派}{Meletians}肆无忌惮，该受诅咒；总之，以最严厉的言辞论及他们。因为当他听到活死人的传闻时，大为震动；当他听闻有人虽未丧命却被说成遭谋害时，大为震动。我们已将信函寄给你们。\par

10. 但这些令人称奇的人物，\Person{欧瑟比乌斯}及其同党，为了装模作样地反驳事实真相和这信函中的声明，便抬出会议的名义，并将其程序建立在皇帝的权威之上。因此，有一位\OriginalTerm{官员}{Count}出席了他们的集会，士兵们充当主教们的护卫，而且还有皇室函件强迫任何他们所需的人出席。但请注意这里他们阴谋的奇特之处，以及他们大胆措施的反复无常，其目的无非是以种种方式将\Person{亚大纳削}从我们中间夺走。因为如果他们作为主教，只为自己要求对该案的审判权，那又何必需要官员和士兵出席呢？或者说，他们何以在皇室函件的准许下集会呢？再者，如果他们需要皇帝的支持并希望从他那里得到权威，那他们又为何废弃他的判决呢？当他在所写的信函中声明梅肋提派是诽谤者，肆无忌惮，而\Person{亚大纳削}最为无辜，并且对那捏造的谋杀活人的事件大加渲染时，他们何以认定梅肋提派说的是实话，而\Person{亚大纳削}犯有罪行呢？他们竟不羞于将使活人变为死人——这人不仅根据皇帝的判决还活着，而且在他们集会时也还活着，甚至直到今日仍在我们中间。关于\Person{阿尔塞尼乌斯}事件，就说这么多。\par

11. 至于那属於奥迹的\OriginalTerm{圣爵}{the cup belonging to the mysteries}，它到底是甚么？\Person{玛加略}又是在哪里把它打碎的？这就是他们到处散播的说法。但至於\Person{亚大纳削}，若非他们收买，连控诉他的人也绝不敢指责他。然而，他们将这罪行的起源归咎於他；其实连\Person{玛加略}也不该承受这指责，因为他对此清白无辜。他们竟然恬不知耻，在\OriginalTerm{慕道者}{Catechumens}面前炫耀\OriginalTerm{神圣奥迹}{sacred mysteries}，更糟的是，甚至在\OriginalTerm{异教徒}{heathens}面前；其实他们应当留意所记载的：\BibleQuote{隐藏君王的秘密是有益的}（\BibleRef{多 12:7}）；并且如同主所吩咐的：\BibleQuote{你们不要把圣物给狗，也不要把你们的珠宝投在猪前}（\BibleRef{玛 7:6}）。因此，我们不该在未入门者面前炫耀神圣奥迹，免得异教徒因无知而加以嘲笑，慕道者因过分好奇而跌倒。然而，那圣爵到底是甚么，是在哪里，在谁面前被打碎的？提出指控的是\OriginalTerm{梅勒提乌斯派}{Meletians}，他们半点也不值得相信，因为他们是\OriginalTerm{分裂者}{schismatics}和教会的仇敌，并非始於今日，而是自从那位蒙福的主教兼殉道者\Person{伯多禄}的时代就已如此。他们结党反对\Person{伯多禄}本人；他们诽谤他的继承人\Person{阿喀拉斯}；他们甚至在皇帝面前控告\Person{亚历山大}；既然他们如此精於此道，现在便将敌意转向\Person{亚大纳削}，其行径完全与昔日的邪恶如出一辙。正如他们曾诽谤那些在他之前的人，如今他们也诽谤了他。然而，他们的诽谤和诬告从未得逞，直到如今，他们以\Person{优西比乌}及其同伙为帮手和靠山，因这些人从\OriginalTerm{亚略}{Arian}的狂徒那里领受了不虔，这便促使他们阴谋反对众多主教，其中包括\Person{亚大纳削}。\par

然而，他们所说圣爵被打碎的地点，并非教会；那里没有\OriginalTerm{司铎}{presbyter}居住；而且他们所说\Person{玛加略}行事的那一天，也不是\OriginalTerm{主日}{the Lord's day}。既然那里没有教堂；既然没有人去举行那神圣职务；既然那日子并不需要用它；那么，这属於奥迹的圣爵到底是甚么，它又在何时、何地被打破呢？显然，无论在私人住宅还是在公共市场上，都有许多杯子；若有人打破其中一个，并不算犯不虔之罪。但那属於奥迹的圣爵，若有意打破，行者便成不虔之人，只有那些合法主持者才能持有。这是唯一能对此类圣爵所作的描述；别无其他；你们依法给予信友饮用；你们是按\OriginalTerm{教会的法典}{the canon of the Church}所领受的；这仅属於那些主持\OriginalTerm{天主教会}{Catholic Church}的人，因为唯有你们有权分施\OriginalTerm{基督的宝血}{the Blood of Christ}，别人都无权。然而，正如打破属於奥迹的圣爵者是不虔之人，那冒犯基督的宝血的人就更加不虔；凡\BibleQuote{如此行}（\BibleRef{格前 11:25}）而违反\OriginalTerm{教会的规则}{the rule of the Church}的人，便是如此。（我们说这话，并非是说\Person{玛加略}打破了分裂派的一只杯子，因为那里根本没有甚么杯子；怎会有呢？那里既没有主的殿，也没有任何教会所有的东西，况且那时也不是举行奥迹的时刻。）正是这个臭名昭著的\Person{伊斯奇拉斯}，他从未由教会委任此职；当\Person{亚历山大}接纳由\Person{梅勒提乌斯}\OriginalTerm{授秩}{ordination}的司铎们时，他甚至没有列在他们当中；因此，他甚至连从那边也没有领受授秩。\par

12. 那么，\Person{伊斯奇拉斯}究竟通过甚么方式成了司铎？是谁祝圣了他？是\Person{科卢图斯}吗？因为只剩下这个假设了。但众所周知，且没有人对此存疑：\Person{科卢图斯}是以司铎身份去世的，他所有的授秩均属无效，凡在\OriginalTerm{分裂}{schism}期间由他授秩者，均已降至\OriginalTerm{平信徒}{laymen}身份，并且以这身份出现在会众中。那么，怎能相信一个占据私人住宅的\OriginalTerm{私人}{private person}，会拥有一只\OriginalTerm{圣爵}{sacred chalice}呢？然而事实是，当时他们称一个私人为司铎，以此称号取悦他，以支持他对我们的不法行径；如今，作为对他指控我们的奖赏，他们又为他谋得了一座教会。因此，此人当初根本没有教会；但他因怀恨和顺从他们指控我们，如今得到了他昔日所没有的；不仅如此，也许他们甚至以\OriginalTerm{主教职}{Episcopate}酬报他的效劳，因为他就是这样四处宣扬的，并且因此对我们极为嚣张。这些奖赏如今由主教们赐给控诉者和诽谤者，但其实，对於共犯而言，理应是：既然他们让他成为他们的伙伴，共同参与行动，那么他们就也该让他成为他们自己主教职的同僚。但这还不是全部；请继续听他们那时的所作所为。\par

13. 他们虽然如此列阵反对真理，却不能胜过真理；而\Person{伊斯奇拉斯}在\Place{提洛}没有证实任何事情，反而被揭露为诽谤者，这诽谤破坏了他们的计谋，于是他们推迟审判，声称为获取新证据，将派遣他们一伙中的一些人前往\Place{马雷奥蒂斯}，仔细调查此事。因此，他们借助世俗权力，秘密派遣了我们曾因多种理由公开反对的那些人，因为他们是\Person{亚略}的同党，因此是我们的敌人；即\Person{迪奥格尼乌斯}、\Person{马里斯}、\Person{狄奥多鲁斯}、\Person{马塞多尼乌斯}，以及另外两位年纪既轻、心智也浅的人，来自\Place{潘诺尼亚}的\Person{乌尔撒奇乌斯}和\Person{瓦伦斯}；这些人长途跋涉，为要审判他们的敌人之后，又从\Place{提洛}出发前往\Place{亚历山大里亚}。他们不惜亲自充当见证人，尽管他们本是审判官，反而公开采用各种手段推进他们的图谋，并且无论付出多少辛劳，奔走多少路程，只求使正在进行的阴谋得以成功。他们将\Person{亚大纳削}主教扣留在异国他乡，而他们自己则进入他们敌人的城市，仿佛要在他的教会和他的人民身上狂欢作乐。更令人发指的是，他们带上了控告者\Person{伊斯奇拉斯}，却不准被告\Person{玛卡里乌斯}随行，而是将他拘留在\Place{提洛}。因为\Quote{\Place{亚历山大里亚}的司铎\Person{玛卡里乌斯}}便远近为此指控负责。\par

14. 因此，他们仅与控告者一同进入\Place{亚历山大里亚}，这控告者是他们住宿、饮食、杯盏的伙伴；他们带上\Place{埃及}总督\Person{菲拉格里乌斯}，前往\Place{马雷奥蒂斯}，在那里伙同前述之人，完全按照他们自己的方式，进行了所谓的调查。虽然众司铎屡次请求允许他们到场，他们却不准。城市及全乡的司铎都希望出席，以便查明那些被\Person{伊斯奇拉斯}收买的人究竟是谁、来自何方。然而他们禁止众\OriginalTerm{圣职人员}{Ministers}在场，自己却在异教徒面前对教堂、圣爵、圣桌和圣物进行审查；更有甚者，他们在调查那属于\OriginalTerm{神圣奥迹}{mysteries}的圣爵时，竟传唤异教徒作见证；他们声称那些被\Person{亚大纳削}借着\OriginalTerm{总税务官}{Receiver-general}的传票而藏匿起来、且不知下落的人，这些人，他们却只叫到他们自己与总督面前，公然使用他们的证词，并不知羞耻地断言这些人已被\Person{亚大纳削}主教隐藏。\par

但在这里，他们的唯一目的仍是置他于死地，于是他们再次假装活人已死，沿用他们在\Person{阿尔塞尼乌斯}一案中采用的方法。然而这些人依然活着，并且可以在他们本乡见到；但对于你们这些远离该地的人，他们却在这事上大造声势，仿佛这些人已经消失，以便因证据远离你们，他们可以诬告我们的同工司铎，仿佛他使用了暴力和世俗权力；而事实上，他们自己却在各方面借助了那权力和他人的支持。因为他们在\Place{马雷奥蒂斯}的行径，与在\Place{提洛}的如出一辙；在那里，有一位\OriginalTerm{伯爵}{Count}带着军事援助到场，不允许任何有违他们心意的话或事；在这里，\Place{埃及}总督同样带着一队人，恐吓教会所有成员，不许任何人作真实的见证。而最奇怪的是，那些来的人，无论是充当审判官还是见证人，或者更可能是为了满足他们自己和\Person{欧瑟比乌斯}的目的，竟与控告者住在同一地方，甚至就住在他家里，并在那里随心所欲地进行调查。\par

15. 我们想你们不会不知道他们在\Place{亚历山大里亚}所施的暴行，因为这些事已传遍各地。他们拔刀攻击圣洁的贞女和弟兄们；鞭打他们那些在天主眼中尊贵的人，以致他们的脚因鞭痕而跛瘸，尽管他们的灵魂在纯洁和一切善工上完整而健全。他们煽动各行业的人反对他们；并驱使异教徒群众去剥光他们的衣服，殴打他们，任意凌辱他们，并用他们的祭坛和祭品威胁他们。有一个粗鄙之徒，仿佛总督为取悦主教们而给他们开了方便之门，抓住一位贞女的手，把她拖向附近的一座祭坛，模仿逼迫时期强迫献祭的做法。这事发生后，贞女们纷纷逃散，异教徒便对教会发出一阵哄笑；主教们正在现场，并且就住在发生这一切的屋子里；而为了讨好他们，贞女们遭受了刀剑的攻击，面临各种危险、凌辱和肆意暴力。她们是在一个守斋日遭受这些待遇的，而且出自那些与主教们在屋内一同宴乐的人之手。\par

16. 预见到这些事，并考虑到敌人进入某地并非寻常灾祸，我们抗议这项任命。\Place{得撒洛尼}的主教\Person{亚历山大}也持有同样看法，他写信给那里的居民，揭发阴谋，证实密谋。他们确实把他算作自己人，视他为计谋的同伙；但这只证明了他们对他施加的暴力。因为甚至无耻的\Person{伊斯奇拉斯}本人也只是因恐惧和暴力而行事，被迫提出控告。作为证明，他亲自写信给我们的弟兄\Person{亚大纳削}，承认那里根本没有发生过所声称的那种事，他是被收买作假见证的。他作出这一声明，尽管\Person{亚大纳削}从未接纳他为\OriginalTerm{司铎}{Presbyter}，也未从他领受如此恩宠的称号，也未作为酬报而委托他建造一座教堂，也未期望得到主教职的贿赂；而所有这些，他都从他们那里获得，作为他承担控告的回报。再者，他的全家都和我们保持\OriginalTerm{共融}{communion}，如果他们在任何程度上受到伤害，就不会这样做了。\par

17. 为了证明这些是事实而非空言，我们有\Place{马雷奥提斯}所有司铎的见证，他们常陪同主教巡视，当时也写信反对\Person{伊斯奇拉斯}。但他们中那些前往\Place{提洛}的人，不被允许宣告真理，留在\Place{马雷奥提斯}的人也未能获得许可反驳\Person{伊斯奇拉斯}的诽谤。\Person{亚历山大}、众司铎以及\Person{伊斯奇拉斯}的信函副本也将证明同一件事。我们也送上了皇帝之父的信函，他在信中表达愤慨，因为有人指控某人谋杀了\Person{阿尔塞尼乌斯}，而\Person{阿尔塞尼乌斯}本人却还活着；他也对他们在\OriginalTerm{圣爵}{cup}一事上指控的变幻无常、自相矛盾表示惊讶；因为他们时而控告司铎\Person{玛加略}，时而又控告主教\Person{亚大纳削}，说他们用手打碎了圣爵。他还一方面宣称梅勒提乌斯派是诽谤者，另一方面宣称\Person{亚大纳削}完全无辜。\par

难道梅勒提乌斯派不是诽谤者吗？尤其\Person{若望}，他进入了教会，与我们共融后，又定自己的罪，不再参与有关圣爵的诉讼，但当他看到\Person{欧瑟比乌斯}及其同伙热心支持亚略派狂人时——虽然他们不敢公开与他们合作，却企图利用他人作面具——便像异教剧场中的演员一样，扮演了一个角色。这出戏的主题是亚略派的斗争；剧本的真正目的是他们的成功，但\Person{若望}和他的同伙被推上台扮演角色，以便在这些人的掩饰下，身着法官外衣的亚略派支持者可以赶走他们不信神的仇敌，牢固建立他们不虔敬的教义，并将亚略派带入教会。那些想驱逐真信仰的人，竭尽全力以不信神行事；他们选择了与\Person{基督}作战的那种不虔敬的角色，竭力消灭其仇敌，仿佛他们是邪恶的人；他们谴责我们打碎了圣爵，为的是显得\Person{亚大纳削}像他们一样，对\Person{基督}犯有不虔敬的罪。\par

对他们来说，提及这属于\OriginalTerm{奥迹}{mysteries}的圣爵有何意义？在那些支持对\Person{基督}不虔敬的人中，这种对圣爵的宗教尊重从何而来？不认识\Person{基督}的人，如何能认识\Person{基督}的圣爵？那些声称尊敬那圣爵的人，如何能侮辱圣爵的天主？或者，那些哀悼那圣爵的人，如何能企图谋杀用这圣爵举行奥迹的主教？因为他们若力所能及，早已谋害了他。那些哀悼\OriginalTerm{主教宝座沦丧}{the loss of the throne that was Episcopally covered}的人，如何能企图毁灭坐在这座位上的主教，致使座位无主教，人民失去敬虔的教义？因此，并非那圣爵，并非谋杀，也非他们谈论的任何骇人听闻的事，促使他们如此行事；而是前面提到的亚略派异端，为了这异端，他们阴谋反对\Person{亚大纳削}和其他主教，并继续对教会发动战争。\par

真正造成谋杀和流放的是谁？难道不是这些人吗？利用外部支持，阴谋反对主教的是谁？难道不正是\Person{欧瑟比乌斯}及其同伙，而非\Person{亚大纳削}，正如他们在信中所说？他和其他人都曾在他们手中受苦。就在我们说话的这时候，\Place{亚历山大利亚}的四位司铎，甚至未动身前往\Place{提洛}，就因他们的手段而被流放。那么，是谁的行为令人泪流满面和悲哭？难道不是他们吗？他们在犯下一次迫害之后，竟悍然追加第二次，不择一切虚假手段，为要毁灭一个不肯屈服于他们不虔敬异端的主教。由此产生了\Person{欧瑟比乌斯}及其同伙的敌意；由此产生了他们在\Place{提洛}的诉讼；由此产生了他们虚假的审判；由此也产生了他们现在甚至在没有任何审判的情况下写的信，言辞中充满对他们的陈述的十足自信；由此产生了他们在皇帝之父和众位最虔诚皇帝面前的诽谤。\par

18. 因为你们必须知道，现今所传的、为陷害我们的同工\Person{亚大纳削}的事，好使你们因此定他们的恶行，并明白他们无非是要杀害他。有一批谷物，由皇帝们的父亲赐给，用以供养某些寡妇，部分来自\Place{利比亚}，部分来自\Place{埃及}。她们至今都已领受，\Person{亚大纳削}从中毫无所得，只落得协助她们的辛苦。但如今，尽管受惠者自己并无怨言，且承认已领受谷物，\Person{亚大纳削}却被指控出售全部谷物，并将收益归为己有；皇帝也为此事写信，因所捏造的诽谤而追究其罪。那么，是谁散布这些诽谤呢？岂不是那些曾行一宗迫害，又毫不顾忌地发起另一宗的人吗？那些据称出自皇帝的信件，作者是谁呢？岂不是那些极力攻击\Person{亚大纳削}，且毫不顾忌地以言语、文字反对他的\OriginalTerm{亚略派}{Arians}吗？没有人会忽略行事如此之人，而去怀疑别人。的确，他们诽谤的证据显得最为明显，因为他们急切地以此为掩饰，将谷物从教会夺走，交给\OriginalTerm{亚略派}{Arians}。此事比其他任何事都更清楚地将矛头指向这计谋的设计者及其主子们，他们既不惜捏造对\Person{亚大纳削}的谋杀指控，作为在皇帝面前陷害他的卑鄙手段，也不惜从教会的圣职人员手中夺走穷人的生计，从而实在地为异端者牟利。\par

19. 我们也送上了我们在\Place{利比亚}、\Place{五城}和\Place{埃及}的同工的证词，从中你们同样可以得知那加在\Person{亚大纳削}身上的虚假指控。他们行这些事，为使真正恭敬天主的人从此因恐惧而保持沉默，好使不虔敬的\OriginalTerm{亚略派}{Arians}的异端得以取而代之。但是，亲爱的弟兄们，感谢你们的虔敬，因为你们常在信函中绝罚\OriginalTerm{亚略派}{Arians}，且从未容他们进入教会。揭露\Person{欧瑟伯}及其同伙也轻而易举，近在眼前。看，他们在先前有关\OriginalTerm{亚略派}{Arians}的信件——我们也已送上副本——之后，如今竟公然煽动\OriginalTerm{亚略狂人}{Arian madmen}反对教会，虽然整个\OriginalTerm{天主教会}{Catholic Church}已绝罚了他们；他们为亚略派按立了一位主教；他们用威胁和惊扰使各教会不安，好能在各处为他们的不虔获得帮手。更有甚者，他们派执事到\OriginalTerm{亚略狂人}{Arian madmen}那里，这些执事公开加入他们的集会；他们写信给他们，也收到回信，从而在教会内制造分裂，并与他们共融；他们派人到处去宣扬他们的异端，并弃绝教会，正如你们从他们写给\Place{罗马}主教的信，以及或许也写给你们本人的信中所见。因此，亲爱的弟兄们，你们看到，这些事并非不应受罚：它们确实可怕，且与基督的教导相背。\par

为此，我们聚集一处，写信给你们，请求你们的基督徒智慧接纳我们的这声明，并同情我们的弟兄\Person{亚大纳削}，并对\Person{欧瑟伯}及其同伙——他们谋划了这些事——表示愤慨，使此等邪恶与凶残不再侵害教会。我们呼吁你们作为此不义之事的复仇者，提醒你们宗徒的命令：\BibleQuote{你们应把那邪恶的人从你们中间除去！} \BibleRef{格前 5:13}。他们的行径实属邪恶，与你们的共融毫不相称。故此，即使他们再写反对主教\Person{亚大纳削}的信件（因为从他们发出的一切都是虚假的），也不要再理会他们；即便他们在信上签上\Place{埃及}主教的名字，也不要理会。因为显然，那写信的不会是我们，而是\OriginalTerm{梅勒提派}{Meletians}，他们一向是裂教者，直至今日仍在各教会中制造扰乱、兴立派系。因为他们祝圣不称职的人，几近外教人；他们的所作所为我们羞于形诸笔墨，但你们可从我们派到你们那里的人得知，他们也会将我们的信交与你们。\par

20. \Place{埃及}的主教们如此写信给所有主教，并致\Place{罗马}主教\Person{犹利乌}。\par

% structure: p0042 source=h2 target=subsection reason=relative-heading-rank; base=h2

\subsection{第二章 \Emph{\Person{犹利乌}致\Place{安提约基雅}的欧瑟伯派信函}}

\Person{欧瑟伯}及其同伙也写信给\Person{犹利乌}，想恐吓我，便请他召开会议，并如果他愿意，可亲自担任审判者。因此，当我上到\Place{罗马}时，\Person{犹利乌}便合宜地写信给\Person{欧瑟伯}及其同伙，并派去他自己的两位司铎，\Person{额尔皮迪乌斯}和\Person{费罗克塞努斯}。但他们一听说我，就张皇失措，未预料到我会去那里；他们拒绝了提议的会议，提出不充分的理由，但事实上他们是害怕\Person{瓦伦斯}和\Person{乌撒齐乌斯}后来所坦白的事情会被证明出来。然而，有五十多位主教聚集，在司铎\Person{维多}主持聚会的地方；他们认可了我的申辩，并给予我共融与爱德的确认。另一方面，他们对\Person{欧瑟伯}及其同伙表示极大的愤慨，并请求\Person{犹利乌}向那些写信给他的人写出以下信函。他照办了，并经由\Person{加比亚努斯伯爵}之手送出。\par

\Emph{\Person{犹利乌}的信函}\par

\Person{犹利乌}致他亲爱的弟兄们：\Person{达尼乌斯}、\Person{弗拉基利乌斯}、\Person{纳尔齐苏斯}、\Person{欧瑟伯}、\Person{马里斯}、\Person{马其顿尼乌斯}、\Person{狄奥多鲁斯}及他们的友人，他们从\Place{安提约基雅}写信给我，我在主内问候他们。\par

21. 我阅读了由我的司铎\Person{埃尔皮迪乌斯}和\Person{斐洛克塞努斯}带来的你的信，并且我惊讶地发现，尽管我以爱德和真诚的善意写信给你，你却以不得体和好争辩的态度回复我；因为写信人的骄傲和傲慢在这封信中表露无遗。然而，这样的情感与基督徒的信仰是不相符的；因为以爱德精神所写的事，同样也应当以爱德而非争辩的精神来答复。派遣司铎去同情受苦之人，并且希望那些曾写信给我的人能到那里去，以便能使争议问题迅速解决，使一切事务得以妥善安排，让我们的弟兄们不再遭受痛苦，并让你避免进一步的诽谤，这难道不是爱德的标志吗？但某些事似乎显示，你的脾性如此，甚至令我们不得不认为，即便是在那些你表面上尊崇我们的措辞中，你也是以讽刺伪装来表达的。我们所派遣到你那里的司铎，原本应当愉快地返回，却因他们在你们中间目睹的所作所为而忧闷归来。而我，在阅读了你的信后，经多方考虑，保留了它，心想毕竟你们中有些人会来，便没有必要把它公布出来，以免公开出示后，会使我们这里的许多弟兄难过。但当没有人前来，而此信又不得不展示时，我向你声明，他们全都感到震惊，几乎不能相信这样一封信竟然是你写的；因为它是以好争辩而非爱德的语言表达出来的。\par

现在，如果这封信的作者写下此信，是出于炫耀其语言能力的雄心，那么这种做法显然更适合其他主题：在教会事务中，需要的并非展示口才，而是遵守宗徒法典，并殷切留意不使教会中的任何一个小子跌倒。因为按照教会的说法，\BibleQuote{与其使一个小子跌倒，倒不如把磨石拴在他的颈上，投在海里。}\BibleRef{玛 18:6}但如果这封信之所以写成，是因为某些人彼此之间气量狭小而感到愤懑（因我不将此事归咎于所有人），那么最好根本不要怀有这类恼怒之情，至少不要含怒到日落；更不可让其有机会形诸文字。\par

22. 然而，究竟发生了何事，竟成为恼怒的正当理由？或者，我写给你们的信，在哪一方面如此冒犯你们？是因我邀请你们出席一次公会议吗？你们本当喜悦地接纳这个提议。那些对自己所处理的事务——或按他们所选定的说法，对自己的决议——有信心的人，通常不会因他人查问这样的决议而生气；他们反倒表现出极大的勇气，因为他们知道，自己既作出公正的决议，它就绝不会被推翻。聚集在\Place{尼西亚}大公会议的主教们，并非没有天主的旨意，一致同意，一个公会议的决议应在另一个公会议中接受审查，目的是让审判者，因着眼于此后的另一个审判，能以最大的审慎去调查事实，并使受他们判决的有关各方确知，他们所领受的判决是公正的，并非受先前审判者的敌意所左右。如今，如果你们不愿在你们自己的事上采用这种做法，尽管它由来已久，且已为大会所留意和推荐，你们的拒绝便是不合宜的；因为一项在教会中曾一度通行，并由公会议确立的习惯，却要被少数人所搁置，这是不合理的。\par

他们基于另一个理由，亦不能在此事上合理地动怒。当你们，即\Person{欧瑟伯}及其同党，派遣持你们信函的人——我是指司铎\Person{玛加略}，及执事\Person{玛尔提里乌}和\Person{赫西基乌}——抵达此间，并发觉他们无力反驳从\Person{亚大纳削}那里来的司铎们的论证，反而在各方面被驳倒并揭露时，他们便要求我召开一次公会议，并写信给在\Place{亚历山大}的主教\Person{亚大纳削}，同样也写给\Person{欧瑟伯}及其同党，以便在所有相关方到场的情况下，能作出公正的判决。他们且承诺在此情形下，会证明针对\Person{亚大纳削}所提出的一切指控。因为\Person{玛尔提里乌}和\Person{赫西基乌}已在我们面前被公开驳倒，而主教\Person{亚大纳削}的司铎们以极大的信心抵挡了他们：诚然，若必须说出真相，\Person{玛尔提里乌}及其同伙已被彻底颠覆；正是这导致他们希望召开一次公会议。现在，假设他们并没有要求召开公会议，而是我本人提议召开，以劝阻那些给我写信的人，并为我们那些抱怨遭受不公的弟兄们的缘故；即使在这种情况下，这项提议也是合理且公正的，因为它合乎教会惯例，且蒙天主悦纳。但是，当那些你们，即\Person{欧瑟伯}及其同党，视为可靠的人，连他们也希望我召集弟兄们时，受邀方却因此动怒，便是自相矛盾了；他们本该表现得极为乐意出席才对。这些考量表明，那些被冒犯者所表现的愤怒是任性的，而拒绝出席会议的人，其行为不仅有失体统，且显得可疑。若有人看到别人做了他自己会允许的事，难道会加以责难吗？若如你们信中所写，\Quote{每一次公会议都具有不可撤销的效力，且若他的判决受到他人复查，则作出判决者便被羞辱}，那么，亲爱的弟兄们，请想一想，是谁在羞辱公会议？是谁在搁置先前审判官的裁决？目前我不想逐一细查每个案件，以免显得对某些人施加过重的压力；然而，最近发生的事例，每一个听见的人都必不寒而栗，便足以证明我所省略的其他事例了。\par

23. 那些因不虔敬而被已故的、令人怀念的\Place{亚历山大}主教\Person{亚历山大}处以绝罚的亚略派人士，不仅在各城被弟兄们谴责，更受到齐聚于\Place{尼西亚}大公会议的全体与会者的绝罚。因为他们的过犯并非寻常，他们也不是得罪了人，而是得罪了我们的主耶稣基督自己，永生天主之子。然而，这些被全世界谴责、在各教会中都留下污名的人，据说如今竟又被接纳共融；我想就连你们听到此事，也必感到愤慨。那么，究竟是谁在羞辱公会议呢？岂不是那些藐视三百位教父的投票，且喜爱不虔敬胜过虔敬的人吗？亚略狂徒们的异端，已为各地全体主教所谴责和摈弃；但主教\Person{亚大纳削}和\Person{马尔赛路}却有许多支持者，为他们发言、替他们执笔。我们已收到有利于\Person{马尔赛路}的证词，证明他在\Place{尼西亚}公会议上曾对抗那些倡导亚略教义的人；也有有利于\Person{亚大纳削}的证词，表明在\Place{提洛}时，没有任何指控坐实于他，并且在\Place{马雷奥蒂斯}，即他被说成是起草了那些指控报告的地方，他当时并不在场。亲爱的弟兄们，你们知道，\Emph{单方面}的诉讼程序毫无份量，反倒带有可疑的表相。尽管如此，事已至此，我们为了准确起见，既不偏袒你们这一方，也不偏袒为另一方执笔的人，便邀请了那些致信于我们的人前来此处；这样，由于有多人为他们执笔，便可在一次公会议中调查一切，既可使无辜者不被定罪，也可使受审者不被贸然认作清白。因此，我们并不是那羞辱公会议的一方；羞辱公会议的，乃是那些即刻且轻率地接纳了已被全体谴责的亚略派的人，他们违背了那些审判官的裁决。那些审判官中的大部分人，如今已离世，与基督同在；但仍有一些人尚在这考验的尘世生命中，当得知某些人竟搁置了他们的判决时，就感到愤慨。\par

24. 我们也从那些在\Place{亚历山大里亚}的人那里得知了以下情况。有一个名叫\Person{卡蓬}的人，因\OriginalTerm{亚略异端}{Arianism}被\Person{亚历山大}绝罚，后由一位\Person{额我略}派遣至此，与他同来的还有一些因同一异端而被绝罚的人。不过，我也从\OriginalTerm{司铎}{Presbyter}\Person{玛加略}，以及\OriginalTerm{执事}{Deacons}\Person{玛尔提留}和\Person{赫西基}那里得知了此事。因为在\Person{亚大纳削}的\OriginalTerm{司铎们}{Presbyters}抵达之前，他们催促我写信给在\Place{亚历山大里亚}的\Person{丕斯笃}，尽管当时\OriginalTerm{主教}{Bishop}\Person{亚大纳削}也在那里。当\OriginalTerm{主教}{Bishop}\Person{亚大纳削}的\OriginalTerm{司铎们}{Presbyters}到来时，他们告诉我，这个\Person{丕斯笃}是\OriginalTerm{亚略派}{Arian}，他曾被\OriginalTerm{主教}{Bishop}\Person{亚历山大}和\OriginalTerm{尼西亚大公会议}{the Council of Nicæa}绝罚，之后由一位\Person{塞恭德}祝圣，而这位\Person{塞恭德}也被\OriginalTerm{大公会议}{the great Council}以\OriginalTerm{亚略派}{Arian}绝罚。\Person{玛尔提留}及其同伴并未反驳这一说法，也未否认\Person{丕斯笃}是从\Person{塞恭德}接受祝圣的。现在请考虑，此后谁最应受责备？是我，不肯被说服写信给\OriginalTerm{亚略派}{Arian}的\Person{丕斯笃}；还是那些劝我侮辱\OriginalTerm{大公会议}{the great Council}，并将不虔敬者当作虔诚者对待的人？而且，当由\Person{欧瑟比}与\Person{玛尔提留}等人一同派来的\OriginalTerm{司铎}{Presbyter}\Person{玛加略}，在听说\Person{亚大纳削}的\OriginalTerm{司铎们}{Presbyters}所提出的反对意见之后——当时我们正期待他与\Person{玛尔提留}和\Person{赫西基}一同出现——他竟不顾身体不适，在夜间离去；这使我们推测，他的离去是因有关\Person{丕斯笃}的真相被揭露而感到羞耻。因为在\OriginalTerm{天主教会}{Catholic Church}内，\OriginalTerm{亚略派}{Arian}\Person{塞恭德}的祝圣绝不能被视为有效。这确实是对\OriginalTerm{大公会议}{the Council}及其组成的\OriginalTerm{主教们}{Bishops}的侮辱，如果他们如同在\OriginalTerm{天主}{God}面前，以极度诚心和谨慎所制定的法令，竟被当作毫无价值而弃置。\par

25. 如你们所写，若依照\Person{诺瓦替安}和\Person{萨摩撒塔的保禄}的先例，所有大公会议的法令均应有效，那么\OriginalTerm{三百位主教}{the Three hundred}的判决就更不应被推翻；确实，一个大公会议不应被少数人置若罔闻。因为\OriginalTerm{亚略派}{Arians}也是\OriginalTerm{异端者}{heretics}，正如他们一样，对双方都已作出了类似的判决。那么，经过这些大胆的行径后，是谁点燃了纷争的火焰？因为你们在信中指责我们做了这事。是我们，同情弟兄们的苦难，并在各方面依照\OriginalTerm{法典}{Canon}行事；还是他们，争闹且违反\OriginalTerm{法典}{Canon}，置\OriginalTerm{三百位主教}{the Three hundred}的判决于不顾，并在各方面侮辱大公会议？因为不仅\OriginalTerm{亚略派}{Arians}被接纳进入共融，而且\OriginalTerm{主教们}{Bishops}也惯于从一处迁往另一处。现在，如果你们真的相信所有\OriginalTerm{主教}{Bishops}拥有相同且平等的权柄，并且如你们所声称，不按城市的大小来衡量他们；那么受托管理小城的人就应留在交付给他的地方，而不应出于对其托付的轻蔑，迁移到从未交托给他的地方；轻看\OriginalTerm{天主}{God}所赐予他的，而看重人的虚浮称赞。因此，亲爱的，你们理应前来而不推辞，以便事情得以解决；因为这是理性所要求的。\par

但或许你们是被为大公会议所定的日期所阻，因为你们在信中抱怨，我们所定之日前的间隔太短。但是，亲爱的，这不过是借口。假若日期真的赶在了旅途中的任何人之前，那么所允许的间隔才会被证明是太短。但若人们不愿前来，甚至将我的\OriginalTerm{司铎们}{Presbyters}扣留至一月，这不过是那些对自己的事业毫无信心之人的借口；否则，如我先前所说，他们会前来，不顾旅程的长短，不考虑时间的短促，而只信赖自己事业的正义与合理。但或许他们是因为时局的状况而不来，因为你们又在信中宣称，我们本应考虑\Place{东方}当前的处境，而不应催促你们前来。现在，如果如你们所说，你们是因时局如此才不来，你们就应当事先考虑这样的时局，而不应成为分裂的始作俑者，也导致各教会内哀伤痛哭。但事实上，那些造成这些事端的人表明，应受指责的不是时局，而是那些不肯参加大公会议之人的顽固。\par

26. 但我也诧异，你们竟能在信中写出那一段，说是我独自写信，不是写给你们全体，而只写给\Person{欧瑟伯}及其同伙。在这抱怨中，人可发现更敏于挑剔，而非顾及事实。我收到反对\Person{亚大纳修}的信件，不是别人，正是\Person{玛尔提里乌斯}、\Person{赫西基乌斯}及其同伙送来的，我理当写信给那些曾写反对他的人。那么，要么\Person{欧瑟伯}及其同伙不该独自写信，与你们众人分开；要么你们没有收到我的信，就不该因我写给那些曾写信给我的人生气。如果我的信本该写给你们全体，你们也当与他们一同写信；但如今，我考虑到何为合理，便写信给那些曾致信于我，并将消息告知我的人。但若你们因我独自写信给他们而不悦，那么按理你们也当因他们独自写信给我而生气。然而，亲爱的，为此也有一个公平且非不通情理的理由。然而我必须让你们知道，我虽写信，我所表达的主张，并非仅是我个人的，而是整个\Place{意大利}及这些地区全体主教的。我实在不愿让他们全都写信，免得对方因人数众多而承受压力。主教们却在指定之日聚集，一致赞同这些意见，我再次写信向你们表明；因此，亲爱的，我虽独自致信你们，你们却可确信这即是所有人的主张。关于某些人为其行为提出的辩解——并非合理，而是不公且可疑的——我就说这么多。\par

27. 现在，虽然前面所述已足够表明，我们并非过于轻易或是不公地接纳了我们的兄弟\Person{亚大纳修}和\Person{马尔切洛}进入我们的共融，然而简略地将事情陈明在你们面前也是公平的。\Person{欧瑟伯}及其同伙从前曾写信反对\Person{亚大纳修}及其同伙，正如你们如今也写信了；但众多来自\Place{埃及}和其他省份的主教，却写信支持他。首先，你们反对他的信件彼此矛盾，后来的与先前的全无一致之处，在许多情形下，前者被后者答复，后者又被前者指摘。如今，信件既有此等矛盾，其中所含的陈述便毫无可信之处。其次，若你们要求我们相信你们所写的，那么按理，我们也不应拒绝相信那些写信支持他的人；尤其考虑到，你们是从远方写信，而他们身居当地，熟识其人及那里所发生的事，并以书面证实他的行事为人，且肯定地说他始终是阴谋的受害者。\par

再者，曾有一时，据说\Person{阿尔塞尼乌斯}主教被\Person{亚大纳修}害死，但我们得知他依然活着，并且与\Person{亚大纳修}保持着友谊。他已肯定地断言，在\Place{马雷奥蒂斯}起草的报告乃是\Emph{\OriginalTerm{单方面}{ex parte}}的；因为被告方\Person{玛加略}长老不在场，他的主教\Person{亚大纳修}本人也不在。此事我们不仅从他亲口得知，也从\Person{玛尔提里乌斯}、\Person{赫西基乌斯}及其同伙带给我们的报告中得知；因为我们在阅读时发现，控告者\Person{伊斯基拉斯}在场，但\Person{玛加略}和主教\Person{亚大纳修}均不在场；而且\Person{亚大纳修}的长老们欲出席，却不被允许。如今，亲爱的，若审判欲公正进行，不仅控告者，被告亦理应出席。正如在\Place{提洛}，被告\Person{玛加略}与控告者\Person{伊斯基拉斯}均曾出席，却一无所证；因此，不仅控告者应当前往\Place{马雷奥蒂斯}，被告亦当前往，好使他本人要么被定罪，要么未被定罪而能显明控告的虚假。但如今，情形却非如此，只有控告者前往，同去的还有那些\Person{亚大纳修}所反对的人，于是诉讼程序便呈现出可疑的面貌。\par

28. 他也抱怨说，那些前往\Place{马雷奥蒂斯}的人是违背他的意愿被派去的，因为他们所派的是\Person{特奥格尼乌斯}、\Person{玛里斯}、\Person{德奥多禄}、\Person{乌尔撒基乌斯}、\Person{瓦伦斯}和\Person{马切多尼乌斯}，这些人均属可疑之人。他不仅凭自己的断言，也借\Place{得撒洛尼}主教\Person{亚历山大}的信证明了这一点；他出示了\Person{亚历山大}写给主持公会议的伯爵\Person{狄约尼削}的一封信，信中极清楚地表明，当时正有一桩针对\Person{亚大纳修}的阴谋在进行。他还提交了一份真实的文件，全由控告者\Person{伊斯基拉斯}亲笔所写，其中他呼号全能天主作证：既无圣爵被打破，也无祭台被推翻，而是他被某些人收买，编造了这些指控。此外，当\Place{马雷奥蒂斯}的长老们抵达时，他们肯定地断言，\Person{伊斯基拉斯}并非天主教会的长老，而\Person{玛加略}也未曾犯下对方指控他的任何罪行。那些前来见我们的长老和执事们，也以最充分的方式支持\Person{亚大纳修}主教，竭力断言针对他的指控全非属实，他乃是阴谋的受害者。\par

埃及和利比亚的众主教都写信并声言，他的祝圣是合法且完全符合教会法规的，而你们所提出的一切指控皆属虚妄，因为既没有发生过杀人事件，也没有任何人因他之故被杀害，更没有任何杯子被打破，一切纯属捏造。不但如此，\Person{亚大纳削}主教也依据在\Place{马雷奥提斯}起草的\Emph{\OriginalTerm{片面}{ex parte}}报告指出，一位\OriginalTerm{慕道者}{catechumen}经询问后声称，当人们说亚大纳削的\OriginalTerm{司铎}{Presbyter}\Person{玛加略}闯入该地时，他正与\Person{伊斯奇拉斯}一同在里面；而其他接受询问的人则说——一位说伊斯奇拉斯当时在一间小室里——另一位则说他当时正身患疾病，躺在门后，就在他们所称玛加略到来的时刻。从他所做的这些陈述中，我们自然要问：一个患病躺在门后的人，怎么可能起身主持礼仪并献祭呢？而且，当有慕道者在内时，又如何能呈献\OriginalTerm{献礼}{Oblations}呢？因为如果慕道者尚在场，那就还不是呈献献礼的时刻。这些陈述，如我所言，皆由亚大纳削主教作出，他还依据这些报告指明，并且与他同在的人也明确证实，伊斯奇拉斯从未在\OriginalTerm{天主教会}{Catholic Church}里担任过司铎，也从未在教会的聚会中以司铎身份出现过；因为即使当\Person{亚历山大}在大公会议的宽免下接纳了那些属于\OriginalTerm{梅勒蒂乌斯分裂}{Meletian schism}的人时，\Person{梅勒蒂乌斯}也从未将他列在自己的司铎之中，正如他们所证实的；这便构成了最有力的论据，证明他甚至不是梅勒蒂乌斯手下的司铎；否则，他必定会被列入其余人之中。此外，亚大纳削也依据这些报告指明，伊斯奇拉斯在其他事上也说了谎：他曾指控，在据称玛加略闯入他们那里时焚烧了某些书籍，但他却被自己带来作证的证人证实是虚假的。\par

29. 当这些事情如此向我们陈明，有这么多证人出面为他作证，他又为自己辩白提出了如此多的理由时，我们应当做什么呢？教会的规条要求我们做什么呢？无非是不要定他的罪，反而要接纳他，并像对待主教那样待他，正如我们所做的那样。此外，除了这一切，他在这里逗留了一年零六个月，期待你们以及任何愿意前来之人的到来，他以自己的在场使每一个人都感到羞愧，因为如果他对自己的案情没有信心，他就不会在这里了；而且他并非自己主动前来，而是应我们信函的邀请，正如我们致函给你们的方式那样。然而，你们却仍然抱怨我们违犯了\OriginalTerm{教规}{Canons}。现在，请想想：是谁如此行事呢？是我们这些凭着如此充分的证据证明他清白而接纳了他的人呢，还是那些在相距三十六驿站之遥的\Place{安提约基雅}，任命一个陌生人为主教，并派军队护送他去\Place{亚历山大里亚}的人呢？这种事情，即便在\Person{亚大纳削}被放逐到\Place{高卢}时也没有做过，尽管如果当时他真的被证明犯有罪行，那时本来会这样做的。但当他回来时，自然发现他的教会仍空着并在等候他。\par

30. 而今，我不知晓这些行动是以何种借口进行的。首先，若必须说实话，当我们已写信召集会议时，任何人都不应预断会议的决定，这是不对的；其次，采取如此新奇的行动来反对教会，也是不合宜的。因为，有哪一条教会教规，或哪一项宗徒传统允许这样做：即当一个教会处于和平之中，且众多主教都与\Place{亚历山大里亚}的主教\Person{亚大纳削}意见一致时，\Person{额我略}竟被派到那里——一个对那城而言是陌生人的人，既未在那里受洗，也不为公众所知，既非司铎们、主教们、也非平信徒们所期望的——他竟在\Place{安提约基雅}受任命，并被派往亚历山大里亚，陪伴他的不是司铎，也不是本城的执事，也不是埃及的主教，而是士兵？因为来到这里的人都抱怨事情正是如此。\par

即便假设\Person{亚大纳削}在公会议后处于罪犯的地位，这项任命也不应以如此非法且违背教会规条的方式作出；而应由该省的主教们在那个教会本身，从那个司祭职、那个圣职班中祝圣一位；而承自宗徒的教规也不应这样被搁置。倘若这种冒犯是加在你们任何一个人身上，难道你们不会大声抗议，并要求对违犯教规的行为予以公正处理吗？极亲爱的诸位，我们在天主面前诚实地说，并声明，这种做法既不虔诚，也不合法，也不符合教会规矩。此外，关于\Person{额我略}进城时行为的描述，清楚地表明了他的任命是何性质。正如从\Place{亚历山大里亚}来的人所宣称的，以及主教们在信中所陈述的，在当时如此和平的时期，教会被人纵火；\OriginalTerm{童贞女}{Virgins}被剥去衣服；\OriginalTerm{隐修士}{Monks}被践踏；司铎和许多百姓被鞭打、遭受暴力；主教被投入监牢；群众被从一处拖到另一处；那些\OriginalTerm{神圣奥迹}{Mysteries}——就是他们用来指控司铎\Person{玛加略}的那事——竟被\OriginalTerm{外教人}{heathens}抓住并扔在地上；而这一切都是为了强迫某些人接受额我略的任命。这种行为清楚地显明了究竟是谁在违犯教规。如果这项任命是合法的，他就不会借助非法的手段来强迫那些以合法方式抵制他的人服从。然而尽管如此，你们却写道，在亚历山大里亚和埃及，完全和平占了上风。当然不是，除非和平的作为完全变了样，你们竟把这类行径称为和平。\par

31. 我也认为必须向你们指出以下事实，即：\Person{亚大纳修}明确声称，\Person{玛加略}在\Place{提洛}被士兵看守，而只有指控他的人陪同前去\Place{玛瑞欧提}的人；并且希望参与调查的司铎们不被允许这样做，而前述关于杯与桌案的调查，却在总督及其随从面前，并有异教徒和犹太人在场的情况下进行。起初这似乎难以置信，但根据报告证明事实确是如此；这令我们大为震惊，亲爱的诸位，我想你们也是如此。司铎们，即\OriginalTerm{奥迹}{Mysteries}的施行者，不被允许参加，而关于基督的血和基督的身体的调查，却在外教法官面前，并有\OriginalTerm{慕道者}{Catechumens}在场，甚至更糟，在异教徒和犹太人面前进行，这些人在基督教信仰方面声名狼藉。即使假定有罪行发生，也应在教会内由\OriginalTerm{圣职界}{Clergy}依法调查，而不是由憎恶\OriginalTerm{圣言}{Word}、不认识真理的异教徒调查。我相信你们和所有人都能认识到这罪的性质和严重性。关于\Person{亚大纳修}就说这么多。\par

32. 关于\Person{玛策禄}，因为你们也指控他对基督不敬，我亟愿告知你们：他在这里时，明确声明你们所写的关于他的内容并非真实；但尽管如此，在我们要求他说明他的信仰时，他本人以极大的勇气作了回答，以致我们确认他并没有持守任何违背真理的主张。他宣认了关于我们的主和救主耶稣基督的同一虔诚道理，正如天主教会所宣认的；并且他证实，他持有这些观点已经很长时间了，并非最近才采纳：实际上，先前曾出席\Place{尼西亚}大公会议的我们的司铎们，也为他的\OriginalTerm{正统信仰}{orthodoxy}作证；因为他当时就反对\OriginalTerm{亚略异端}{Arianism}，正如现在所做的一样（关于这些事，理应劝诫你们，以免你们中有人接纳这种\OriginalTerm{异端}{heresy}，而不厌恶它，因它背离了健全道理\BibleRef{弟前 1:10}）。那么，既然他宣认正统信仰，并有为其正统性的证词，我再问，在他的事上，我们除了接他为主教，如我们所做的那样，而不拒绝他于我们的\OriginalTerm{共融}{communion}之外，还应做什么呢？我写这些，不是为替他们辩解，而是为使你们相信，我们接纳这些人既合乎正义又合乎\OriginalTerm{法典}{Canon}，而你们的争论是无端的。但你们有责任尽心竭力，以一切方法纠正违反法典所发生的混乱，并确保各教会的和平；好使我们的主所赐给我们的平安\BibleRef{若 14:27}得以存留，教会不致分裂，你们也不致被指责为\OriginalTerm{分裂}{schism}的制造者。因为我承认，你们过去的行径是分裂的时机，而非和平的时机。\par

33. 因为不但主教\Person{亚大纳修}和\Person{玛策禄}及他们的同道来到这里，申诉他们所遭受的不公，还有许多其他主教，来自\Place{色雷斯}、\Place{科厄勒叙利亚}、\Place{腓尼基}和\Place{巴勒斯坦}，以及不少司铎，还有来自\Place{亚历山大里亚}和其他地方的人，出席了此地的会议；除了其他陈述之外，他们在所有与会主教面前哀叹各教会所遭受的暴力和不公，并证实，类似于在\Place{亚历山大里亚}发生的暴行，在他们自己的教会以及其他教会中也发生了。最近又有司铎带着来自\Place{埃及}和\Place{亚历山大里亚}的信函前来，他们申诉说，许多希望来参加会议的主教和司铎被阻止了；因为他们说，自\Person{亚大纳修}离开后直到此时，身为\OriginalTerm{宣信者}{confessors}的主教们被鞭打，其他人被投入监牢，而且就在最近，一些年事已高、担任\OriginalTerm{主教职}{Episcopate}已极长时间的老人，被交出去从事公共劳役，而几乎所有天主教会的\OriginalTerm{圣职界}{Clergy}连同其人民，都成了阴谋和迫害的对象。此外他们还说，某些主教和其他弟兄被流放，没有别的原因，只是为了强迫他们违心地与\Person{额我略}及其亚略同党\OriginalTerm{共融}{communion}。我们还从别人那里听到，并得到主教\Person{玛策禄}的证词证实，在\Place{加拉太}的\Place{安西拉}也发生了许多类似于在\Place{亚历山大里亚}所犯的暴行。除此之外，来参加会议的人还报告了你们中某些人（因为我不提名字）的一些如此严重的控诉，以致我拒绝将其形诸笔墨：或许你们也从别人那里听说了。正是为此原因，我才特别写信希望你们前来，好让你们在场聆听，并使一切混乱得以纠正，分歧得以愈合。那些为此目的被召唤的人，本不该拒绝，而应更加踊跃前来，免得因不来而招致对他们所受指控的怀疑，并被认为无法证明他们所写的内容。\par

34. 现在，根据这些描述，既然众教会如此受苦并受奸诈的攻击，正如我们的通报者所确证的，那么是谁点燃了纷争的火焰？是我们这些为这种状况忧伤、同情弟兄们苦难的人，还是那些制造这些事端的人？当每个教会都处在如此极度的混乱中时——这正是那些来见我们的人到这里的原因——我惊讶你们怎么能写出众教会普遍和谐的话来。这些事不是为教会的建树，而是为它的毁灭；那些以这些事为乐的人，不是和平之子，而是混乱之子：然而，我们的天主不是混乱的天主，而是平安的天主。\BibleRef{格前 14:33}因此，正如我们的主耶稣基督的天主和父所知，我是顾念你们的好名声，并祈祷众教会不陷入混乱，而能延续宗徒们所立的规范，才认为必须这样写信给你们，为叫你们最终能使那些因他们彼此间的仇视而把众教会带到这种地步的人蒙羞。因为我听说，只是极少数人是这一切的始作俑者。\par

现在，既怀着慈悲的心肠，要留意纠正，正如我先前所说，那些违反\OriginalTerm{法典}{Canon}的违规行为，以便如果已有什么祸害发生，可以借着你们的热心得到医治。也不可写什么，说我宁愿要\Person{马赛路}和\Person{亚大纳削}的共融而不要你们的，因为这类抱怨不是和平的表示，而是好争辩和仇恨弟兄的表现。为此我写了前述的那些话，为使你们明白，我们接纳他们进入共融并没有不当，这样你们或可停止纷争。如果你们曾到这里来，他们已被定罪，而且显得无法为自己的事提供合理的证据，那么你们这样写信是好的。但事实是，如我先前所说，我们与他们保持共融是符合法典的，并非不当，因此我因基督的缘故恳求你们，不要使基督的肢体被撕裂，也不要信赖偏见，反而寻求主的平安。为了满足少数人的意气，而拒绝从未被定罪的人，并因此使圣神忧郁，这既不圣善也不正义。\BibleRef{弗 4:30}但若你们认为能指证他们什么，并当面驳倒他们，你们中愿意的可以到这里来；因为他们也保证，他们已准备就绪，要完全证实他们向我们报告的那些事的真实性。\par

35. 因此，亲爱的各位，请通知我们此事，好使我们能写信给他们以及那些将要再次聚集的主教们，以便被告能在众人面前被定罪，混乱不再在教会中蔓延。已经发生的事已足够；确实，主教们在主教面前被判放逐，这已经是够了；对此我不宜多谈，免得我似乎过于苛责那些当时在场的人。但如果必须说真话，事态本不该发展到如此地步；他们的小气私忿也不应被容忍到如今的程度。就算我们承认，如你们所写，\Person{亚大纳削}和\Person{马赛路}被\Quote{调离}他们的地方，那么对于我从各处来的其他主教和长老们的情形，又该怎么说呢？他们抱怨自己也被强行驱逐，并遭受了类似的伤害。啊，亲爱的，教会的决定不再符合福音，而只趋向放逐和死亡。假如，如你们所断言的，那些人确有某罪过，那么针对他们的案件办理，不应采取这种方式，而应根据教会的法典。本应写信告知我们全体，以便公正的判决能从全体发出。因为这些受害者是主教，是绝非普通的教会，而是那些宗徒们曾亲自治理的教会。\par

而且为什么有关\Place{亚历山太}的教会，尤其没有告知我们？你们难道不知，惯例是先写信给我们，然后从此处作出公正的判决吗？因此，如果那个地方的主教有任何此类嫌疑，理应将此告知本地的教会；然而，他们既忽略了通知我们，又按自己的喜好自作主张，现在却想要我们同意他们的决定，尽管我们从未定他的罪。\Person{保禄}的规条不是这样，教父们的传统也不是这样指示的；这是另一种行事方式，是一种新的做法。我恳请你们，耐心地容忍我：我所写的是为了公益。因为我从真福宗徒\Person{伯多禄}所领受的，我向你们表明；如果不是这些行径如此地惊扰了我们，我本不会写这些，以为这些事情对众人都是显而易见的。主教们被强行赶离自己的教座，被放逐，而来自其他地方的其他人被任命接替；另有些人被阴险地攻击，以致民众不得不为那些被强行带走的人悲伤，至于那些被派来接替的人，他们被迫放弃寻找他们想要的人，而接纳他们不想要的。\par

我请求你们，使此类事情不再发生，并且你们要以书面谴责那些企图这样做的人；使教会不再如此受苦，不再有任何主教或长老受凌辱，不再有人被迫违背自己的判断行事，正如他们向我们所述的那样，免得我们成为外教人的笑柄，尤其是，免得我们招惹天主的义怒临到我们身上。因为我们每个人，在审判之日，\BibleRef{玛 12:36}都要为今生所做的一切交账。愿我们所有人都怀有天主的心意！好使教会能恢复她们自己的主教，并在我们的主耶稣基督内永远喜乐；愿光荣借着祂归于圣父，至于无穷之世。阿们。\par

我在主内为你们的安康祈祷，至亲爱的弟兄们，我深切思念你们。\par

36. \OriginalTerm{罗马公会议}{Council of Rome}由\Place{罗马}的主教\Person{犹利乌斯}如此记下。\par

% structure: p0072 source=h2 target=subsection reason=relative-heading-rank; base=h2

\subsection{第三章 \Emph{\OriginalTerm{撒底加公会议}{Council of Sardica}致\Place{埃及}及\Place{亚历山大里亚}各教会以及所有教会的信函}}

然而，尽管这样，\Person{欧瑟比乌斯}及其同伙仍无耻地行事，扰乱各教会，并图谋毁灭许多人。最虔诚的皇帝们\Person{君士坦提乌斯}和\Person{君士坦斯}得知此事，命令来自西方和东方的主教们到\Place{撒底加}城聚会。在此期间\Person{欧瑟比乌斯}死了；但从各地聚集了一大批人，我们要求\Person{欧瑟比乌斯}的同伙及其党羽接受审判。但他们，鉴于自己过去所行之事，并察觉控告他们的人已来到大公会议，便不敢这样做；反而，当其余所有人都怀着诚实的意图聚集时，他们又带来了\OriginalTerm{侍卫官}{Counts}\Person{穆索尼亚努斯}和\OriginalTerm{御营官}{Castrensian}\Person{赫西基乌斯}，为的是按他们的惯例，利用他们的权威达到自己的目的。但当大公会议在没有侍卫官出席的情况下召开，且不许任何士兵在场时，他们便惶恐不安，良心受到谴责，因为他们再也无法获取他们所想要的判决，只能得到理性与真理所要求的判决。然而我们一再提出挑战，众主教的大公会议要求他们出来，说：\Quote{你们来此本为接受审判，为何现在却要逃避？要么你们就不该来，来了就不该躲藏。这种行为将成为你们最大的定罪。看，\Person{亚大纳削}及其同伙就在这里，你们趁他们不在时控告他们；因此，如果你们认为自己对他们有任何指控，可以当面定罪他们。但如果你们假装不愿这样做，而其实是不能，你们就清楚地表明自己是诽谤者，这便是大公会议要给你们的判决。}他们听到这话，自己定罪了自己（因为他们意识到他们对我们的阴谋和捏造），羞于露面，从而证明自己犯了许多卑劣的诽谤罪。\par

因此，圣公会议谴责了他们不合礼节且令人怀疑的逃跑，并准许我们为自己辩护；当我们陈述了他们对我们所行的，并且通过见证人和其他证据证实了我们陈述的真实性后，他们满心惊讶，所有人都承认我们的对手害怕面对公会议是有充分理由的，唯恐他们的罪行当面被证实。他们还说，或许这些人从东方来，原以为\Person{亚大纳削}及其同伙不会出现，但看到他们对自己的案件充满信心，并要求审判时，他们就逃跑了。于是，他们接纳我们这些被诬告的受害者，并且更加坚定地与我们共享情谊和爱。然而，他们罢免了\Person{欧瑟比乌斯}那些作恶的同党——他们甚至比欧瑟比乌斯更无耻，即\Place{赫拉克莱亚}的\Person{德奥多鲁斯}、\Place{尼洛尼亚斯}的\Person{纳尔基苏斯}、\Place{凯撒勒亚}的\Person{阿卡丘斯}、\Place{安提约基}的\Person{斯特法努斯}、\Place{潘诺尼亚}的\Person{乌尔撒丘斯}和\Person{瓦伦斯}、\Place{厄弗所}的\Person{梅诺凡图斯}，以及\Place{劳狄刻雅}的\Person{乔治}；他们并致函世界各处的主教，以及每一位受害者的教区，内容如下。\par

\Emph{\OriginalTerm{撒底加公会议}{Council of Sardica}致\Place{亚历山大里亚}教会的信函}\par

\OriginalTerm{圣公会议}{Holy Council}，因\OriginalTerm{恩宠}{grace}在\Place{撒底加}聚集，包括来自\Place{罗马}、\Place{西班牙}、\Place{高卢}、\Place{意大利}、\Place{坎帕尼亚}、\Place{卡拉布里亚}、\Place{阿普利亚}、\Place{非洲}、\Place{撒丁尼亚}、\Place{潘诺尼亚}、\Place{默西亚}、\Place{达基亚}、\Place{诺里库姆}、\Place{西斯西亚}、\Place{达尔达尼亚}、\Place{另一个达基亚}、\Place{马其顿}、\Place{德撒里}、\Place{阿卡亚}、\Place{厄皮鲁斯}、\Place{色雷斯}、\Place{罗多佩}、\Place{巴勒斯坦}、\Place{阿拉伯}、\Place{克里特}和\Place{埃及}的代表，致亲爱的弟兄们，即\OriginalTerm{司铎}{Presbyters}和\OriginalTerm{执事}{Deacons}，以及所有住在\Place{亚历山大里亚}的\OriginalTerm{天主的圣教会}{Holy Church of God}，在主内致候。\par

37. 我们并非不知，而是在收到你们敬诚的信函之前就清楚知道，可憎的\OriginalTerm{亚略派}{Arians}\OriginalTerm{异端}{heresy}的拥护者们正在施行许多危险的阴谋，这与其说伤害教会，不如说毁灭他们自己的灵魂。因为这始终是他们无原则诡计的目的；他们不断从事的正是这种致命的企图，即如何以最佳方式驱逐并迫害任何地方拥有\OriginalTerm{正统}{orthodox}信仰、并持守由\OriginalTerm{教父}{Fathers}传授给他们的\OriginalTerm{天主教会}{Catholic Church}教义的人。对一些人，他们捏造虚假指控；对另一些人，他们驱逐流放；对又一些人，他们以所施加的刑罚加以毁灭。无论如何，他们企图以暴力和专制突袭我们弟兄兼同侪主教\Person{亚大纳削}的清白，因此在对他的案件进行审查时，毫无谨慎、毫无信实、毫无任何正义。因此，他们对自己在那次事件中所扮演的角色毫无信心，对他们所散播的关于他的流言蜚语也毫无信心，反而察觉自己无法就案件提出任何确凿的证据，于是当他们来到\Place{撒底加}城时，便不愿与全体圣主教的公会议会面。由此显明，我们的弟兄兼同侪主教\Person{犹利乌斯}的裁决是公正的；因为他经过审慎的考量与关注后断定，我们丝毫不可犹豫与我们弟兄\Person{亚大纳削}\OriginalTerm{共融}{communion}。因为他拥有八十位主教的可信见证，并且还能提出这样一个公平的论点来支持他，即仅仅通过我们至亲爱的弟兄——他自己的司铎们——以及通信，他就挫败了\Person{欧瑟比乌斯}及其同伙的计谋，他们更多依赖暴力而非司法审查。\par

因此，各地的主教们一致认定，应因\Person{亚大纳削}无罪而与他保持共融。也让你们的爱德注意，当他来到在\Place{撒底加}召开的神圣大公会议时，正如我们先前所说，我们既通过书信，也通过口头的命令，将此事告知了东方的主教们，并邀请他们出席。然而，他们因受自己良心的谴责，便以不当的借口推脱，并设法逃避审查。他们要求将一个无辜的人作为罪犯从我们的共融中予以弃绝，却不考虑此举是何等不当，甚至可以说是何等地不可能。至于那些由某些最为邪恶、最为堕落无耻的青年在\Place{玛琉提斯}炮制的调查报告——这些人连教职中最卑微的职务都不可委托——，可以确知，它们都是\Emph{\OriginalTerm{单方面}{ex parte}}的陈述。因为我们的弟兄\Person{亚大纳削}主教当时并不在场，被他们控告的\Person{玛加略}司铎也不在场。此外，他们的调查，或更确切地说，他们对事实的歪曲，伴随着最可耻的行径。他们时而讯问外教人，时而讯问慕道者，目的不是让他们陈述自己所知道的事，而是让他们说出别人教给他们的谎言。而你们司铎们——在你们主教缺席期间负责事务的人——希望到场参加调查，好能揭示真理，驳斥谎言，但他们却丝毫不尊重你们；他们不允许你们到场，反而以侮辱的方式将你们赶走。\par

虽然这些情况已将他们的诽谤在众人面前清晰地揭露出来；但我们在阅读调查报告时也发现，那个极其邪恶的\Person{依斯希拉}——他因这诬告而从他们那里获得了主教的虚衔作为赏报——已自己证明了自己的诬告。他在调查报告中声称，正当他肯定地说\Person{玛加略}进入他的小室时，他正卧病在床；然而\Person{欧瑟比乌斯}及其同伙竟胆敢写道，当\Person{玛加略}进入时，\Person{依斯希拉}正站着献祭。\par

38. 他们随后对他提出的卑劣而诽谤的指控，已为众人所周知。他们大声鼓噪，断言\Person{亚大纳削}犯了谋杀罪，杀害了一位名叫\Person{阿尔瑟尼乌斯}的\OriginalTerm{梅勒廷派}{Meletian}主教，并假装悲痛，以虚假的哀悼和伪装的泪水来哀吊他的丧亡，且要求将那活人如同死人一般的身体交给他们。然而他们的欺诈并未不被察觉；大家全都知道那人活着，并且算是活人之一。当这些一有机会就伺机而动的人察觉到他们的谎言被揭穿后（因为\Person{阿尔瑟尼乌斯}显现自己活着，从而证明他既未被杀害，也并未死亡），他们却仍不肯罢休，反而在原先的诽谤之外又添加其他诽谤，并采取新的手段来诋毁此人。然而，我们亲爱的弟兄\Person{亚大纳削}并未惊惶，在此案中他再次以极大的勇气邀请他们来验证，我们也祈祷并劝勉他们前来对质，如果他们能够，就证实他们对他的指控。何等的狂傲！何等可怕的骄傲！或者更确切地说，如果必须讲真话，是何等的邪恶和控告良心的！因为这就是众人都持有的看法。\par

因此，亲爱的弟兄们，我们劝诫并敦促你们，首要的是保持天主教会正确的信仰。你们经历了众多严峻而沉重的考验；天主教会蒙受了诸多的侮辱和伤害，但\BibleQuote{唯独坚持到底的，才可得救。} \BibleRef{玛 10:22}，所以，即使他们仍然肆无忌惮地攻击你们，你们也当以你们的磨难为喜乐。因为这样的苦难是一种殉道，你们这样的宣认和折磨不会没有赏报，你们必将由天主那里获得奖赏。因此，你们要全力维护纯正的信仰，以及你们的主教、我们同工\Person{亚大纳削}的清白。我们也没有沉默，也未忽略关切你们的安慰，而是商议并实践了爱德所要求的一切。我们与受苦的弟兄们同心同感，将他们的苦难视如己出。\par

39. 因此，我们已致函恳请我们最虔敬和虔诚的皇帝们，恳请他们的仁慈下令释放那些仍遭受磨难和压迫的人，并命令那些职责仅限于审理民事案件的地方官员，不得对圣职人员作出审判，今后也不得以照顾教会为借口，在任何方面对弟兄们有所企图；而是愿每个人都能如同他所祈祷和渴望的那样，免于迫害、暴力和欺诈，在安宁与和平中遵循大公及宗徒的信仰。至于\Person{额我略}，他享有由异端者非法委任的名声，并被他们派往你们的城市，我们希望你们的全体一致知晓，他已被整个神圣大公会议的一项判决予以罢免，尽管事实上他从未在任何时候被视为真正的主教。因此，欢欢喜喜地接待你们的主教\Person{亚大纳削}吧，因为为此目的我们已平安地将他遣回。我们也劝勉所有那些或出于恐惧、或由于某些人的阴谋而与\Person{额我略}共融的人，现在经过我们的忠告、劝勉和说服后，要从他那可憎的共融中退出，并立即归向天主教会。\par

40. 但鉴于我们得知，我们的同僚\OriginalTerm{司铎}{Presbyters} \Person{阿弗托尼乌斯}、\Person{卡皮托}之子\Person{亚大纳削}、\Person{保禄}和\Person{普鲁提奥}，也因\Person{优西比乌斯}及其同党的阴谋而受苦，以致他们中有人遭受流放之考验，有人为生命安危而逃亡，我们因此认为有必要让你们知道此事，好叫你们明白，我们也已接纳并宣判他们无罪，深知\Person{优西比乌斯}及其同党对正统派所做的一切，都趋向于使受他们攻击者得到荣耀和称赞。你们的主教、我们的弟兄\Person{亚大纳削}理应将他们的事告知你们，向自己人说明自己人的事；但为了更充足的见证，他愿这神圣的公会议也给你们写信，我们并未迟延，而是赶紧向你们传达此事，好使你们如同我们一样接纳他们，因为他们也堪受赞美，由于他们对基督的虔诚，竟被异端者施以暴力。\par

至于神圣的公会议对那些为首倡导亚略异端、曾得罪你们及其余各教会的人所作出的判决，你们将从附呈的文件中得知。我们已将文件送达你们，好叫你们从中明白，\OriginalTerm{天主教会}{Catholic Church}必不忽略那些冒犯她的人。\par

\Emph{\Place{萨底卡}公会议致\Place{埃及}和\Place{利比亚}主教书。}\par

神圣的公会议，赖天主的恩宠在\Place{萨底卡}聚集，致\Place{埃及}和\Place{利比亚}的主教们，他们同为牧职者及至爱弟兄，在主内致候。\par

41. 我们并非不知——其实在我们接获你们虔诚书信之前便已清楚知道——那可憎的亚略异端的支持者们正行许多危险的阴谋，与其说是损害教会，不如说是毁灭他们自己的灵魂。这正是他们奸计与邪恶的一贯目的：他们不断从事的致命图谋，即如何最好地将所有持有正统观点、并坚持\OriginalTerm{天主教会}{Catholic Church}由教父们传下来之教义的人，从他们的位所驱逐并加以迫害。他们对有些人捏造假指控；将另一些人逐入流放；再有一些人则因施加的刑罚而遭毁灭。无论如何，他们试图以暴力和专制突然袭击我们的弟兄和同僚主教\Person{亚大纳削}的清白，因此他们审理他的案件时毫无诚信与任何公允可言。因此，他们既对自己在那场合扮演的角色，乃至对自己散布的针对他的报告毫无信心，又察觉到无法就案件提出任何确凿证据，当他们来到\Place{萨底卡}城时，便不愿与全体神圣主教组成的公会议会面。由此可见，我们的弟兄和同僚主教\Person{尤利乌斯}的裁决是公正的；因为他经过谨慎审议和关切后断定，我们丝毫不该犹豫与我们弟兄\Person{亚大纳削}保持共融。他不但拥有八十位主教的可信见证，还能提出以下合理论据支持他：仅凭借我们至爱的弟兄，也就是他自己的\OriginalTerm{司铎}{Presbyters}们，以及通信往来，就挫败了\Person{优西比乌斯}及其同党的图谋，这些人倚靠暴力远胜于司法审查。\par

因此，各地的主教们基于他的无辜而决定与\Person{亚大纳削}保持共融。也让你们的爱德注意到，当他来到在\Place{萨底卡}聚集的神圣公会议时，正如我们先前所说，东方的主教们既通过书信，又通过口头传达的命令得知此事，并受我们邀请出席。但他们被自己的良心谴责，便求助于不合宜的托辞，开始回避审查。他们要求将一个无辜者作为犯人从我们的共融中排斥，而不考虑这样的程序多么不合宜，甚至多么不可能。至于在\Place{马雷奥蒂斯}由某些极邪恶无耻的青年人所炮制的报告——这些人即使交给他们圣职中最卑下的职分也不可信任——这些报告无疑是\Emph{单方面}的陈述。因为我们的弟兄\Person{亚大纳削}主教既未在场，被他们指控的\Person{马卡里乌斯}\OriginalTerm{司铎}{Presbyter}也未在场。况且，他们的调查，或不如说对事实的捏造，伴随着最可耻的情节。有时审查的是异教徒，有时是慕道者，不是要他们说出自己所知道的，而是要他们说出别人教给他们的谎言。当你们这些\OriginalTerm{司铎}{Presbyters}——在你们主教不在时负责管理——希望到场参加调查，以便能展示真相、驳斥谎言时，他们根本不理会你们；不许你们在场，反而将你们侮辱地赶走。\par

如今尽管他们的诽谤藉着这些情形已在众人面前最显明地暴露出来；但我们阅读报告后也发现，那最不义的\Person{依斯基拉斯}——他从他们那里获得了空虚的主教头衔作为虚假指控的赏报——已自证其诬告之罪。他在报告中宣称，据他确切断言，\Person{马卡里乌斯}进入他的小屋时，他正卧病在那里；然而\Person{优西比乌斯}及其同党竟大胆写道：当\Person{马卡里乌斯}进来时，\Person{依斯基拉斯}正站着行祭献。\par

42. 他们所捏造的卑鄙污蔑的指控，如今已为众人所共知。他们大声喧哗，声称\Person{亚大纳削}犯有谋杀罪，杀害了一位名叫\Person{阿尔塞纽斯}的\OriginalTerm{梅勒提派}{Meletian}主教，并虚伪地悲叹其失落，假惺惺地流泪，要求将一个活人的身体当作死人交出。然而他们的骗局并未得逞；众人皆知此人尚在人世，且名列生者之中。当这些一有机会便肆意妄为的人发现谎言已被揭穿（因为\Person{阿尔塞纽斯}现身于世，证明自己并未被害，没有死），他们仍不肯罢休，反而更添新的诽谤，另设诡计来污蔑此人。然而，我们亲爱的弟兄\Person{亚大纳削}并未慌乱，再次在此事上以大无畏的精神挑战他们拿出证据；我们亦一同祈祷并劝诫他们前来对质，若有能力，便证实对他的指控。何其狂妄！何其可憎的骄傲！更确切地说，若须道出真相，实乃邪恶而责难人的良心！这正是众人对此事的看法。\par

因此，亲爱的弟兄们，我们劝勉你们，最要紧的是坚持天主教会的正确信德。你们曾经历许多严峻而痛苦的考验；天主教会遭受了许多侮辱和伤害，但\BibleQuote{唯独坚持到底的，才可得救。}\BibleRef{玛 10:22}。为此，即使他们仍会肆无忌惮地攻击你们，你们的磨难也成为你们的喜乐。因为这种苦难是一种殉道，像你们这样的宣认和折磨不会没有赏报，你们必将从天主那里领受赏赐。所以，你们要全力维护健全的信德，以及你们的主教、我们的弟兄\Person{亚大纳削}的清白。我们也没有沉默，没有忽视你们的安慰，而是商议并做了爱德所要求的一切。我们同情受苦的弟兄们，将他们的磨难视为自己的磨难，我们的眼泪与你们的混合在一起。弟兄们，你们并非唯一的受苦者：我们中的许多弟兄，就是与我们同工的人，也来到此处，为此事痛苦哀哭。\par

43. 因此，我们已写信恳请我们最虔敬、至圣的皇帝陛下，期望他们仁慈地下令释放那些仍处于患难和压迫中的人，并命令那些只应处理民事案件的官吏不得审判圣职人员，从此以后也不得以供应教会为借口，对兄弟们有任何企图，而是让每个人都能按照自己的祈祷和愿望生活，免受迫害、暴力和欺诈，并在平静与和平中遵循公教及宗徒的信德。至于\Person{额我略}，他被人视为由异端者非法任命，并被他们派到你们的城市，我们希望你们一致认识到，他已被整个神圣大公会议的判决所罢免，尽管事实上他从未在任何时候被视为真正的主教。因此，欣然迎接你们的主教\Person{亚大纳削}吧；因为我们正是在此和平中将他打发回来。我们也劝勉所有曾因恐惧或因某些人的阴谋而与\Person{额我略}共融的人，在受到我们的告诫、劝勉和说服后，脱离他那可憎的共融，立即与天主教会联合。\par

神圣大公会议针对\Person{德奥多鲁斯}、\Person{纳尔基索斯}、\Person{斯德望}、\Person{阿卡基乌斯}、\Person{梅诺方图斯}、\Person{乌尔撒齐乌斯}、\Person{瓦伦斯}和\Person{若尔日}——这些人是亚略异端的首领，曾对你们及众教会犯下罪过——所作出的判决，你们将从随附的文件中得知。我们已将文件寄给你们，以便你们的虔敬能赞同我们的决定，并从中理解到，天主教会绝不会忽视那些冒犯她的人。\par

\Emph{\Place{撒尔底加}大公会议通函}\par

圣大公会议，赖天主的恩宠，在\Place{撒尔底加}聚集，致书于普世天主教会内最亲爱的弟兄、主教及同工，谨在主内致以安好。\par

44. 那些亚略派的疯人曾屡次胆敢攻击维护正统信德的天主的仆人；他们企图代之以虚假的教义，并驱逐正统派；最后，他们对信德发起了如此猛烈的攻击，以至于连我们皇帝陛下的虔敬也得知了此事。于是，靠天主恩宠的助佑，我们皇帝陛下亲自将我们从不同的省份和城市召集起来，并准许在\Place{撒尔底加}城举行这次神圣大公会议，为要消除一切分歧，驱除一切虚假教义，使基督的虔敬得以由所有人共同维护。东方的众主教也应皇帝陛下的劝勉出席了会议，主要是因为他们曾屡次散布关于我们亲爱的弟兄及同工\Place{亚历山大里亚}主教\Person{亚大纳削}和\Place{安基拉—迦拉达}主教\Person{马尔塞鲁斯}的谣言。他们的诽谤或许已经传到你们耳中，也许他们曾企图扰乱你们的听闻，好让你们相信他们对无辜者的指控，并从你们心中抹去对他们自身邪恶异端的任何怀疑。然而，他们并未能大肆得逞；因为主是祂众教会的护佑者，祂为我们众人并为他们舍命，并藉着祂自己为我们众人开辟了通往天堂的道路。因此，当\Person{欧瑟比乌斯}及其同党早先写信给我们的弟兄、\Place{罗马}教会主教\Person{犹利乌斯}，反对前述的弟兄，亦即\Person{亚大纳削}、\Person{马尔塞鲁斯}和\Person{阿斯克勒帕斯}时，其他地区的主教也写信作证，证明我们的同工\Person{亚大纳削}的清白，并声明\Person{欧瑟比乌斯}及其同党所陈述的，无非是虚假与诽谤。\par

事实上，他们的\OriginalTerm{诽谤}{calumnia}由以下事实清楚证明：当我们至爱同僚\Person{尤利乌斯}邀请他们参加一次\OriginalTerm{公会议}{concilium}时，他们竟不前来；而\Person{尤利乌斯}本人写给他们的信也证明此事。因为倘若他们对自己针对我们弟兄所采取的手段和行为有信心，他们早就来了。此外，他们在这次伟大而神圣的公会议上的表现，更加明显地证明他们的阴谋。因为他们抵达\Place{撒狄卡}城，见到我们的弟兄\Person{亚大纳削}、\Person{马尔塞鲁斯}、\Person{阿斯克列帕斯}等人之后，却不敢前来受审；尽管一再受邀出席，他们仍不服从召唤。虽然我们所有主教齐聚一堂，尤其是那位享有至福高龄的\Person{何西乌斯}——以其年龄、\OriginalTerm{宣认}{confessio}以及所经历的诸多劳苦，堪受一切尊敬——虽然我们等候并力劝他们前来受审，以便在众同僚面前，就他们于弟兄缺席时散布及写信指控的罪名，证实其真实性；然而，正如我们先前所言，他们受此邀请时仍不愿前来，从而证实了他们的诽谤，并藉此拒绝几乎向世界宣布了他们所从事的阴谋与密谋。凡对自己断言的真相有信心者，必能当着对手的面证实其言。但他们既不肯与我们相见，我们以为，如今无人能再怀疑——无论他们可能再次借助其恶行——他们对自己的同僚毫无证据，只在同僚缺席时诽谤他们，却回避面对他们。\par

45. 亲爱的弟兄们，他们逃跑，不仅因为自己散布的诽谤，更因为他们看到那些带着种种指控前来控告他们的人已经到场。因为他们使用过的锁链和铁铐被拿了出来；从流放中返回的人出现了；我们的弟兄也来了，他们是那些仍被扣留在放逐中之人的亲属，以及那些因他们的手段而丧生者的朋友。最具分量的指控是，有主教在场，其中一位拿出了他们迫使他佩带的铁铐锁链，其他人则诉说到因他们的诽谤而导致的死亡。他们竟疯狂到如此地步，甚至企图杀害主教；若非主教们逃脱其手，他们早已得逞。我们可敬的同僚，已故的\Person{狄奥杜鲁斯}，因他们的诬告而正在逃亡时去世；正是因这些诬告，官方下达了处死他的命令。另一些人展示了刀伤；还有一些人抱怨说，自己因他们的手段而饱受饥饿之苦。为这些事作证的，并不是普通人，而是整个教会，有\OriginalTerm{特使}{legatus}代表他们出面，向我们述说士兵手持刀剑、暴徒挥舞棍棒、法官的威胁、伪造信件的行径。因为他们宣读了一些由\Person{狄奥格尼乌斯}及其同伙伪造、用以鼓动皇帝敌视我们同僚\Person{亚大纳削}、\Person{马尔塞鲁斯}和\Person{阿斯克列帕斯}的信件；而当时曾是\Person{狄奥格尼乌斯}的\OriginalTerm{执事}{diaconus}的人证实了此事。从这些人那里，我们听说有贞女被剥光衣服、教堂被焚烧、圣职人员被拘押——这一切无非是为了那受诅咒的\OriginalTerm{亚略}{Arian}派狂人的\OriginalTerm{异端}{haeresis}；凡拒绝与之\OriginalTerm{共融}{communio}者，都被迫遭受这些事。\par

当他们看清了事情的状况后，便陷入困境，不知该选择哪条路。他们羞于承认自己的所作所为，却再也不能掩饰下去。于是他们来到\Place{撒狄卡}城，想借到场之机，显得自己似乎可以消除对这些罪行的嫌疑。但当他们见到自己诽谤过的人，以及那些曾遭他们毒手的人；当他们面对控告者以及他们的罪证时，就不愿上前来，尽管我们的同僚\Person{亚大纳削}、\Person{马尔塞鲁斯}及\Person{阿斯克列帕斯}一再邀请，这几位同僚带着极大的自由控诉他们的行径，并力劝、挑战他们前来受审，承诺不仅要驳斥他们的诽谤，还要提出他们对自己教会所犯罪行的证据。但他们被良心的恐惧所攫，以致逃跑；他们这样做，就暴露了自己的诽谤，并借逃跑承认了自己所犯的罪行。\par

46. 但是，尽管他们的恶意与诽谤在此次以及先前的场合已显露无遗，但为避免他们因逃跑而再施恶计，我们认为宜依照真理原则审查他们所扮演的角色；这正是我们的意图，而我们发现，他们以其行为成了诽谤者，不过是针对我们在圣职上的弟兄们的一个阴谋的策划者。因为，他们声称被\Person{亚大纳削}杀害的\Person{阿尔塞尼乌斯}，仍然活着，被列入活人之列；由此我们可推断，他们在其他事上所散布的传闻也是捏造的。再者，他们散布谣言，提及一只\OriginalTerm{圣爵}{calix}，说是被\Person{亚大纳削}的\OriginalTerm{司铎}{presbyter}\Person{玛加略}打碎的；但从\Place{亚历山大里亚}、\Place{马雷奥蒂斯}及其他地方来的人，都证实从未发生过此类事情。而且，那些致信我们同僚\Person{尤利乌斯}的埃及主教们，明确肯定他们中间甚至连对此事的丝毫怀疑都不曾有过。\par

\Emph{\OriginalTerm{单方面陈述}{ex parte}}；甚至在这些报告的形成过程中，他们还审问了外教人和慕道者；其中一位慕道者在受审时说，\Person{玛加略}闯入他们房间时他正在场；另一人则声称，他们常谈及的那位\Person{依斯基拉斯}那时正卧病在床；由此可见，\OriginalTerm{圣事}{Mysteries}从未举行，因为慕道者当时在场，而且\Person{依斯基拉斯}并不在场，而是卧病在床。况且，这个最卑劣的\Person{依斯基拉斯}，曾妄称（且已被定罪）\Person{亚大纳削}焚烧了一些圣书，但他本人已承认，\Person{玛加略}来时他正卧病在床；由此显见他是个诽谤者。然而，作为对他这些谗言的酬报，他们竟授予这\Person{依斯基拉斯}主教头衔，尽管他连司铎都不是。因为两位司铎——他们曾与\Person{梅勒提乌斯}来往，后来由真福\Place{亚历山卓}的主教\Person{亚历山大}接纳，如今与\Person{亚大纳削}在一起——来到大公会议前，作证说他连\Person{梅勒提乌斯}的司铎都不是，且\Person{梅勒提乌斯}在\Place{马雷奥蒂斯}从未有过教会或圣职人员。可是，这个从未担任过司铎的人，如今竟被他们推为主教，企图借此名号在听他诽谤之人中占得上风。\par

47. 我们的共同司牧者\Person{马尔塞禄}的著作也被宣读，藉此\Person{欧瑟比乌斯}及其同伙的欺诈昭然若揭。因为\Person{马尔塞禄}以探寻方式提出的观点，被他们歪曲为他所宣认的意见；但当著作的后半部分以及这些提问之前的段落也被宣读之后，他的信仰被证实纯正。他从未如他们所断言的那般，声称\OriginalTerm{圣言}{Word of God}以圣玛利亚为起始，也未曾声称祂的国度有终结；恰恰相反，他写道，祂的国度既无起始，亦无终结。我们的共同司牧者\Person{阿斯克勒帕斯}也出示了在\Place{安提约基}当着控告者和\Place{凯撒勒亚}的\Person{欧瑟比乌斯}之面拟就的报告，并通过审判其案件的众主教的声明，证明了自己的清白。因此，亲爱的弟兄们，他们完全有理由不理睬我们屡次的召唤，并逃离大公会议。他们是被自己的良心驱使至此；而他们的逃亡，反而证实了他们的诽谤行径，使得在场控告者所陈述和辩说的那些事，令人信以为真。除此之外，他们不仅收留了先前因\OriginalTerm{亚略异端}{Arian heresy}而被免职和驱逐的人，甚至还将他们提升到更高的位置，把执事升为司铎，由司铎中拣立主教，其目的无他，就是为了散播邪恶，败坏正统信仰。\par

48. 继\Person{欧瑟比乌斯}及其同伙之后，他们的领袖如今是\Place{赫拉克勒亚}的\Person{德奥多鲁斯}、\Place{基里基雅}的\Place{内罗尼亚斯}的\Person{纳尔齐苏斯}、\Place{安提约基}的\Person{斯德望}、\Place{劳狄刻雅}的\Person{乔治}、\Place{巴勒斯坦}的\Place{凯撒勒亚}的\Person{阿卡齐乌斯}、\Place{亚细亚}的\Place{厄弗所}的\Person{梅诺凡图斯}、\Place{默西亚}的\Place{辛吉杜努姆}的\Person{乌尔撒齐乌斯}，以及\Place{潘诺尼亚}的\Place{穆尔萨}的\Person{瓦伦斯}。这些人不准随他们从东方来的人参加神圣大公会议，甚至不许他们接近\OriginalTerm{天主的教会}{Church of God}；反而在前往\Place{萨迪卡}的途中，他们自己在各地召开会议，并胁迫约定，等到了\Place{萨迪卡}，他们不仅不会出席审判，连神圣大公会议的集会也不参加，只消形式上抵达并通报到达后，便立即逃走。从我们的共同司牧者、\Place{巴勒斯坦}的\Person{玛加略}和\Place{阿拉伯}的\Person{阿斯泰里乌斯}那里，我们得以确证这事。这两位原本与他们同行，后来脱离了他们的不信。他们来到大公会议，申诉所遭受的暴力，并说那些人毫无公正之举；还补充说，他们中间有许多人持守正统，却被那些人用威逼利诱的手段阻止前来，这些手段专用来阻挠想与他们分开的人。正因如此，那些人急切地要求所有人同住一处，甚至不许他们独自待上哪怕片刻。\par

49. 因此，既然我们不应缄默不言，也不该忽视他们的诽谤、囚禁、谋杀、伤害、利用假信函密谋、暴行、剥去贞女衣服、放逐、破坏教会、纵火焚烧、将主教从小城调往大教区，尤其是他们借机兴起的恶名昭彰的\OriginalTerm{亚略异端}{Arian heresy}对抗正统信仰；我们遂宣告我们亲爱的弟兄和共同司牧者\Person{亚大纳削}、\Person{马尔塞禄}、\Person{阿斯克勒帕斯}，以及与他们一同服事主的人，清白无罪，并已致函各个教区，使各教会的信友得知自己主教的清白，视他为主教并期待他的归来。\par

至于那些像狼一样侵入他们教会的人，即\Place{亚历山卓}的\Person{额我略}、\Place{安基拉}的\Person{巴西略}和\Place{迦萨}的\Person{昆提亚努斯}，既不可称他们为主教，亦不可与之有任何共融，不可收他们的信函，亦不可写信给他们。至于\Person{德奥多鲁斯}、\Person{纳尔齐苏斯}、\Person{阿卡齐乌斯}、\Person{斯德望}、\Person{乌尔撒齐乌斯}、\Person{瓦伦斯}、\Person{梅诺凡图斯}和\Person{乔治}，虽然最后一位因恐惧未从东方前来，但他既然已被真福\Person{亚历山大}免职，且他和其他人都与亚略的疯狂有染，加之他们所负的种种指控，神圣大公会议一致决定剥夺他们的主教职，我们认定他们不仅不是主教，而且不配与信友共融。\par

因为那些离间圣子、使圣言远离圣父的人，自己理当被排除于\OriginalTerm{天主教会}{Catholic Church}之外，与\OriginalTerm{基督徒之名}{Christian name}无分。因此，你们当视他们受\OriginalTerm{绝罚}{anathema}，因为他们\BibleQuote{错乱了真理的言语}（\BibleRef{格后 2:17}）。这是宗徒的训令（\BibleRef{迦 1:9}）：\BibleQuote{若有人给你们讲授别的福音，与你们所领受的不同，此人当受诅咒。}你们要嘱咐你们的子民，谁也不要与他们相通，因为光明与黑暗没有什么相通；\Person{基督}与\Person{贝里雅耳}也没有什么和谐（\BibleRef{格后 6:14}）。亲爱的诸位，你们要谨慎，既不要给他们写信，也不要接受他们的来信；却要切望，弟兄们和同工们，就如你们亲身在我们的公议会中（\BibleRef{格前 5:3}），以你们的署名赞同我们的决议，好使各地所有我们的同工都保持和谐。愿\OriginalTerm{天主眷顾}{Divine Providence}在圣化与喜乐中保护并保守你们，亲爱的弟兄们。\par

\Signature{我，\Person{霍修斯}主教，签署了此信，其余所有主教也同样签署。}\par

这就是\Place{萨尔底加}公议会致那些未能出席者的信函，而他们也相应地表明了他们的判断；以下是那些在公议会中署名的主教以及其他人的名录。\par

\Place{西班牙}的\Person{霍修斯}、\Place{罗马}的\Person{犹利}（由他的司铎\Person{阿基达穆}和\Person{斐洛色诺}代表）、\Place{萨尔底加}的\Person{普罗托格内斯}、\Person{戈登提乌斯}、\Person{马切多尼乌斯}、\Person{塞维鲁斯}、\Person{普雷特克塔图斯}、\Person{乌尔西基乌斯}、\Person{卢基卢斯}、\Person{欧根尼乌斯}、\Person{维塔利乌斯}、\Person{卡莱波迪乌斯}、\Person{弗洛伦提乌斯}、\Person{巴苏斯}、\Person{文森提乌斯}、\Person{斯特科里乌斯}、\Person{帕拉迪乌斯}、\Person{多米提安努斯}、\Person{卡尔比斯}、\Person{格隆提乌斯}、\Person{普罗塔西乌斯}、\Person{欧罗古斯}、\Person{波菲里乌斯}、\Person{迪奥斯库鲁斯}、\Person{佐西穆斯}、\Person{雅努阿里乌斯}、\Person{佐西穆斯}、\Person{亚历山大}、\Person{欧提基乌斯}、\Person{苏格拉底}、\Person{迪奥多鲁斯}、\Person{马尔提里乌斯}、\Person{欧特里乌斯}、\Person{欧卡尔普斯}、\Person{阿提诺多鲁斯}、\Person{伊雷内乌斯}、\Person{犹利安努斯}、\Person{阿里皮乌斯}、\Person{约纳斯}、\Person{埃提乌斯}、\Person{雷斯提图图斯}、\Person{马尔切利努斯}、\Person{阿普里安努斯}、\Person{维塔利乌斯}、\Person{瓦伦斯}、\Person{赫尔莫格内斯}、\Person{卡斯图斯}、\Person{多米提安努斯}、\Person{福尔图纳提乌斯}、\Person{马尔库斯}、\Person{安尼安努斯}、\Person{赫利奥多鲁斯}、\Person{穆塞乌斯}、\Person{阿斯特里乌斯}、\Person{帕雷戈里乌斯}、\Person{普鲁塔尔库斯}、\Person{希梅奈乌斯}、\Person{亚大纳修斯}、\Person{卢基乌斯}、\Person{阿曼提乌斯}、\Person{阿里乌斯}、\Person{阿斯克勒皮乌斯}、\Person{迪奥尼修斯}、\Person{马克西穆斯}、\Person{特里丰}、\Person{亚历山大}、\Person{安提哥努斯}、\Person{埃利安努斯}、\Person{佩特鲁斯}、\Person{辛弗鲁斯}、\Person{穆索尼乌斯}、\Person{欧提库斯}、\Person{菲洛罗吉乌斯}、\Person{斯普达西乌斯}、\Person{佐西穆斯}、\Person{帕特里基乌斯}、\Person{阿多利乌斯}、\Person{撒普里基乌斯}。\par

来自\Place{高卢}的有以下：\Person{马克西米安努斯}、\Person{维里西穆斯}、\Person{维克图鲁斯}、\Person{瓦伦提努斯}、\Person{德西德里乌斯}、\Person{欧罗吉乌斯}、\Person{萨尔巴提乌斯}、\Person{迪斯科利乌斯}、\Person{苏佩里奥尔}、\Person{梅尔库里乌斯}、\Person{德克洛佩图斯}、\Person{欧塞比乌斯}、\Person{塞维里努斯}、\Person{萨提鲁斯}、\Person{马尔提努斯}、\Person{保卢斯}、\Person{奥普塔提安努斯}、\Person{尼卡西乌斯}、\Person{维克托尔}、\Person{森普罗尼乌斯}、\Person{瓦勒里努斯}、\Person{帕卡图斯}、\Person{耶塞斯}、\Person{阿里斯顿}、\Person{辛普利基乌斯}、\Person{梅提安努斯}、\Person{阿曼图斯}、\Person{阿米利安努斯}、\Person{犹斯提尼安努斯}、\Person{维克托里努斯}、\Person{萨托尔尼卢斯}、\Person{阿布恩丹提乌斯}、\Person{多纳提安努斯}、\Person{马克西穆斯}。\par

来自\Place{非洲}的有：\Person{内苏斯}、\Person{格拉图斯}、\Person{梅加西乌斯}、\Person{科尔代乌斯}、\Person{罗加提安努斯}、\Person{孔索尔提乌斯}、\Person{鲁菲努斯}、\Person{曼尼努斯}、\Person{切西利安努斯}、\Person{赫伦尼安努斯}、\Person{马里安努斯}、\Person{瓦勒里乌斯}、\Person{迪纳米乌斯}、\Person{米佐尼乌斯}、\Person{犹斯图斯}、\Person{策莱斯蒂努斯}、\Person{基普里安努斯}、\Person{维克托尔}、\Person{霍诺拉图斯}、\Person{马里努斯}、\Person{潘塔加图斯}、\Person{斐利克斯}、\Person{鲍迪乌斯}、\Person{利贝尔}、\Person{卡皮托}、\Person{米内尔瓦利斯}、\Person{科斯穆斯}、\Person{维克托尔}、\Person{赫斯佩里奥}、\Person{斐利克斯}、\Person{塞维里安努斯}、\Person{奥普坦提乌斯}、\Person{赫斯佩鲁斯}、\Person{菲登提乌斯}、\Person{撒鲁斯提乌斯}、\Person{帕斯卡西乌斯}。\par

来自\Place{埃及}的有：\Person{利布尔尼乌斯}、\Person{阿曼提乌斯}、\Person{斐利克斯}、\Person{伊斯希拉蒙}、\Person{罗穆卢斯}、\Person{提贝里努斯}、\Person{孔索尔提乌斯}、\Person{赫拉克利德斯}、\Person{福尔图纳提乌斯}、\Person{迪奥斯库鲁斯}、\Person{福尔图纳提安努斯}、\Person{巴斯塔蒙}、\Person{达提卢斯}、\Person{安德烈亚斯}、\Person{塞雷努斯}、\Person{阿里乌斯}、\Person{泰奥多鲁斯}、\Person{埃瓦戈拉斯}、\Person{赫利亚斯}、\Person{提摩泰乌斯}、\Person{奥里翁}、\Person{安德洛尼库斯}、\Person{帕弗努提乌斯}、\Person{赫尔米亚斯}、\Person{阿拉比翁}、\Person{普塞诺西里斯}、\Person{阿波罗尼乌斯}、\Person{穆伊斯}、\Person{撒拉帕姆蓬}、\Person{菲洛}、\Person{菲利普斯}、\Person{阿波罗尼乌斯}、\Person{帕弗努提乌斯}、\Person{保卢斯}、\Person{迪奥斯库鲁斯}、\Person{尼拉蒙}、\Person{塞雷努斯}、\Person{阿奎拉}、\Person{阿奥塔斯}、\Person{哈尔波克拉提翁}、\Person{伊撒克}、\Person{泰奥多鲁斯}、\Person{阿波洛斯}、\Person{阿莫尼安努斯}、\Person{尼卢斯}、\Person{赫拉克利乌斯}、\Person{阿里翁}、\Person{阿塔斯}、\Person{阿尔塞尼乌斯}、\Person{阿加撒蒙}、\Person{泰翁}、\Person{阿波罗尼乌斯}、\Person{赫利亚斯}、\Person{帕尼努提乌斯}、\Person{安德拉加提乌斯}、\Person{内梅西翁}、\Person{撒拉皮翁}、\Person{阿莫尼乌斯}、\Person{阿莫尼乌斯}、\Person{克塞农}、\Person{格隆提乌斯}、\Person{昆图斯}、\Person{莱奥尼德斯}、\Person{森普罗尼安努斯}、\Person{菲洛}、\Person{赫拉克利德斯}、\Person{希埃拉基斯}、\Person{鲁弗斯}、\Person{帕索菲乌斯}、\Person{马切多尼乌斯}、\Person{阿波罗多鲁斯}、\Person{弗拉维安努斯}、\Person{普萨埃斯}、\Person{叙鲁斯}、\Person{阿普弗斯}、\Person{撒拉皮翁}、\Person{埃撒伊亚斯}、\Person{帕弗努提乌斯}、\Person{提摩泰乌斯}、\Person{埃卢里翁}、\Person{盖乌斯}、\Person{穆塞乌斯}、\Person{皮斯图斯}、\Person{赫拉克拉蒙}、\Person{赫隆}、\Person{赫利亚斯}、\Person{阿纳加姆弗斯}、\Person{阿波罗尼乌斯}、\Person{盖乌斯}、\Person{菲洛塔斯}、\Person{保卢斯}、\Person{提托埃斯}、\Person{欧代蒙}、\Person{犹利乌斯}。\par

意大利道上的主教有：\Person{普罗巴提乌斯}、\Person{维亚托尔}、\Person{法昆迪努斯}、\Person{若瑟}、\Person{努梅迪乌斯}、\Person{斯佩兰提乌斯}、\Person{塞维鲁斯}、\Person{赫拉克利安努斯}、\Person{福斯提努斯}、\Person{安东尼努斯}、\Person{赫拉克利乌斯}、\Person{维塔利乌斯}、\Person{斐利克斯}、\Person{克里斯皮努斯}、\Person{保利安努斯}。\par

来自\Place{塞浦路斯}的有：\Person{奥克西比乌斯}、\Person{福提乌斯}、\Person{杰拉西乌斯}、\Person{阿弗罗迪西乌斯}、\Person{伊雷尼库斯}、\Person{努内基乌斯}、\Person{亚大纳修斯}、\Person{马切多尼乌斯}、\Person{特里菲利乌斯}、\Person{斯皮里东}、\Person{诺尔巴努斯}、\Person{索西克拉特斯}。\par

来自\Place{巴勒斯坦}的有：\Person{马克西穆斯}、\Person{埃提乌斯}、\Person{阿里乌斯}、\Person{泰奥多西乌斯}、\Person{杰尔马努斯}、\Person{西尔瓦努斯}、\Person{保卢斯}、\Person{克劳迪乌斯}、\Person{帕特里基乌斯}、\Person{埃尔皮迪乌斯}、\Person{杰尔马努斯}、\Person{欧塞比乌斯}、\Person{泽诺比乌斯}、\Person{保卢斯}、\Person{佩特鲁斯}。\par

以上是那些签署公会议决议者的名字；但除此之外，还有很多人，来自\Place{亚细亚}、\Place{弗里吉亚}和\Place{伊扫利亚}，他们在我之前，即在这公议会举行之前，便写信支持我，他们的名字，大约有六十三位，可在他们自己的信函中找到。总计共有三百四十四人。\par

% structure: p0117 source=h2 target=subsection reason=relative-heading-rank; base=h2

\subsection{第四章 \Emph{萨尔底加公议会议决之后所发生的帝国与教会行动}}

当最虔诚的皇帝\Person{君士坦提乌斯}听闻这些事后，他派人来找我，此前他已私下致函给他那荣福记忆的兄弟\Person{君士坦斯}，也三次分别致函给我，内容如下。\par

\Signature{凯旋者奥古斯都\Person{君士坦提乌斯}致\Person{亚大纳修斯}。}我们的仁慈悲悯不容你再被大海的狂浪颠簸；因为我们不倦的虔敬从未遗忘你，尽管你失去故土，被剥夺财产，在荒蛮的旷野中漂泊。诚然，我久已延宕，未曾以信函表达我对你的心意，主要原因为我期望你会主动前来见我，寻得你痛苦的缓解；然而，或许由于恐惧，你未能实现你的意愿，因此我们致函你的坚毅，满载我们的慷慨，望你急速、无畏地前来我所面前，从而获享你所愿，使你既经验了我们的仁慈，便可重返你自己的地方。为此，我已为你恳求我的主上、兄弟凯旋者奥古斯都\Person{君士坦斯}，请他允许你前来，以便经我们二人同意，将你复归故土，以此作为我们眷顾的保证。\par

\Emph{第二封信}\par

尽管我们在前一封信中已向你们清楚表明，你们可以毫不犹豫地来到我们的宫廷，因为我们非常希望送你们回家；然而，我们又再给你们写了这封信，为劝勉你们的坚毅，不要有任何疑虑或恐惧，要搭乘公共交通工具，尽快来到我们这里，好使你们的愿望得以实现。\par

\Emph{第三封信}\par

当我们住在\Place{埃德萨}时，你们的司铎们也在那里，我们乐意让其中一人被派到你们那里，你们应赶快来到我们的宫廷，以便能见我们的面，然后立刻前往\Place{亚历山大里亚}。但自从你们收到我们的信以来，过了很久，你们还没有来，因此我们赶紧再次提醒你们，使你们现在仍努力尽快来到我们面前，这样你们就可以重返故乡，并使你们的祈求得以实现。为使你们更充分了解情况，我们已派遣执事\Person{阿基塔斯}，你们将从他得知我们的心意，从而现在就能确保你们所祈求的事。\par

这就是皇帝信函的内容；收到这些信后，我前往\Place{罗马}，向教会和主教告别：因为写这些信时我在\Place{阿奎莱亚}。教会充满了喜乐，主教\Person{尤利乌斯}因我返回而与我同乐，并写信给教会；我们一路经过时，各地的主教都平安地送我们上路。尤利乌斯的信如下。\par

52. \Person{尤利乌斯}致居住在\Place{亚历山大里亚}的司铎、执事和民众。\par

我祝贺你们，亲爱的弟兄们，因为你们如今亲眼看见你们信德的果实；任何人都能看到，关于我的弟兄和同为主教的\Person{亚大纳削}，确实如此，天主因着他生活的无罪，并由于你们的祈祷，再次把他归还给你们。因此不难看出，你们一直向天主献上纯洁和充满爱心的祈祷。你们记念天上的应许，以及那通往这些应许的交谈，这些你们是从我前述弟兄的教导中学来的；你们凭借在你们内的正确信仰，确实知道并明白，那位你们一向以最虔诚的心怀视为临在的，绝不会永远与你们分离。因此，我无需在写信给你们时多费言辞；因为你们的信仰已经预先包括了我要对你们说的一切，并借着天主的恩宠，实现了你们众人共同的祈祷。所以，我再次祝贺你们，因为你们在信德中保全了你们的灵魂，使它们不受征服；我也同样祝贺我的弟兄\Person{亚大纳削}，因为他虽经历许多磨难，却从未忘记你们对他的爱和热切的愿望。虽然他一时似乎与你们在形体上分离，但他一直在精神上与你们同在，犹如始终临在一样。\par

53. 因此，他现在回到你们这里，比他离开时更加光辉。火能试验和炼净贵重的材料，如金和银：但人怎能描述这样一个人的价值呢？他经过如此多患难的危险，胜利归来，现在被宣告无罪，不仅凭我们的声音，更是凭整个会议的声音？因此，亲爱的弟兄们，请以一切虔敬的尊崇和喜乐，接纳你们的主教\Person{亚大纳削}，以及那些在许多劳苦中与他合作的人。并且要喜乐，因为你们现在获得了你们祈祷的实现，你们曾在有益的信中，赐给你们的牧者食粮和饮品，他可谓是渴慕你们的虔诚。因为当他旅居异乡时，你们是他的安慰；在他受迫害时，你们用最忠信的心灵和精神使他振作。我现在欣喜地想象和揣摩你们每个人在他返回时的喜乐，那些聚集的人群的虔诚致意，以及那些跑着迎接他的人的荣耀庆祝。当我的弟兄回来，你们先前的苦难结束，他那备受珍视和期盼的归来使你们全都充满圆满的喜乐，这对你们将是何等的一天啊！这同样的喜乐我们也在极大程度上感受到，因为天主赐予我们能够结识如此卓越的人。因此，我宜以祈祷结束此信。愿全能的天主，以及他的圣子我们的主和救主耶稣基督，赐给你们持续的恩宠，因你们为你们的主教所作的高尚宣认中表现的可嘉信德而赏报你们，使他能将那些更好的事物，即\BibleQuote{眼所未见，耳所未闻，人心所未想到的，就是天主为爱他的人所预备的}\BibleRef{格前 2:9}，在此生及来世，赐予你们和与你们同在的人；以上所求是靠我们的主耶稣基督，愿光荣藉着他归于全能的天主，至于无穷世之世。阿们。亲爱的弟兄们，我祈求你们在主内健康强壮。\par

54. 当我带着这些信来到皇帝面前时，他亲切地接待了我，并送我回我的故乡和教会，同时向主教、司铎和民众发表了如下讲话。\par

君士坦提乌斯，常胜者，伟大者，奥古斯都，致天主教会的主教和司铎。\par

\Quote{\Person{亚大纳削}至可敬者，并未被天主的\OriginalTerm{恩宠}{grace}所离弃；他虽曾短暂经受人性所难免的考验，却已从遍察万有的\OriginalTerm{天主眷顾}{Providence}那里获得了适以应其祈祷的恩允，且依照至高者的旨意及我们的裁定，同时被归还给他的国家与教会——他曾荣膺天授而执掌该教会。因此，依此而行，理应由我们的仁慈作出规定：凡此前针对那些与他共融者所颁布的一切法令，均应予以遗忘，同时从此消除对他们的所有嫌疑，并且，与他联合的圣职人员先前所享有的各种豁免权，应照例赐予他们确认。此外，在我们对他施予的其他恩惠之外，我们还乐意补充如下：凡列入圣职人员名录者，均应知晓，凡归附他的人，无论是主教或其他圣职人员，均已获得安全保证。并且，与他的结合，对任何人而言，都将是其意图正直的充分担保。因为，凡基于更佳判断与抉择、选择与他共融的人，我们命令，效法那先行的天主眷顾，让所有这样的人都能享有那按照至高者的旨意，如今由我们赐予他们的恩宠的益处。愿天主保佑你们。}\par

\Emph{第二封信函。}\par

\Quote{\Person{君士坦提乌斯}、\Person{维克托}、\Person{马克西穆斯}，\OriginalTerm{奥古斯都}{Augustus}，致\Place{亚历山大里亚}的天主教会的民众。}\par

\Quote{55. 我们关心你们在各方面福祉，并了解你们长时间缺少主教监督，遂决定将你们的主教\Person{亚大纳削}差回你们那里；他是一个以内在正直和个人品性善良而为人所共知的人物。请以适宜的方式接待他，一如你们惯常接待每个人那样，并以他的代祷作为你们向天主祈祷时的援助，努力依照教会的秩序，保持始终如一的和谐与和平，这既对你们适宜，也对我们最为有利。因为，在你们中间兴起任何纷争或派系，违背我们时代的繁荣，是不合宜的。我们愿意将这种罪行从你们中间完全消除，并劝勉你们坚持惯常的祈祷，并如前所言，以他为你们在天主面前的代祷者和援助者。这样，当你们的这种决心，亲爱的各位，影响了所有人的祈祷时，就连那些至今仍沉迷于偶像虚假敬拜的外教人，也会切愿前来认识我们神圣的宗教。因此，我们再度劝勉你们在这方面持之以恒，并欣然接纳你们的主教——他奉至高者的谕令和我们的决定被差回你们这里——且要全心全力地、真诚地迎接他。因为这既适合你们，也合乎我们的仁慈。为消除那些心怀恶意者制造骚乱和叛乱的种种机会，我们已通过信函命令你们那里的长官，对所有被发现有党派行为的人，依法惩处。所以，你们要考虑到这些：我们遵循至高者旨意的决定、我们对你们及你们彼此和睦的关切、以及等待滋事者的惩罚；你们要遵守那合乎我们神圣宗教秩序所适宜和恰当的事，并以十足的敬仰和尊崇迎接上述主教，还要留心与他一起为你们自己、以及为你们整个生活的康泰，向天主、万有之父献上祈祷。}\par

56. 写罢这些信函后，他还下令，把他先前因\Person{优西比乌}及其同党的诬告而发出的、针对我的那些谕令，从督军和埃及总督的谕令档中撤销并删除；并派遣什长\Person{优西比乌}去将这些谕令从谕令册中抽回。他为此事发出的信函如下：\par

\Quote{\Person{君士坦提乌斯}、\Person{维克托}，\OriginalTerm{奥古斯都}{Augustus}，致\Person{乃斯托留}。\Emph{并以相同措辞致\Place{奥古斯塔姆尼卡}、\Place{特拜}和\Place{利比亚}的总督。}}\par

\Quote{凡是先前颁布的、旨在损害和侮辱那些与\Person{亚大纳削}主教共融之人的任何谕令，我们均希望现在予以删除。因为我们愿意，他的圣职人员从前所享有的一切豁免权，他们应重新享有。我们希望这一谕令得到遵守：即当\Person{亚大纳削}主教被恢复至他的教会时，与他共融的人可享有他们素来享有、且其余圣职人员也享有的各项豁免权；这样他们就可以满意地处于与他人同等的地位。}\par

57. 我如此踏上旅程，途经\Place{叙利亚}的时候，遇见了\Place{巴勒斯坦}的主教们；他们在\Place{耶路撒冷}召开了一次公会议，热诚接待了我，并且自己也以平安送我上路，还向教会和主教们发出下面的信函。\par

\Quote{在\Place{耶路撒冷}聚集的圣公会议，致在\Place{埃及}和\Place{利比亚}的同工们，以及\Place{亚历山大里亚}的长老、执事和民众，我们深爱的、极为思念的弟兄们，谨从主内问安。}\par

我们无法向万有的天主献上相称的感谢，最亲爱的诸位，为祂在历世历代所行的奇妙之事，尤其在此刻为你们的教会所做的，就是将你们的牧者与主、我们的同工\Person{亚大纳削}归还给你们。因为谁曾期望自己的眼睛看见你们如今实际所得的？确实，你们祈祷已被万有的天主垂听，祂眷顾自己的教会，注视了你们的眼泪与叹息，因此俯允了你们的恳求。因为你们曾像没有牧人的羊群，困苦流离而昏倦。\BibleRef{玛 9:36} 因此，那位真正的牧者，眷顾自己羊群的，从天上探望了你们，将你们所渴望的归还了你们。看，我们也将一切为教会的和平而预备，并被与你们同样的热情所催促，已在他之前向你们致意；并通过他与你们相通，我们向你们送上这些问候，及我们的感恩奉献，好让你们知道我们也同样连结在那将你们与他相连的爱德纽带中。你们还有责任为我们最蒙天主喜爱的皇帝们的虔敬祈祷，他们知道了你们对他的深切渴慕和他是无辜的，就决定以一切尊荣将他归还给你们。因此，举起手来迎接他，并要谨慎，为他向那位赐予你们这些福乐的天主献上相称的感谢；好使你们能偕同天主，常常喜乐，并光荣我们的主，在于我们的主耶稣基督内，借着祂，光荣归于圣父，至于永世。阿们。\par

我在此列出这封信的签署人名单，虽然我在前面已经提到过他们。他们是：\Person{马克西穆}、\Person{埃提乌}、\Person{亚略}、\Person{德奥多若}、\Person{革尔曼}、\Person{西尔望}、\Person{保禄}、\Person{帕特里克}、\Person{埃尔皮狄}、\Person{革尔曼}、\Person{优西比乌}、\Person{泽诺比乌}、\Person{保禄}、\Person{马克里努}、\Person{伯多禄}、\Person{克劳狄}。\par

58. 当\Person{乌尔撒齐}与\Person{瓦伦斯}看到这一切时，他们立刻谴责自己所为，去到\Place{罗马}，承认自己的罪行，表示悔改，并请求宽恕，向古罗马城主教\Person{犹利乌}和我们寄出了以下信件。这些信件的抄本由\Place{特雷维里}的主教\Person{保理诺}寄给了我。\par

\Emph{一封拉丁文信件的译文，致\Person{犹利乌}，论\Person{乌尔撒齐}与\Person{瓦伦斯}的悔改。}\par

\Person{乌尔撒齐}与\Person{瓦伦斯}致至可颂扬的主、教宗\Person{犹利乌}。\par

众所周知，我们先前曾在信中提出许多严重的指控，反对主教\Person{亚大纳削}；而当我们受到你的仁善之信的纠正时，我们未能对自己所作的陈述给出交代；我们现在在你的仁善面前，并在所有我们兄弟司铎的在场下，承认所有先前传到你耳中的关于前述\Person{亚大纳削}的事件的报告，都是虚假和捏造的，与他的品行完全不符。因此，我们热切渴望与前述\Person{亚大纳削}共融，尤其因你的虔敬，以你特有的慷慨，已俯允宽恕我们的过失。但我们也声明，倘若任何时候东方的主教们，甚或\Person{亚大纳削}本人，不宽仁地想因此事将我们交付审判，我们决不会违背你的判决。至于异端者\Person{亚略}及其党羽，他们说圣子曾不存在，圣子是由无中造成的，并否认基督在万世之前是天主及天主之子，我们现在并直到永远绝罚他们，正如我们先前在\Place{米兰}的声明中所陈述的。我们亲手写了这些，并再次宣认，我们一如前所述，已永远弃绝了亚略异端及其创始者。\par

\Signature{我\Person{乌尔撒齐}亲笔签署此我的认罪书；同样，我\Person{瓦伦斯}也签署。}\par

主教\Person{乌尔撒齐}与\Person{瓦伦斯}，致他们的主内弟兄，主教\Person{亚大纳削}。\par

趁着有我们的兄弟和同僚司铎\Person{穆塞乌斯}，前往你的爱德前来的机会，我们从\Place{阿奎莱亚}，通过他亲切地问候你，亲爱的弟兄，并祈求你——我们相信你身体健康——阅读我们的信。你若能书面回覆我们，也将给予我们信心。因为你知道我们与你和平相处，并与教会共融，这封信开头的问候即是证明。愿天主眷顾保存你，我主，我们亲爱的弟兄！\par

这便是他们的信件，以及主教们为我所作的定断和判决。但为了证明他们这样做并非为了讨好，或是出于任何方面的胁迫，我愿意，若蒙你许可，将整件事从头讲述，好使你们明白，主教们写信时的作为，是出于正直和公义的意向，且\Person{乌尔撒齐}与\Person{瓦伦斯}，尽管行动迟缓，终究坦白了真相。\par

\SourceRecord{Translated by M. Atkinson and Archibald Robertson.；Nicene and Post-Nicene Fathers, Second Series；Vol. 4.；Edited by Philip Schaff and Henry Wace.；Buffalo, NY: Christian Literature Publishing Co.,；1892.；<http://www.newadvent.org/fathers/28081.htm>.}

% structure: p0001 source=h1 target=section reason=page-title

\section{\WorkTitle{驳亚略派}（第二部分）}

% structure: p0002 source=h2 target=subsection reason=relative-heading-rank; base=h2

\subsection{第五章 \Emph{与\OriginalTerm{默肋提派}{Meletians}控告\Person{圣亚大纳修}有关的文件}}

59. 教难之前，\Person{伯多禄}曾是我们的主教，并在教难中殉道。当在\Place{埃及}拥有主教头衔的\Person{默肋提}被证实犯有许多罪行，其中包括向偶像献祭时，\Person{伯多禄}在一次主教全体会议中将其罢免。于是\Person{默肋提}没有向另一会议申诉，也没有试图在后来者面前为己辩护，反而制造了分裂，以致拥护他的人至今仍被称为默肋提派而非基督徒。他立即开始辱骂主教们，并进行诬告，首先针对\Person{伯多禄}本人，然后针对他的继任著\Person{阿基拉斯}，之后又针对\Person{亚历山大}。他效法\Person{阿贝沙隆}的榜样，狡猾行事，以期既然因被罢免而遭受耻辱，便可借诽谤误导无知之人。正当\Person{默肋提}如此行事之时，\OriginalTerm{亚略异端}{Arian heresy}也兴起了。但在\OriginalTerm{尼西亚大公会议}{Council of Nicæa}上，一方面该异端被绝罚，\OriginalTerm{亚略派}{Arians}被驱逐，另一方面出于某些理由（此处无需提及）默肋提派被接纳了。然而，还未满五个月，真福\Person{亚历山大}去世，那些本该保持安静、为能蒙接纳而心存感激的默肋提派，却像狗无法忘记它们的呕吐物一样，又开始扰乱各教会。\par

\Person{欧瑟比乌斯}获悉这些后，便以重金承诺收买\OriginalTerm{默肋提派}{Meletians}，成为他们的秘密朋友，并与他们商定，在他需要的时候可借其助力。起初，他先写信给我，催促我接纳\Person{亚略}及其同伙共融，口头上威胁我，而在信中却[仅仅]提出请求。当我拒绝，并宣布那些违背真理捏造异端、已被\OriginalTerm{大公会议}{Ecumenical Council}绝罚的人，不应被接纳共融时，他便促使真福记忆的\Person{君士坦丁}皇帝也写信给我，威胁我说，若我不接纳\Person{亚略}及其同伙，我必将遭受我之前已经历、现仍在忍受的那些苦难。以下是他的信的一部分。\Person{辛克肋提}和\Person{嘎乌登提}，宫廷官员，是送信的人。\par

\Emph{君士坦丁皇帝信函的一部分}\par

因此，你既知道我的旨意，就当允许一切愿意进入教会的人自由进入。因为我若得知你阻止或排斥任何要求被接纳进入教会共融的人，我将立即派人，奉我的命令罢免你，并把你从你的位置上赶走。\par

60. 于是，当我写信并力求说服皇帝，那反基督的异端与天主教会绝无共融时，\Person{欧瑟比乌斯}立即利用他与\OriginalTerm{默肋提派}{Meletians}事先约定的机会，写信怂恿他们编造借口，以便如他们曾对待\Person{伯多禄}、\Person{阿基拉斯}和\Person{亚历山大}那样，也设计并散布谣言中伤我们。因此，他们长期寻找却一无所获后，终于在\Person{欧瑟比乌斯}及其同伙的建议下达成一致，通过\Person{依息盎}、\Person{欧德蒙}和\Person{加里尼各}，首次捏造了关于\OriginalTerm{麻布圣衣}{linen vestments}的控告，声称我向埃及人强加了一条法律，并要求他们首先遵守。但当我的几位司铎在场，皇帝审理此事时，他们被判有罪（这两位司铎是\Person{阿丕斯}和\Person{马加略}），皇帝便写信谴责\Person{依息盎}，并命我前去见他。他的信如下。\par

\Person{欧瑟比乌斯}得知此事，怂恿他们等待；等我抵达后，他们接下来又指控\Person{马加略}打碎了圣爵，并对我提出了一项极为恶毒的控告，即声称我是皇帝之敌，曾给一个名叫\Person{斐路门}的人送去一袋黄金。于是皇帝在\Place{普撒玛提亚}就这一指控也听取了我们的申诉，结果他们照例被定罪，并被逐出御前；当我返回时，他给民众写了以下这封信。\par

\Person{君士坦丁}·\Person{马克西穆}·\Person{奥古斯都}致\Place{亚历山大里亚}天主教会全体教友。\par

61. 亲爱的主内弟兄们，我向你们问安，并呼求天主，祂是我心意的最主要见证，并呼求唯一圣子，我们律法的创立者，祂掌管万民的生命，憎恨纷争。我还能对你们说什么呢？说我身体健康吗？不，如果你们彼此相爱，消除敌意——正是这些敌意，因好争之人所起的风暴，使我们离开了兄弟友爱的港湾——我本当享有更好的健康和力量。可叹！这是何等乖谬！你们中间所激起的嫉妒骚动，天天产生怎样的恶果！正因如此，恶言蜚语已笼罩着天主的子民。正义的信仰又往何处去了？因为我们不仅因重重错谬，也因忘恩负义之人的过失而深深陷入黑暗迷雾，以至我们竟容忍那些喜好愚昧的人，虽明知如此，却不理会那些摒弃良善与真理的人。这是何等奇怪的矛盾！我们不定敌人的罪，却效法他们给我们立下的抢劫恶例，于是最有害的错谬，既无人抵挡，便轻易——若可这么说——为自己开辟了道路。难道我们竟毫无悟性，不顾我们共同人性的信誉，如此忽视律法的训令吗？\par

但有人或许会说，爱是生而有之的本性。然而我要问：我们既有天生的优势，又以天主的法律为指引，为何竟容忍敌人挑起的纷争扰乱？他们仿佛被火把点燃，怒火中烧。为什么我们有眼睛却看不见，有理智却不明白，尽管我们被法律的智慧所环绕？怎样的麻木攫住了我们的生命，使我们如此疏忽自己，即使天主警告我们？这岂非难以容忍的灾祸？我们岂不该视这样的人为敌人，而非天主的家人和子民？因为他们被遗弃，就疯狂攻击我们：他们把重罪加在我们身上，好像敌人一样来攻击我们。\par

62. 我愿你们仔细想想，他们这样做是多么疯狂。这些愚昧的人口无遮拦，恶意脱口而出。他们怀着铅块般的怒气，以致彼此攻击，并连累我们以加重他们自己的惩罚。好导师被当作敌人，而那身披嫉妒之恶的人，却违背一切公义，利用民众温和的性情谋取私利；他破坏、消耗，炫耀自己，用虚假的赞美自荐；他颠覆真理，败坏\OriginalTerm{信德}{faith}，直到为他的良心找到藏身之处。这样，他们的邪恶使他们不幸，同时他们无耻地把自己推举到尊位，尽管完全不配。啊！这是何等的祸害！他们说：\Quote{这个人太老了；那个人不过是个孩子；这职位该是我的；既然他被剥夺，便应当归我。我要争取一切人支持我，然后设法用我的权力毁灭他。} 他们疯癫的宣告昭然于世；就是那些结党、聚众、并在可憎的阴谋中指挥划桨者的集会。唉！若我可这样说，我们的行为是何等荒谬！他们是在天主的教会中展示自己的愚昧吗？他们还不羞愧吗？他们还不自责吗？他们的良心不被触动吗？以至于现在他们至少表现出对诡诈和争斗有所自觉？他们的力量不过是嫉妒暴力，得到与之相随的邪恶影响的支撑。但那些可怜人对你们的主教无能为力。相信我，弟兄们，他们的努力只有一种效果，就是消耗了我们的岁月，却使自己今生不再有忏悔的余地。因此，我恳求你们，自助自救；亲切接受我们的爱，全力赶走那些想从我们中间抹去同心合意之恩宠的人；仰望天主，彼此相爱。我欣然接纳你们的主教\Person{亚大纳削}，并以我认为他属于天主的人的方式与他交谈。这些事情应由你们判断，而非由我裁决。我以为最可敬的\Person{亚大纳削}本人应将我的问候传达给你们，我知道他慈爱地关怀你们，这种关怀与我自己宣认的和平\OriginalTerm{信德}{faith}相称，他不断致力于宣讲救恩知识的美善事业，并能给你们适当的劝勉。愿天主保佑你们，亲爱的弟兄们。\par

这就是\Person{君士坦丁}的信。\par

63. 这些事件之后，默肋提派安静了一小段时间，但随后又显露出敌意，策划了以下阴谋，为要讨好那些雇用了他们的人。\Place{玛瑞奥提斯}是\Place{亚历山大里亚}的乡村地区，\Person{默肋提}未能在此制造分裂。当时各教会仍在指定的界域内，所有\OriginalTerm{司铎}{presbyter}在其中都有会众，而民众也和平生活，有一个名叫\Person{依斯希辣}的人，并非圣职人员，品性卑劣，企图误导自己村庄的人，宣称自己是圣职人员。得知此事后，当地的本堂\OriginalTerm{司铎}{presbyter}在我巡视各教会时通知了我，我便派\Person{玛加略}\OriginalTerm{司铎}{presbyter}随他前往传唤\Person{依斯希辣}。他们发现他卧病在床，躺在一间小室里，于是嘱咐他的父亲告诫儿子，不要再有人们举报的那种行为。但他病愈后，被朋友和父亲阻止继续同样的行径，便投奔了默肋提派；他们与\Person{欧瑟比乌}及其同党勾结，最后编造了如下诽谤：说\Person{玛加略}打破了一只圣爵，又说一位名叫\Person{亚尔塞尼}的主教被我杀害。他们将\Person{亚尔塞尼}藏匿起来，好使他似乎已从人间消失；他们还到处展示一只手，谎称他被碎尸。至于那不认识的\Person{依斯希辣}，他们开始散布流言，称他是\OriginalTerm{司铎}{presbyter}，以便他关于圣爵的话能误导民众。然而，\Person{依斯希辣}遭到友人的责备，哭着来见我，声称\Person{玛加略}根本没有做过他们所说的那种事，并说自己是被默肋提派收买来捏造这诽谤的。于是他写了如下信函。\par

致真福\OriginalTerm{主教}{pope}\Person{亚大纳削}，\Person{依斯希辣}在主内请安。\par

64. 主教大人，当来到您面前，渴望被接纳进入教会时，您因我先前所说的话责备了我，仿佛我是自愿做出那等事的。因此，我向您呈上这份书面的自诉，好让您明白，是\Person{依撒格}、\Person{赫辣克里德}、\Place{勒托颇里斯}的\Person{依撒格}以及他们的同党对我施加暴力、殴打。我声明，并在此事上呼求天主为证，就他们所陈述的事，我知您没有任何罪责。因为从未发生过打破圣爵或倾覆\OriginalTerm{祭台}{Holy Table}的事，是他们用暴力手段迫使我如此供述。我作此辩护，并书面呈交给您，渴望并请求被接纳为您信众中的一员。我祈祷您在主内安康。\par

我将此亲笔函呈递给您，\Person{亚大纳削}主教，在以下诸位面前：司铎：\Place{狄切拉}的\Person{阿摩纳斯}、\Place{法斯科}的\Person{赫拉克利乌斯}、\Place{切内布里}的\Person{博孔}、\Place{米西纳}的\Person{阿奇拉斯}、\Place{塔福西里斯}的\Person{迪迪莫斯}，以及来自\Place{波莫修斯}的\Person{尤斯特斯}；执事：\Place{亚历山大里亚}的\Person{保禄}、\Person{伯多禄}、\Person{奥林比乌斯}，和\Place{马雷奥提斯}的\Person{阿摩尼乌斯}、\Person{皮斯图斯}、\Person{德默特留斯}、\Person{盖乌斯}。\par

65. 尽管\Person{伊斯奇拉斯}有此声明，他们仍到处散布同样的指控，并向\Person{君士坦丁}皇帝报告。皇帝先前在\Place{普萨玛提亚}时已听说过那杯爵的事，那时我在场，并已识破敌人的虚假。但现在他写信给\Place{安提约基}的监察官\Person{达玛提乌斯}，要求他就那起谋杀案进行司法调查。于是监察官通知我准备为指控辩护。收到信后，虽然我起初不予理睬，因我知道他们所说的无一属实，但见皇帝已有所动，我便写信给在\Place{埃及}的同工，并派了一名执事，希望能打听到\Person{阿尔塞尼乌斯}的消息，因我有五六年未见到此人了。长话短说，\Person{阿尔塞尼乌斯}先是在\Place{埃及}被找到，随后我的朋友们又在\Place{提洛}发现他再次藏匿。最值得注意的是，即便被发现后他仍不肯承认自己是\Person{阿尔塞尼乌斯}，直到在法庭上当着时任\Place{提洛}主教的\Person{保禄}的面被定罪，最后出于羞惭，才不得不承认。\par

他如此行，是为履行他与\Person{欧瑟比乌斯}及其同党的约定，以免若他被发现，他们所玩的把戏终会被戳穿；事实也确是如此。因当我写信给皇帝，告知\Person{阿尔塞尼乌斯}已被找到，并提醒他先前在\Place{普萨玛提亚}所听关于司铎\Person{玛加略}之事后，他便停止了监察官法庭的程序，并写信谴责那些针对我的控告纯属诬蔑，且命令正要前往东方指控我的\Person{欧瑟比乌斯}及其同党返回。为表明他们曾控告我谋杀\Person{阿尔塞尼乌斯}（无须引出许多人的信件），这里只引述\Place{得撒洛尼}主教\Person{亚历山大}的一封信便足矣，从中可推知其余信函的基调。他得知了那也被称为\Person{若望}的\Person{阿尔恰夫}就谋杀一事对我散布的谣言，并听说\Person{阿尔塞尼乌斯}尚在人世，便写了如下这封信。\par

\Emph{亚历山大的信函}\par

致亲爱的儿子、志同道合的同工、主\Person{亚大纳削}，主教\Person{亚历山大}在主内致以问候。\par

66. 我为最卓越的\Person{萨拉匹翁}感到欣喜，他如此勤奋地力求以圣善的习惯装饰自己，并以此使其父的纪念德愈显高尚。因为，正如圣经某处所说：\BibleQuote{他的父亲虽然死了，却好像没有死一样}（\BibleRef{德 30:4}），因他留下了生命的纪念。至于我对永志不忘的\Person{索佐}的感情，我主，您自己并非不知，因您深知他神圣的纪念，以及这年轻人的良善。我只收到您尊敬的一封信，是由这年轻人带来的。我主，我提此事是让您知道。我们亲爱的主内兄弟及执事\Person{玛加略}，从\Place{君士坦丁堡}写信来，使我深感欣慰，他告知那诬告者\Person{阿尔恰夫}已蒙羞辱，因其当众宣称一个活人已被谋杀。至于他将和其所有同党从公义审判者领受那与其罪行相称的惩罚，无谬的圣经向我们作出了保证。\newline{}愿万有之主保全您多年，我主，您至为仁慈的。\par

67. 与\Person{阿尔塞尼乌斯}同住的人作证说，他之所以被藏匿，正是为了让他们可假称他已死；因为我们在搜寻他时找到了那个[做了此事的人]，他于是给\Person{若望}写了下面这封信，而\Person{若望}正是这次诬告的主谋。\par

致亲爱的兄弟\Person{若望}，\Place{安特奥波利斯}地区\Place{普特门基尔基斯}隐修院的司铎\Person{品内斯}，致以问候。\par

我希望您知道，\Person{亚大纳削}曾派他的执事到\Place{特拜德}地区遍寻\Person{阿尔塞尼乌斯}；司铎\Person{佩基西乌斯}、\Person{赫里亚斯}的兄弟\Person{西尔瓦努斯}、\Person{塔佩纳克拉梅乌斯}，以及\Place{希普塞勒}的隐修士\Person{保禄}，这些他最先遇见的人都承认\Person{阿尔塞尼乌斯}在我们这里。我们得知后便让人将他送上船，与隐修士\Person{赫里亚斯}一同驶往下埃及地区。后来那执事又带着一些人突然返回，进入我们的隐修院搜寻同一位\Person{阿尔塞尼乌斯}，却未找到，因为正如我刚才所说，我们已将他送往下埃及地区；但他们却将我和护送他离开的隐修士\Person{赫里亚斯}带往\Place{亚历山大里亚}，并将我们带到都督面前；当时我无法否认，只好承认他还活着，并未被谋杀：那位护送他离开的隐修士也承认了同样的事。因此，神父，我将这些事告知您，免得您决定去控告\Person{亚大纳削}；因为我已声明他还活着，并曾藏匿在我们这里，而这一切已在\Place{埃及}广为人知，再也无法隐瞒。\par

\Signature{我，\Person{帕弗努提乌斯}，同隐修院的隐修士，写这封信，衷心问候您。我祈求您健康。}\par

以下是皇帝得知\Person{阿尔塞尼乌斯}被找到尚在人世时所写的信函。\par

\Person{君士坦丁}，\Person{维克托尔}，\Person{马克西姆斯}，\Person{奥古斯都}，致主教\Person{亚大纳削}。\par

68. 拜读了您智慧的来信，我遂有意回信给您坚毅的胸怀，并劝勉您致力于恢复天主子民的安宁与慈悲之心。因为在我心中，我认为最重要的事，莫过于我们应培育真理，常将正义存念于心，并特别喜悦那些行走在正直生活道路上的人。至于那些应受万般咒骂的人，我指的是那些最乖谬、不虔敬的\OriginalTerm{默肋提派}{Meletians}，他们最终因自己的愚昧而自显愚蠢，如今正借着嫉妒、骚乱和喧嚷发起无理的纷扰，从而显明他们不虔敬的性情。对此我只想说这么多。你们看，那些他们谎称已被刀剑杀害的人，如今仍在我们中间，享受着生命。那么，还有什么比这更强烈的推定，能如此明白清楚地定他们的罪，胜过那些他们宣称已被谋杀的人，如今却仍享受生命，并且将能为他们自己辩护呢？\par

但是，同样是这些默肋提派提出了更进一步的控告。他们断然肯定地说，您以无法无天的暴力冲进至圣所，抢走并打碎了一个存放于此的杯子；如果这种罪行真的发生，那无疑是极为严重的指控和极其恶劣的罪行。但这究竟是什么样的控告呢？其情形变化不定、前后不一的用意何在？竟至现在他们将同样的控告转嫁到另一个人身上，这一事实可以说比光还要清楚地表明，他们意在陷害您的智慧。在此之后，谁还愿意追随这些人呢？他们编造这样的指控以伤害他人，况且他们是在急速奔向毁灭，并且意识到自己在用虚假捏造的罪名控告您。那么，正如我所言，谁还会追随他们，从而直冲向毁灭之路呢？那条路似乎只有他们自己以为有得救和帮助的希望。然而，如果他们愿意按照纯洁的良心行事，接受最卓越智慧的引导，走上健全心志的道路，他们就很容易明白，只要他们沉溺于此类行为，自招毁灭，从天主的眷顾中就不会得到任何援助。我不应称此为对他们的苛刻判断，而只是简单的真相。\par

最后，我还要补充说，我希望这封信能由您的智慧常在公开场合宣读，好让所有人都知道，尤其要传到处事煽乱之人的耳中；因为我按公道所表达的判断，也由真实的事实所证实。因此，既然这种行为中有如此大的过犯，就让他们明白，我作出了这样的判断；并已作出这样的决定：如果他们再煽动任何类似的骚乱，我将亲自审理此事，且不是按教会的法律，而是按世俗的法律来审理，并在将来找到他们，因为他们显然可说是盗贼，不只是危害人类，更是危害天主圣道本身。愿天主永远保佑您，亲爱的弟兄！\par

69. 但为使诬告者的邪恶更充分地暴露，请看，\Person{阿尔塞尼乌斯}在他藏身之所被发现之后，也写信给我；正如\Person{依斯基拉斯}写的那封信承认了他们的控告是虚假的，阿尔塞尼乌斯的这封信更加完全地证明了他们的恶毒。\par

致蒙福的教宗\Person{亚大纳削}，\Person{阿尔塞尼乌斯}，原属\Person{默肋提乌斯}治下\Place{希普塞利特}城的那些人的主教，偕同众长老与执事，在主内致以万分的安康。\par

我们切望与天主教会和平合一（您因天主的恩宠主持此教会），并愿遵照古时规矩，服从教会的\OriginalTerm{教规}{Canon}，因此写信给您，至亲爱的教宗，并奉主名声明：我们今后绝不与那些继续处于\OriginalTerm{分裂}{schism}状态、尚未与天主教会和好的人共融，无论是主教、长老或执事。若他们想在会议中确定任何事，我们也不参与；也不与他们互通和平书信，或收受他们的来函；未经您，\OriginalTerm{都主教区}{metropolis}的主教同意，我们也不发表任何有关主教或其他一般教会问题的决议；相反，我们将遵循以往所制定的一切教规，效法\Person{阿莫尼安}、\Person{提兰诺}、\Person{普鲁西安}等主教的榜样。因此，我们恳求您的慈仁，请速赐回音，并将我们的事通知我们的同工，告诉他们我们将遵守前述决定，与天主教会和平共处，并与各区的同工团结一致。我们相信，您的祈祷蒙天主悦纳，必能藉以成就，使此和平按众主之主天主的旨意，藉着我们的主耶稣基督，坚定直至终局，永不解体。\par

您属下的神圣圣职者，我们和与我们同在的人向你致意。若天主允许，我们将很快来拜访您的慈仁。我，\Person{阿尔塞尼乌斯}，在多年内祈求您在主内安康，最蒙福的教宗。\par

70. 但对他们攻击我们的诬告，一个更有力、更明晰的证据是\Person{若望}的撤回声明。对此，蒙受天主喜爱、永志不忘的皇帝\Person{君士坦丁}可作见证，因为他知道若望如何自己指控自己，并且因为他收到了若望表达悔改的信函，遂写信给他如下。\par

\Person{君士坦丁}、\Person{马克西穆斯}、\Person{奥古斯都}致\Person{若望}。\par

从您的明智的信函所获，使我极为喜悦，因为我从这些信中学到了我久盼听到的事：您已放下一切琐屑的情感，如同您所应当的那样加入了教会的\OriginalTerm{共融}{Communion}，并且现已与最可敬的主教\Person{亚大纳削}完全和谐一致。因此，请确信我全然赞许您的行为；因为您已放弃一切纷争，行了天主所喜悦的事，并拥抱了祂的教会的合一。为使您能实现您的愿望，我认为理应准许您使用公共驿车，前来我仁慈的宫廷。请您务必不要耽搁；既然此信授予您使用公共驿车的权力，请立即到我这里来，好使您的愿望得以满足，并藉着出现在我面前，享受那您当得的喜乐。愿天主永远保守您，至亲爱的弟兄。\par

% structure: p0039 source=h2 target=subsection reason=relative-heading-rank; base=h2

\subsection{第六章。\Emph{与\Place{提洛}会议有关的文件}。}

71. 于是，阴谋就此告终。梅勒提派被击退，蒙受羞辱；然而尽管如此，\Person{优西比乌}及其同党仍未停止活动，因为他们所关心的并非梅勒提派，而是\Person{亚略}及其同党，他们担心，一旦前者的诉讼被制止，他们将再也找不到人扮演角色，借助这些人，他们就可以引入那种\OriginalTerm{异端}{heresy}。因此，他们再次煽动梅勒提派，并说服皇帝下令在\Place{提洛}重新召开一次会议。于是，\Person{狄奥尼修斯}伯爵被派往那里，\Person{优西比乌}及其同党被派给了一支卫队。\Person{马卡里乌斯}也在士兵押解下，作为囚犯被送往\Place{提洛}。皇帝写信给我，对我下了严厉的命令，因此，无论我多么不愿意，也只好动身。整个阴谋可从埃及主教们所写的信中加以理解；但首先有必要叙述他们是如何从起初策划这一切的，这样便能认清他们对我所施的恶意与邪恶。在埃及、利比亚和五城，有将近一百位主教；他们中没有一个人指控我；没有一位司铎对我有所非难；百姓中也没有一个人说我的坏话；而是那些被\Person{伯多禄}所驱逐的梅勒提派和亚略派，共同分派了阴谋，前者声称有权控告我，后者则有权审判此案。我反对\Person{优西比乌}及其同党，因为他们是我的仇敌，因着异端的缘故；其次，我以如下方式指出，那个所谓的控告者根本就不是一位司铎。当\Person{梅勒提乌斯}被接纳进入\OriginalTerm{共融}{communion}时（巴不得他从未被接纳！），有福的\Person{亚历山大}知道他的狡诈，要求他提供一份名册，列出他声称在埃及拥有的主教，以及在\Place{亚历山大里亚}本城的司铎和执事，也列出他在乡间地区所拥有的人。教宗\Person{亚历山大}这样做，是免得\Person{梅勒提乌斯}在领受了教会的自由后，提出许多人，并这样通过欺诈程序，不断把他所喜悦的人强塞给我们。因此，他开列了以下有关埃及人的名册。\par

\Emph{\Person{梅勒提乌斯}向主教\Person{亚历山大}提交的一份名册}。\par

我，\Place{利科波利斯}的\Person{梅勒提乌斯}，\Place{安提诺波利斯}的\Person{路爵}，\Place{赫尔莫波利斯}的\Person{法西莱乌斯}，\Place{库塞}的\Person{阿喀琉斯}，\Place{狄奥斯波利斯}的\Person{阿摩尼乌斯}。\par

在\Place{托勒密斯}，\Place{滕提拉}的\Person{帕基梅斯}。\par

在\Place{马西米安诺波利斯}，\Place{科普托斯}的\Person{狄奥多鲁斯}。\par

在\Place{底比斯}，\Place{赫尔梅特斯}的\Person{卡莱斯}，\Place{上基诺波利斯}的\Person{科卢图斯}，\Place{俄克西林库斯}的\Person{佩拉吉乌斯}，\Place{赫拉克利奥波利斯}的\Person{伯多禄}，\Place{尼洛波利斯}的\Person{特翁}，\Place{莱托波利斯}的\Person{依撒格}，\Place{尼基奥波利斯}的\Person{赫拉克利德斯}，\Place{克娄巴特里斯}的\Person{依撒格}，\Place{阿尔塞诺伊提斯}的\Person{梅拉斯}。\par

在\Place{赫利奥波利斯}，\Place{莱昂托波利斯}的\Person{阿摩斯}，\Place{阿特里比斯}的\Person{伊西翁}。\par

在\Place{法尔贝图斯}，\Place{布巴斯图斯}的\Person{哈尔波克拉提翁}，\Place{法库塞}的\Person{梅瑟}，\Place{培琉息乌姆}的\Person{卡利尼库斯}，\Place{塔尼斯}的\Person{欧德蒙}，\Place{特姆伊斯}的\Person{厄弗辣因}。\par

在\Place{塞斯}，\Place{基诺波利斯}和\Place{布西里斯}的\Person{赫尔梅翁}，\Place{塞贝尼图斯}的\Person{索特里库斯}，\Place{普特内吉斯}的\Person{皮尼努特斯}，\Place{梅特利斯}的\Person{克罗尼乌斯}，\Place{亚历山大里亚}地区的\Person{阿伽特亚蒙}。\par

在\Place{孟斐斯}，\Person{若翰}，他是奉皇帝之命与总主教在一起的。这些是埃及的。\par

而他在\Place{亚历山大里亚}拥有的圣职人员如下：司铎\Person{阿波罗尼乌斯}、司铎\Person{依勒内}、司铎\Person{狄奥斯库若}、司铎\Person{提兰努斯}。执事：执事\Person{弟茂德}、执事\Person{安蒂诺乌斯}、执事\Person{赫菲斯提翁}。以及\Place{帕雷姆波莱}的司铎\Person{马卡里乌斯}。\par

72. \Person{Meletius} 确实亲自向 \Person{Alexander} 主教提交了这些，但他并未提及 \Person{Ischyras} 此人，也从未声称他在 \Place{Mareotis} 有任何神职人员。然而我们的敌人并未停止其企图，那个本非司铎的人仍被冒充为司铎，因为伯爵准备对我们施压，士兵们也将我们驱来赶去。但即便如此，天主的恩宠仍然得胜；因为他们无法在圣爵一事上定罪 \Person{Macarius}；而 \Person{Arsenius}，就是他们宣称被我杀害的那位，竟活着站在他们面前，显明了他们指控的虚假。因此，当他们无法定罪 \Person{Macarius} 时，\Person{Eusebius} 及其同伙因为失去了他们一直追捕的猎物而恼怒不已，便说服了同属他们一伙的伯爵 \Person{Dionysius} 派人前往 \Place{Mareotis}，看看能否在那里找到一些对司铎不利的事，更确切地说，好让他们能在我们缺席的情况下，远距离地任意编造他们的阴谋：因为这就是他们的目的。然而——当我们指出前往 \Place{Mareotis} 实属多此一举（因为他们不应假装那些他们已盘问了如此之久的陈述尚有欠缺，现在也不应再拖延此事；他们想说什么，他们以为能说的，都已说完，如今不知如何是好，才找借口）；或者，如果你们非要去 \Place{Mareotis} 不可，那至少不要派那些有嫌疑的人去——伯爵被我的论证说服，同意不派那些有嫌疑的人；但他们所做的与我提议的完全相反，因为我曾因 \OriginalTerm{亚略异端}{Arian heresy} 而反对的那些人，正是这伙人即刻出发了，即 \Person{Diognius}、\Person{Maris}、\Person{Theodorus}、\Person{Macedonius}、\Person{Ursacius} 与 \Person{Valens}。他们又写信给 \OriginalTerm{埃及总督}{Prefect of Egypt}，并提供了军事护卫；而最引人注目也最为可疑的是，他们将被指控的一方 \Person{Macarius} 留了下来，交由士兵看守，却带上了那个原告。如今，有谁还看不穿这阴谋？有谁还不清楚 \Person{Eusebius} 及其同伙的邪恶？因为若真有必要在 \Place{Mareotis} 进行司法调查，那么被告也应被派往那里。但如果他们此行不为调查，为何要带上原告？他未能证明其事，这已足够。但他们这样做，是为了能够对他们无法在当面定罪的那位缺席司铎实施阴谋，并任意捏造一个方案。因为当 \Place{Alexandria} 和整个地区的司铎们因他们单独行事而指责他们，并要求自己也能参与他们的诉讼时（因为他们说，自己既了解案情，也熟知 \Person{Ischyras} 其人的底细），他们却不允许；尽管他们身边有 \OriginalTerm{埃及总督}{Prefect of Egypt} \Person{Philagrius}，一个背教者，以及异教士兵，做着一项连 \OriginalTerm{慕道者}{Catechumens} 都不宜旁观的调查，他们却不接纳神职人员，生怕在那边也像在 \Place{Tyre} 一样，有人会揭穿他们。\par

73. 但尽管采取了这些预防措施，他们仍未能逃脱败露：因为城中和 \Place{Mareotis} 的司铎们察觉到了他们的邪恶用心，于是向他们递交了以下抗议书。\par

致从 \Place{Tyre} 来的众主教 \Person{Theognius}、\Person{Maris}、\Person{Macedonius}、\Person{Theodorus}、\Person{Ursacius} 与 \Person{Valens}，来自 \Place{Alexandria} 天主教会、隶属于至可敬的主教 \Person{Athanasius} 的众司铎与执事。\par

你等既然来到此地，且带来了原告，就理当同时带来司铎 \Person{Macarius}；因为圣经规定，审判应如此设立，使原告与被告一同站立。但你等既未带来 \Person{Macarius}，我们的至可敬主教 \Person{Athanasius} 也未曾与你等同来，我们便为自己主张参与调查的权利，以便亲眼得见调查公正进行，也让自己确信真相。然而你等拒绝允许此事，竟只愿与\OriginalTerm{埃及总督}{Prefect of Egypt}和原告一起，为所欲为，我们坦白相告：我们在这事上看到了邪恶的嫌疑，并察觉你等的到来不过是阴谋集团与共谋的举动。因此我们向你等寄出此信，作为一个真实公会议前的证词，好让众人皆知，你等进行的是一项\Emph{\OriginalTerm{单方面}{ex parte}}的、为私利而行的程序，除了图谋反对我们之外，别无所求。此信副本，恐被你等藏匿，我们已交给\OriginalTerm{奥古斯都财务官}{Controller of Augustus} \Person{Palladius}。因为你等已有的所作所为使我们对你等生疑，并推测你等今后亦将有类似行径。\par

\Signature{司铎 \Person{Dionysius}} 已递交此信。\Signature{司铎 \Person{Alexander}}、\Signature{司铎 \Person{Nilaras}}、\Signature{司铎 \Person{Longus}}、\Signature{司铎 \Person{Aphthonius}}、\Signature{司铎 \Person{Athanasius}}、\Signature{司铎 \Person{Amyntius}}、\Signature{司铎 \Person{Pistus}}、\Signature{司铎 \Person{Plution}}、\Signature{司铎 \Person{Dioscorus}}、\Signature{司铎 \Person{Apollonius}}、\Signature{司铎 \Person{Sarapion}}、\Signature{司铎 \Person{Ammonius}}、\Signature{司铎 \Person{Gaius}}、\Signature{司铎 \Person{Rhinus}}、\Signature{司铎 \Person{Æthales}}。\par

执事：\Signature{执事 \Person{Marcellinus}}、\Signature{执事 \Person{Appianus}}、\Signature{执事 \Person{Theon}}、\Signature{执事 \Person{Timotheus}}、\Signature{执事 另一位 \Person{Timotheus}}。\par

74. 这就是那封信，这些是城中神职人员的名字；以下是 \Place{Mareotis} 的神职人员所写的信，他们知道那个原告的品行，并且曾在我的巡访期间与我同在。\par

致天主教会真福众主教的神圣公会议，\Place{Mareotis} 的全体司铎与执事，在主内问候。\par

我们深知经上记载着：\BibleQuote{讲论你眼所见}，又载着：\BibleQuote{作假见证的，必不免受罚}，因此我们见证亲眼所见的事，尤其是因为针对我们主教\Person{亚大纳削}的阴谋，使得我们的见证成为必要。我们诧异\Person{伊斯基拉斯}竟会被列为教会圣职者之一，这是我们首先要提及的一点。\Person{伊斯基拉斯}从来就不是教会的圣职者；以前他自称是\Person{科卢托}的\OriginalTerm{司铎}{Presbyter}，除自己的亲属外，无人相信他。因他从未有过教堂，也从未被离他村庄不远的人们认作神职人员，正如前述，只有他自己的亲属认他。尽管他僭取了这称号，但在\Place{亚历山大里亚}召开的会议上，当着我们教父\Person{何西乌斯}的面，他被废黜，并以\OriginalTerm{平信徒}{Layman}身份被接纳\OriginalTerm{共融}{Communion}，此后他一直如此，从他妄称的司铎品级上跌落下来。至于他的品行，我们认为无须多论，因为众人皆可了解。但他诬告我们主教\Person{亚大纳削}打破一只杯并掀翻一张桌子，我们不得不就此向您们说明。我们已说过，他在\Place{马雷奥提斯}从未有过教堂；我们在天主面前作证，我们的主教或任何随同他的人，都没有打破杯，也没有掀翻桌子；所有关于此事的指控纯属诽谤。我们这样说，并非因当时不在主教身边，因为他巡视\Place{马雷奥提斯}时，我们全都与他在一起，他从不独自出行，而是有我们全体司铎和\OriginalTerm{执事}{Deacon}，以及相当多的民众陪同。因此，我们作为他每次巡视我们时的在场者，作出这些断言，并证明从未有过杯被打破或桌子被掀翻，整个故事是假的，控告者自己也曾亲笔作证。因为当他与\OriginalTerm{梅勒提派}{Meletian}同去，并向他们报告了这些攻击我们主教\Person{亚大纳削}的事后，他想要被接纳共融，却没有被接纳，尽管他亲笔写下并承认所有这些事都不真实，他是受了某些人的唆使才这样说。\par

75. 因此，\Person{提奥格尼乌斯}、\Person{狄奥多若}、\Person{玛里斯}、\Person{马切多尼乌斯}、\Person{乌尔撒奇乌斯}、\Person{瓦伦斯}及其同伙也来到\Place{马雷奥提斯}，当他们发现这些事全非事实，而很可能他们构陷我们主教\Person{亚大纳削}的虚假指控会被揭穿，\Person{提奥格尼乌斯}等人既是主教的对头，便指使\Person{伊斯基拉斯}的亲属和某些\OriginalTerm{亚略派}{Arian}的狂徒随意说他们想说的话。因为民众中没有人反对主教；但这些人因惧怕\Place{埃及}总督\Person{菲拉格里乌斯}，并在威胁之下，靠着亚略派狂徒的支持，便为所欲为。当我们前去驳斥诽谤时，他们不许我们那样做，反将我们赶出，却让他们所中意的人参与图谋，与他们密谋，并以总督\Person{菲拉格里乌斯}的威势影响他们。通过他的手段，他们阻止我们到场，以免我们发现那些受他们唆使的人到底是教会的成员还是亚略派的狂徒。而且，亲爱的教长们，你们也知道，正如你们教导我们的，敌人的见证毫无价值。我们所说的真实，\Person{伊斯基拉斯}的亲笔供词可证，事实本身亦然，因为当我们清楚并未发生过所假称的事时，他们带着\Person{菲拉格里乌斯}，为能靠刀剑的威慑和恐吓，编织他们所图谋的任何诡计。这些事我们在天主台前作证；我们作这些断言，是因深知将来有天主的审判；我们虽渴愿众人亲赴你们面前，但只派部分代表即可，这样此信便可代替未到之人的出席。\par

\Signature{我，司铎因格尼乌斯，在主内向你们致敬，亲爱的教长们。}\newline{}\Signature{司铎狄昂}\newline{}\Signature{司铎阿摩纳斯}\newline{}\Signature{司铎赫拉克利乌斯}\newline{}\Signature{司铎博孔}\newline{}\Signature{司铎特里丰}\newline{}\Signature{司铎伯多禄}\newline{}\Signature{司铎希厄拉克斯}\newline{}\Signature{司铎撒拉皮翁}\newline{}\Signature{司铎玛尔谷}\newline{}\Signature{司铎托勒迈翁}\newline{}\Signature{司铎加约}\newline{}\Signature{司铎狄奥斯库若}\newline{}\Signature{司铎德默特琉斯}\newline{}\Signature{司铎迪尔索斯}\par

\Signature{执事：}\newline{}\Signature{执事皮斯图斯}\newline{}\Signature{执事阿颇罗}\newline{}\Signature{执事塞拉斯}\newline{}\Signature{执事皮斯图斯}\newline{}\Signature{执事波利尼库斯}\newline{}\Signature{执事阿摩尼乌斯}\newline{}\Signature{执事茂鲁斯}\newline{}\Signature{执事赫菲斯托斯}\newline{}\Signature{执事阿颇罗}\newline{}\Signature{执事梅托帕斯}\newline{}\Signature{执事阿颇罗}\newline{}\Signature{执事塞拉帕斯}\newline{}\Signature{执事梅利弗通古斯}\newline{}\Signature{执事路基乌斯}\newline{}\Signature{执事格雷戈拉斯}\par

\Emph{76. 致审计官及当时埃及总督菲拉格里乌斯的同一信函。}\par

致\Person{弗拉维乌斯·菲拉格里乌斯}、杜塞纳留斯兼宫廷官员兼审计官\Person{弗拉维乌斯·帕拉迪乌斯}，以及粮食专员、我主最显赫的神圣御前总督的百人长\Person{弗拉维乌斯·安东尼努斯}；来自\Place{马雷奥提斯}教区、隶属于最可敬主教\Person{亚大纳削}的天主教会，由下列署名司铎及执事呈上此证：\par

鉴于\Person{提奥格尼乌斯}、\Person{玛里斯}、\Person{马切多尼乌斯}、\Person{狄奥多若}、\Person{乌尔撒奇乌斯}和\Person{瓦伦斯}，仿佛受在\Place{提洛}集会全体主教派遣，来到我们教区，声称奉命调查某些教会事务，其中提及擘开主的杯之事，此信息由他们带来的\Person{伊斯基拉斯}提供，\Person{伊斯基拉斯}自称是司铎，实则并非——因他是受僭取主教职的司铎\Person{科卢托}祝圣，后来在一次由\Person{何西乌斯}及同他在一起的主教们出席的全体会议上，被命恢复先前司铎之职，随即所有由\Person{科卢托}祝圣者均恢复原品级，\Person{伊斯基拉斯}亦证实为平信徒——至于他自称拥有的教堂，根本不是教堂，而是一间属于名叫\Person{伊西翁}的孤儿的极小私宅；——为此，我们呈上此证，恳求你们，藉着全能天主，以及我主\Person{君士坦丁·奥古斯都}和其最显赫的儿子凯撒们，将这一切呈报于他们的虔敬御览。因他既非天主教会的司铎，亦未拥有教堂，更从未有过杯被打破；整个故事纯属虚假捏造。\par

于最显赫的贵族\Person{儒略·君士坦提乌斯}——最虔诚的皇帝\Person{君士坦丁·奥古斯都}的兄弟——与最显赫的\Person{鲁菲努斯·阿尔比努斯}同任执政官之年，\OriginalTerm{托特月}{Thoth}第十日。\par

这些便是\OriginalTerm{长老}{presbyter}们的信函。\par

77. 以下也是那些随我们同来\Place{提洛}的主教们，在得知阴谋与诡计后所写的信函和抗议书。\par

致聚集于\Place{提洛}的诸位主教，至可敬的主教们：从\Place{埃及}随同\Person{亚大纳削}前来的\OriginalTerm{天主教会}{Catholic Church}人士，在主内向你们致意。\par

我们认为，\Person{犹西比乌斯}、\Person{狄奥格尼乌斯}、\Person{马里斯}、\Person{纳尔基索斯}、\Person{狄奥多鲁斯}、\Person{帕特罗斐鲁斯}及其同党所策划针对我们的阴谋，已不再隐晦。从一开始，我们便通过我们的同工\Person{亚大纳削}，一致反对在他们面前进行聆讯，因为我们深知，即使只有一个敌人，更不用说众多敌人，在场就足以扰乱和损害案件的审理。你们自己也清楚，他们所怀的敌意，不仅针对我们，也针对所有\OriginalTerm{正统}{orthodox}信徒：他们因着\Person{亚略}的疯狂及其不虔敬的教义，向所有人发起攻击，图谋陷害众人。当我们凭着对真理的信心，想要揭露\OriginalTerm{默肋提安派}{Meletians}如何以虚伪控告教会时，\Person{犹西比乌斯}及其同党千方百计拦阻我们的陈述，竭力排除我们的证词，威胁那些凭良心判断的人，侮辱其他人，其唯一目的便是实施他们对我们的阴谋。至可敬的主教们，你们敬虔的心或许尚不知晓他们的阴谋，但我们认为如今已是昭然若揭。事实上，他们已经自己显露无遗：他们渴望把那批我们怀疑的人派往\Place{马雷奥提斯}，以便趁我们不在场、滞留此地的时机，扰乱民众，达成他们的企图。他们知道\OriginalTerm{亚略派狂人}{Arian madmen}、\OriginalTerm{科路提安派}{Colluthians}和\OriginalTerm{默肋提安派}{Meletians}乃是\OriginalTerm{天主教会}{Catholic Church}的敌人，因此急欲派遣他们，好在敌人面前随意策划对付我们的计谋。而且，在此地的一些\OriginalTerm{默肋提安派}{Meletians}，早在四天前（他们知悉聆讯即将举行），便于傍晚派出同党充当信使，为要从\Place{埃及}各地召集\OriginalTerm{默肋提安派}{Meletians}涌入\Place{马雷奥提斯}——因为那里原本一个也没有——并从其他地方召集\OriginalTerm{科路提安派}{Colluthians}和\OriginalTerm{亚略派狂人}{Arian madmen}，预备他们作证攻击我们。你们也知道，\Person{依希拉斯}本人曾在你们面前承认，他的会众不超过七人。因此，当我们听说，他们在按自己的心意做了种种准备、派出了这些可疑之人以后，便挨个走访你们，要求你们签署，好使这一切看似是经过你们全体同意的，为此我们火速写信给你们，呈上我们的这份证词，声明我们是阴谋的受害者，这阴谋由他们一手造成的，我们所受的苦皆为他们的缘故，并要求你们存敬畏天主的心，谴责他们未经我们同意擅自派遣人员的行为，拒绝签署，免得他们假装那些事是奉你们之命而行，实际上不过是他们私下策划的勾当。凡在基督里的人，的确不应顾及人情，而应置真理于万有之上。不要害怕他们用来威胁所有人的恐吓，也不要怕他们的阴谋，而要敬畏天主。如果确实有必要派人前往\Place{马雷奥提斯}，我们也应当在那里与他们同去，好能指证教会的敌人，指出那些异己分子，并使调查保持公正。你们知道，\Person{犹西比乌斯}及其同党设法呈递了一封信件，声称来自\OriginalTerm{科路提安派}{Colluthians}、\OriginalTerm{默肋提安派}{Meletians}和\OriginalTerm{亚略派}{Arians}，矛头直指我们。然而，很明显，这些\OriginalTerm{天主教会}{Catholic Church}的敌人并没有述说关于我们的任何真话，而是竭尽所能地反对我们。天主的法律禁止仇敌充当证人，也不许他扮演审判官。因此，鉴于你们终将在审判之日交账，请接收这份证词，认清那针对我们的阴谋，若有他们向你们提出请求，万勿行任何有害于我们的事，或是参与\Person{犹西比乌斯}及其同党的图谋。因为你们知道，正如我们先前所说，他们是我们的仇敌，你们也深知\Person{凯撒利亚的犹西比乌斯}去年缘何变成如此。我们祈祷你们身体健康，至深敬爱的主教们。\par

78. 致最显赫的伯爵\Person{弗拉维·狄奥尼修斯}，寄自从\Place{埃及}前来\Place{提洛}的\OriginalTerm{天主教会}{Catholic Church}众主教。\par

我们认为，由\Person{优西比乌}、\Person{德奥尼乌}、\Person{马里斯}、\Person{纳尔基索}、\Person{德奥多若}、\Person{帕特洛斐洛}及其同伙所策划的针对我们的阴谋，已不再有任何疑问。从一开始，我们便通过我们的同工司铎\Person{亚大纳削}，一致反对在他们面前进行审理，因为我们深知，哪怕只有一个敌人，更不用说许多敌人，在场都足以扰乱和损害案件的审理。他们对我们的敌意是明显的，不仅针对我们，而且针对所有正统人士，因为他们攻击所有人，对所有人策划阴谋。当我们因着对真理的信心，渴望揭露\OriginalTerm{梅勒提派}{Meletians}对教会所施行的谎言时，\Person{优西比乌}及其同伙千方百计地打断我们的陈述，竭力压制我们的见证，威胁那些作出诚实判断的人，并侮辱其他人，其唯一目的就是实现他们对我们所定的计谋。你们的良善或许不知道他们对我们所策划的阴谋，但我们相信现在已经显明了。事实上，他们自己已清楚地表露了这一点：他们想要派遣他们那一伙人中我们有所怀疑的人到\Place{玛瑞奥提斯}去，好让我们不在场而留在这里时，他们可以去扰乱民众，达成他们的意愿。他们知道\OriginalTerm{亚略派狂人}{Arian madmen}、\OriginalTerm{科鲁提派}{Colluthians}和\OriginalTerm{梅勒提派}{Meletians}是教会的敌人，因此急于派遣他们，以便在敌人面前，随心所欲地设计各种计谋来对付我们。而在这里的\OriginalTerm{梅勒提派}{Meletians}，甚至在四天前（因为他们知道这次审讯即将举行），就在傍晚派出了他们党内的两个人作为信使，去从\Place{埃及}各地召集\OriginalTerm{梅勒提派}{Meletians}到\Place{玛瑞奥提斯}，因为那里原本完全没有他们的人，并从其他地方召集\OriginalTerm{科鲁提派}{Colluthians}和\OriginalTerm{亚略派狂人}{Arian madmen}，并预备他们发言反对我们。你们的良善清楚知道，他本人曾在你们面前承认，他的会众中不超过七个人。因此，当我们听说，他们在按自己的意愿做好了对付我们的准备，并派遣了这些可疑之人后，便四处奔走于每一位主教，要求他们签署，以便显得这是出于众人同意的时候；我们便急忙将此事提交给你们的尊长，并呈上这份我们的见证，表明我们是阴谋的受害者，我们正是因着他们并通过他们而遭受这一切，并要求你们，心中怀着对天主的敬畏，以及我们最尊贵的皇帝虔诚的诫命，不再容忍这些人，而是谴责他们未经我们同意就随意派遣人员的行为。\par

\Signature{我，\Person{亚当安提}主教，已签署这封信，\Person{伊斯基拉斯}、\Person{阿蒙}、\Person{伯多禄}、\Person{阿摩尼安}、\Person{提兰诺}、\Person{陶里诺}、\Person{撒拉帕蒙}、\Person{厄琉良}、\Person{哈波克拉提翁}、\Person{梅瑟}、\Person{欧普塔托}、\Person{阿努比翁}、\Person{撒普良}、\Person{阿波罗纽}、\Person{伊斯基良}、\Person{阿尔贝提翁}、\Person{波塔蒙}、\Person{帕弗努提}、\Person{赫拉克利德}、\Person{德奥多若}、\Person{阿加塔蒙}、\Person{加约}、\Person{庇斯托}、\Person{阿塔斯}、\Person{尼孔}、\Person{佩拉纠}、\Person{德翁}、\Person{帕尼努提}、\Person{诺诺}、\Person{阿里斯托}、\Person{德奥多若}、\Person{伊利奈}、\Person{布拉塔蒙}、\Person{斐理伯}、\Person{阿颇罗}、\Person{狄奥斯科若}、\Place{狄奥斯波利}的\Person{提摩德}、\Person{玛加略}、\Person{赫拉克拉蒙}、\Person{克若尼}、\Person{米斯}、\Person{雅各伯}、\Person{阿里斯托}、\Person{阿尔德米多}、\Person{斐内斯}、\Person{普撒伊斯}、\Person{赫拉克利德}。}\par

\Emph{同一会议的另一封信。}\par

79. 从\Place{埃及}来到\Place{提洛}的天主教会主教们，致最尊贵的\Person{弗拉维·狄奥尼修}伯爵。\par

我们察觉到，通过\Person{优西比乌}、\Person{纳尔基索}、\Person{弗拉基洛}、\Person{德奥尼乌}、\Person{马里斯}、\Person{德奥多若}、\Person{帕特洛斐洛}及其同伙的阴谋诡计（我们起初曾想对他们提出反对，但未获准许），许多阴谋和诡计正针对我们形成，因此我们不得不诉诸这份呼吁。我们也注意到，有人在为\OriginalTerm{梅勒提派}{Meletians}大发热心，并且一场针对我们在\Place{埃及}的天主教会的阴谋正在我们身上酝酿。因此，我们将这封信呈交于你，恳请你铭记那保护我们最尊贵、最虔敬的皇帝\Person{君士坦丁}王国的全能天主权能，并将涉及我们的事务的审理留待最尊贵的皇帝本人。因为，既然你是受陛下委派，那么在我们向他的虔敬上诉时，将此事留待他处置，是合理的。我们再也无法忍受继续成为上述\Person{优西比乌}及其同伙的阴险诡计的对象，因此我们要求将案件留待最尊贵和天主所爱的皇帝，在他面前我们将能够陈述我们自身和教会的正当要求。我们深信，当他的虔敬听到我们的案件时，他不会定我们的罪。因此，我们再次藉着全能的天主，并藉着我们最尊贵的皇帝——他连同他虔敬的子孙，多年来一直如此常胜繁荣——恳请你不要再继续行事，也不要容许自己在公会议中对我们的事务作出任何行动，而是将其审理留待他的虔敬。我们也已向我的主人们，正统主教们，作了同样的陈述。\par

80. \Place{得撒洛尼}主教\Person{亚历山大}收到这些信后，写信给伯爵\Person{狄奥尼修}如下。\par

主教\Person{亚历山大}致我的主人\Person{狄奥尼修}。\par

我看出显然有人合谋反对\Person{亚大纳削}；因为他们执意派遣所有他所反对的人，却没有任何通知我们，虽然先前已商定，关于应派遣谁，我们应一起考虑。因此，要小心不可鲁莽行事，（因为他们大为惊恐地来到我这里，说那些野兽已经醒来，将要扑向他们；因他们听到报告说，\Person{若望}已经派遣了某些人），免得他们抢先于我们，并任意图谋。你们知道，那与教会为敌的科卢提派，以及亚略派，和梅肋提安派，都联合在一起，能作许多恶事。所以，考虑什么是最好的行动，免得发生祸患，我们因判断不公而遭受责备。这些人也大受怀疑，恐怕他们因忠于梅肋提安派，会经过那些其主教在此地的教会，在他们中间引起恐慌，从而扰乱整个\Place{埃及}。因为他们看到，此事已经在很大程度上发生。\par

因此，\Person{狄奥尼修}伯爵写信给\Person{优西比乌}及其同伙如下。\par

81. 这是我已经向我的主人们，\Person{福拉基路}及其同伙提到的，\Person{亚大纳削}曾前来诉说，他所反对的那些人正是被派遣的那些人；并大声疾呼他受了冤屈和欺骗。我灵魂的主\Person{亚历山大}也为此事写信给我；为使你们看到他的仁慈所说的是合理的，我附上他的信，供你们阅读。也请记起我先前给你们写过的话：我向你们，我的主人们，强调过，所派遣的人应当由全体共同投票和决定而委任。因此，要小心我们的行径不要遭受指责，也不要给那些乐于挑剔我们的人留下公道的把柄。因为正如控方不应遭受任何压迫，辩方同样不应。我认为，当我的主\Person{亚历山大}明显不赞成我们所做的事时，对我们便有不小的指责理由。\par

82. 事情正在这样进行时，我们离开了他们，如同离开奸诈人的集会 \BibleRef{耶 9:2}，因为他们任意妄为，而世上无人不知，\OriginalTerm{单方面的}{ex parte}程序不能成立。这是神圣法律所决定的；因为当那位真福宗徒遭受类似的阴谋并被带上公堂时，他要求说：\BibleQuote{那些从\Place{亚细亚}来的犹太人，本当到你面前控告我，若他们有什么反对我的事。}\BibleRef{宗 24:18} 那时，\Person{斐斯托}也一样，当犹太人想要对他施行像这些人如今对我施行的阴谋时，说：\BibleQuote{罗马人的惯例，是在被告人与控告者当面对质，并准予针对所受的指控为自己辩护之前，不将任何人交付处死。}\BibleRef{宗 25:16} 但\Person{优西比乌}及其同伙竟胆敢曲解法律，比那些作恶者更不义。因为他们起初并没有私下行事，而是因我们到场而发现自己势弱，便立即退出去，如同犹太人一样，独自商议如何除灭我们并引入他们的异端，就像那些人要求释放\Person{巴拉巴}一样。正是为了这个目的，正如他们自己承认的，他们做了这一切。\par

83. 这些情形虽已足以使我们辩白，但为使这些人的邪恶和真理的自由得到更充分的彰显，我不反对再予重复，以表明他们的行事方式自相矛盾，如同暗中策划的人攻击自己的朋友，并且在想要毁灭我们时，如同疯人一样自伤己身。因为他们在调查\OriginalTerm{圣体奥迹}{Mysteries}的事时，询问了犹太人，盘问了慕道者；他们问：\Quote{当\Person{玛加略}进来推翻祭台时，你们在哪里？}他们回答：\Quote{我们在里面。}但是如果慕道者当时在场，就不能有祭献。再者，虽然他们到处写信说，\Person{玛加略}进来推翻了所有东西，当时司铎正站着举行圣体奥迹，然而当他们任意询问人，问他们：\Quote{当\Person{玛加略}冲入时，\Person{依斯基拉}在哪里？}那些人回答说他正在一间小屋中卧病。那么，躺着的人不可能站着，在小屋里卧病的人也不可能奉献祭献。此外，\Person{依斯基拉}声称一些书籍被\Person{玛加略}烧毁，但那些被教唆作证的人声明，根本没有这类事发生，而是\Person{依斯基拉}说谎。最引人注目的是，虽然他们再次到处写信说，能够作证的人被我们藏匿起来了，但这些人却出来露面，他们询问他们，当看到各方面都证实他们是诽谤者，行这事是暗地里且随心所欲时，他们并不感到羞耻。因为他们用暗号暗示证人，而总督威胁他们，士兵用剑刺他们；但主揭示了真理，显明他们是诽谤者。因此，他们也将自己保留的审讯记录藏匿起来，并命令执笔的人把它们藏起来，不要交给任何人。但这事他们也落空了；因为执笔的人是\Person{鲁福}，他现在是\Place{奥古斯塔里卡总督区}的公共行刑人，能够为此作证；而\Person{优西比乌}及其同伙通过他们自己的朋友将这些记录送到\Place{罗马}，\Person{犹利乌}主教将其转交给我。现在他们恼怒，因为我们获得并阅读了他们想要隐藏的东西。\par

84. 他们的阴谋既是如此，他们的所作所为的缘由也很快就明白显示出来。因为他们离开时，带着亚略派一同前往\Place{耶路撒冷}，在那里接纳他们共融，并为此发出一封信，其部分内容及开头如下。\par

由天主恩宠聚集于\Place{耶路撒冷}的神圣公会议，致书于在\Place{亚历山大里亚}的天主的教会，以及在\Place{埃及}全地、\Place{特拜德}、\Place{利比亚}、\Place{五城地区}并普世的主教、司铎和执事，在主内祝你们平安。\par

我们来自不同行省，聚集参加为祝圣救主殉道堂而举行的大会——该堂乃因我们最蒙天主喜爱的\Person{君士坦丁皇帝}的热心，被指定用于敬拜万王之王的天主及其\Person{基督}——天主的恩宠赐予我们心中更丰盛的喜乐；我们最蒙天主喜爱的皇帝本人也通过信函带给我们这种喜乐，他在信中激励我们行正直的事，从天主的教会中除去一切嫉妒，远远驱除我们中间所有恶意（天主的肢体此前因恶意而被分裂），并敦促我们以纯朴和平和的心接纳\Person{亚略}及其同伙，他们因那一切美善之敌——嫉妒——而暂时被排除在教会之外。我们最虔诚的皇帝也在信中证明他们信仰的正统性，这是他亲自从他们那里得知的，他亲自从他们口头领受了他们的信仰宣认，并现在通过在他的信函后附上这些人的正统信仰书面声明，向我们显明。\par

85. 凡听闻这些事的人，必能看穿他们的诡诈。因为他们对自己所行的事毫不隐瞒；除非他们并非故意地承认了真理。因为，若我是阻止\Person{亚略}及其同伙进入教会的障碍，而在我因他们的阴谋受苦时他们却被接纳，那么除了这些事是为他们的缘故而行，以及他们对我的一切攻击，并他们捏造的关于打碎圣爵和谋杀\Person{亚尔塞尼乌斯}的故事，其唯一目的就是将不虔敬引入教会，并防止他们被定为异端者之外，还能得出什么别的结论呢？因为这就是皇帝从前在给我的信中所威胁的。他们竟不耻于如此写信，并断言那些被整个大公会议绝罚的人持有正统信仰。他们既然肆无忌惮地言行任何事，也就不怕\Quote{在一个角落里}聚集，为要尽他们所能地推翻如此伟大公会议的权威。\par

此外，他们为买得伪证所付出的代价，更充分地显明了他们的邪恶与不虔的意图。\Place{马雷奥提斯}，如我已说过，是\Place{亚历山大里亚}的一个乡区，那里从未有过主教或\OriginalTerm{郊区主教}{Chorepiscopus}；全区的教会都隶属于亚历山大里亚主教，每位司铎负责一个较大的村庄，数目约有十个或更多。而\Person{伊斯奇拉斯}居住的村庄非常小，居民极少，从未在那里建过教堂，仅在邻近村庄建有。然而，他们竟不顾古老惯例，决定为此地任命一名主教，不仅限于此，甚至任命了一位连司铎都不是的人。他们明知此举极不寻常，但因受制于他们为换取他对我作伪证而许下的承诺，不得不妥协至此，以免那个无耻之徒，若被他们忘恩负义地对待，会揭露真相，从而暴露\Person{优西比乌}及其同伙的邪恶。尽管如此，他没有教堂，也没有民众服从他，反而被众人藐视，\Quote{像狗一样}；虽然他们甚至促使皇帝写信给\OriginalTerm{总税务官}{Receiver-General}（凡事都在他们掌握之中），命令为他修建一座教堂，以便等他拥有教堂后，他关于圣爵和祭台的陈述能显得可信。他们还立即使他被任命为主教，因为若他没有教堂，甚至连司铎都不是，他就会显得是个诬告者和整件事的捏造者。无论如何，他没有民众，甚至他自己的亲属也不服从他，而他保留的名号徒有虚名，同样地，以下他保存的信函也毫无效力，他将之展示出来，作为自己和\Person{优西比乌}及其同伙彻底邪恶的暴露。\par

\Emph{总税务官的信函}\par

\Person{弗拉维乌斯·赫梅里乌斯}向\Place{马雷奥提斯}的收税员致意。\par

\Person{伊斯奇拉斯}司铎曾向我等主上——\OriginalTerm{奥古斯都}{Augusti}与\OriginalTerm{凯撒}{Cæsars}——的敬虔祈求，希望在属于\Place{塞康塔鲁鲁斯}的\Place{伊雷内}地区修建一座教堂，他们的神圣威严已命此事尽速成就。因此，你们一收到神圣敕令的副本（该敕令以应有的敬意置于上方）以及呈递于卑职面前的报告后，务须快速为其制作摘要，并将其录入诏令书，以使这神圣命令得以执行。\par

86. 当他们如此策划阴谋时，我前去向皇帝陈述了\Person{优西比乌}及其同伙的不义行为，因为正是他命令召开公会议，并由他的伯爵主持。他听了我的报告后，大为震动，并写信给他们如下：\par

\Person{君士坦丁}，胜利者，至大者，奥古斯都，致聚集于\Place{提洛}的主教们。\par

我不知道你们在喧闹混乱的公会议中所作决议为何：但看来真理因某些混乱无序而受到歪曲，以致你们因执意让相互争论取胜，未能察觉何者悦乐天主。然而，将由\OriginalTerm{天主眷顾}{Divine Providence}来驱散这好争精神明显引发的诸般恶事，并向我们清楚显明，你们聚集于彼时，是否顾念过真理，是否在作出决议时未受偏袒或敌意的影响。因此，我愿你们全体尽速聚集到我的敬虔面前，以便你们能够亲自为你们的所作所为作出真实交代。\par

我之所以认为宜以书信告知你们此事，并召唤你们前来见我，你们将从我的下述话语中得知。前不久我正进入我们那以我的名字命名的、无比幸福的都城\Place{君士坦丁堡}（当时我恰巧骑马），突然，\Person{亚大纳削}主教与同他在一起的某些人，在路中央迎面向我走来，如此出人意料，令我大为惊愕。洞悉一切的天主为我作证：若非我从随从那里自然打听，得知他是谁、正遭受何等不公，我初见他时甚至辨认不出他。那时我并未与他交谈，也未准他进见；他请求陈情，我则拒绝，几乎下令将其带走；然而他愈发大胆，唯一请求的恩典便是召你们前来，好让他在你们面前向我申诉他所遭受的虐待。在我看来，这是合情合理的请求，且合乎时宜，因此我欣然下令修此书予你们，以便组成\Place{提洛}会议的你们全体，从速赶赴我的仁慈法庭，用事实证明你们曾作出公正无私的判决。我说，你们务必当着我——你们也不至否认我是\OriginalTerm{天主的忠仆}{true servant}——的面完成此事。\par

诚然，因着我对天主的虔敬，和平得以处处保全，甚至那些向来不识真理的蛮族，也因我而真正敬拜天主之名。显而易见，不识真理者，同样也不认识天主。尽管如此，一如我先前所言，如今连蛮族也经由我——\OriginalTerm{天主的忠仆}{true servant}——得以认识天主，并学会敬畏祂，他们从事实中看出祂处处是我的盾牌与护佑。因此，他们主要是由于对我的畏惧而敬畏天主。然而，我们这些本应阐述（我甚至不愿说守护）祂美善的神圣奥秘的人，我们，我说，所做的无非是趋向纷争与仇恨，简言之，一切导致人类毁灭的事。但你们要赶紧，如我先前所说，你们全体尽快到我这里来，要确信我必竭尽全力纠正错谬，尤其要使那些无咎无辱的事物得以保存，并稳固建立在天主的法律之中；至于法律的敌人，他们假托祂的圣名引入种种亵渎，必将被驱散、彻底粉碎，全然消灭。\par

87. \Person{尤西比乌斯}及其同伙读了这封信后，由于深知自己所作所为，便阻止其他主教前去，只他们自己去了，即：\Person{尤西比乌斯}、\Person{狄奥格尼乌斯}、\Person{帕特罗斐卢斯}、另一个\Person{尤西比乌斯}、\Person{乌尔萨修斯}和\Person{瓦伦斯}。他们绝口不再提及杯子和\Person{阿尔塞尼乌斯}（因为他们已无此胆量），却捏造一项涉及皇帝本人的新指控，在他面前宣称\Person{亚大纳削}曾威胁要扣留从\Place{亚历山卓}运往皇帝故乡的谷物。主教\Person{亚大曼提乌斯}、\Person{阿努比翁}、\Person{亚加塔孟}、\Person{阿尔贝提翁}和\Person{伯多禄}在场，亲耳听到这话。皇帝的愤怒也证实了此事；因为尽管他前此曾修书谴责他们的不义，但一听到这般控诉，便立刻被激怒，不但不准我申诉，反而将我流放至\Place{高卢}。这再次显露他们的邪恶；因为当堪受怀念的\Person{小君士坦丁}忆起其父所写的内容，将我送回故乡时，他也写了如下书信。\par

\Person{君士坦丁}凯撒，致\Place{亚历山卓}城天主教会子民。\par

我想，你们虔诚的心智不会不知，那可敬律法的阐释者\Person{亚大纳削}，曾被遣往\Place{高卢}一段时间，用意在于：因那嗜血成性、积怨已久的仇敌残忍迫害他，危及他的神圣生命，他可借此逃避恶人们乖张伎俩所导致的无法挽救的灾祸。为此，他被从攻击者的口中夺出，奉命在我的治下度过一段时日，并在此城居住期间获得了充足的供应；尽管他那卓著的德行全然倚靠天主的助佑，对逆境的苦难不屑一顾。如今，我主\Person{君士坦丁}奥古斯都——我的父亲——固定意旨，要将该主教送回原处及你们至爱之虔诚那里，但他却被那世人共有的命运夺去，未能实现其愿望便已安息；我遂决意承行这位神圣记忆的皇帝所遗留的旨意。当他亲临你们面前时，你们将得知他受到了何等尊敬。的确，我为他所做的一切不足为奇；因为你们对他思念之情，以及如此伟大人物的现身，都触动我的灵魂，催迫我如此行事。愿\OriginalTerm{天主眷顾}{Divine Providence}常常护佑你们，亲爱的弟兄们。\par

发自\Place{特里尔}，7月朔日前第十五日。\par

88. 这就是我被遣往\Place{高卢}的原因，我再问一次，有谁不能清楚看出皇帝的意图，以及\Person{欧瑟伯}及其同党的杀人意图？皇帝这样做，是为了防止他们策划更疯狂的计划。因为他单纯地听从了他们。\Person{欧瑟伯}及其同党的行径就是如此，他们对我的阴谋也是如此。有谁亲眼见到这些，还会否认我所蒙受的一切并非出于偏私，而是那众多主教，无论是单独还是集体，都为我的案件如此写信，并公正地依照真理谴责我敌人的谎言呢？有谁观察到这些事情，还会否认\Person{瓦伦斯}和\Person{乌尔撒西}有充分理由定自己的罪，并照他们所做的写下那些话，在悔改时控告自己，宁愿短暂受辱，也不愿永远遭受诬告者的惩罚呢？\par

89. 因此，我的有福的同工们，行事公正并按教会法律，当有些人断言我的案情存疑，并试图强迫他们撤销对我有利的判决时，他们如今已忍受了种种苦难，并宁愿被放逐，也不愿看到如此众多主教的裁决被推翻。如果那些真正的主教只用言语抵挡那些谋害我、企图推翻为我所做一切的人；或者他们只是普通人，而不是显赫城市的主教、伟大教会的领袖，那么就有理由怀疑在这种情形下，他们也是出于争竞并为了取悦我而行事。但他们不仅努力以理服人，也忍受了放逐，其中一位是\Place{罗马}主教\Person{利贝里乌斯}（虽然他未忍受放逐之苦至终，但他流亡了两年，知道有反对我们的阴谋正在进行），还有伟大的\Person{何西乌斯}，以及\Place{意大利}、\Place{高卢}的主教们，还有来自\Place{西班牙}、\Place{埃及}、\Place{利比亚}和\Place{五城地区}的所有主教（尽管在一段时间内，出于对\Person{君士坦提乌斯}威胁的恐惧，他看似未抵抗他们，然而\Person{君士坦提乌斯}所施的巨大暴力和专制权力，以及他遭受的许多侮辱和鞭打，证明了他并非放弃我的案子，而是因年迈体弱，无法承受鞭打，才暂时屈服），因此，我说，所有人都应当，既已充分确信，就憎恨并厌恶他们对我的不公和暴行；尤其众所周知，我遭受这些事，别无他因，只是由于\OriginalTerm{亚略的不虔敬}{Arian impiety}。\par

90. 现在若有人愿意了解我的案情，以及\Person{欧瑟伯}及其同党的虚假，就让他阅读为我所写的内容，并聆听见证人，不是一位、两位或三位，而是那众多主教；再让他关注这些诉讼的见证人，\Person{利贝里乌斯}和\Person{何西乌斯}及其同伴，当他们看到敌人对我们的企图时，宁愿忍受种种苦难，也不放弃真理和为我们作出的有利判决。他们这样做是出于可敬和公义的意图，因为他们所遭受的证明了其他主教被逼到了何等困境。他们是反对\OriginalTerm{亚略异端}{Arian heresy}和诬告者邪恶的纪念和记录，并为后来者提供了榜样和模范，让他们为真理奋斗至死\BibleRef{德 4:28}，并憎恶那与基督为敌、是假基督前驱的\OriginalTerm{亚略异端}{Arian heresy}，不去相信那些企图诋毁我的人。因为有如此众多个性正直的主教提出的辩护和作出的判决，是为我们所作的可靠且充分的见证。\par

\SourceRecord{Translated by M. Atkinson and Archibald Robertson.；Nicene and Post-Nicene Fathers, Second Series；Vol. 4.；Edited by Philip Schaff and Henry Wace.；Buffalo, NY: Christian Literature Publishing Co.,；1892.；<http://www.newadvent.org/fathers/28082.htm>.}
